% Leçon 42 — Expression écrite : caractères complexes liés à la gestion des réseaux sociaux — Chinois 1ère A — Cours signé OnBuch+
% Programme officiel MINESEC. Compiler avec : tectonic ZHA42-ecrit-caracteres-complexes-lies-a-la-gestion-des-r.tex
\newcommand{\DOCMATIERE}{Chinois}
\newcommand{\DOCNIVEAU}{1ère A}
\newcommand{\DOCMODULE}{Module 5 : Mass médias et télécommunication}
\newcommand{\DOCLECON}{ZHA42}
\newcommand{\DOCTITRE}{Leçon 42 — Expression écrite : caractères complexes liés à la gestion des réseaux sociaux}
% OnBuch+ — préambule « cours signés » ENRICHI (compilé avec tectonic / XeLaTeX)
% Système visuel « Pop » d'OnBuch+ : fond crème (jamais blanc), bordures encre
% épaisses, ombres « dures » décalées, titres Archivo Black en capitales.
% Version MATHÉMATIQUES : graphes (pgfplots/tikz), boîtes pédagogiques
% (définition, propriété, méthode, attention, le savais-tu, exercice résolu,
% exercice, TP, bilan), page de garde et signature OnBuch+ ancrée en bas.
%
% Macros attendues dans main.tex AVANT \input{preamble} :
%   \DOCMATIERE  \DOCNIVEAU  \DOCMODULE  \DOCLECON (numéro)  \DOCTITRE
\documentclass[11pt]{article}

\usepackage{fontspec}
\usepackage[french]{babel} % français = langue principale ; le chinois via \zh{..} / textebox (xeCJK)
\usepackage{xeCJK}
\usepackage{pifont} % \ding{51} \ding{55} utilisés par les modèles
\usepackage[a4paper,margin=2cm,top=3.4cm,bottom=2.8cm,headheight=1.4cm,footskip=1.3cm]{geometry}
\usepackage{amsmath,amssymb,amsfonts}
\usepackage{mathtools} % \xmapsto, \coloneqq... souvent utilisés par les modèles
\usepackage{cancel} % simplifications barrées (bilans de forces)
% \abs{z} (module, valeur absolue) et \norm{u} (norme), utilisés par les modèles
\providecommand{\abs}{}\let\abs\relax\DeclarePairedDelimiter{\abs}{\lvert}{\rvert}
\providecommand{\norm}{}\let\norm\relax\DeclarePairedDelimiter{\norm}{\lVert}{\rVert}
\usepackage[table]{xcolor} % [table] : fournit aussi \rowcolors (alternance de lignes)
\usepackage{pagecolor}
\usepackage{tikz}
\usetikzlibrary{arrows.meta,positioning,calc,shapes.geometric,shapes.misc,shapes.arrows,decorations.pathmorphing,decorations.pathreplacing,decorations.markings,patterns,shadows,mindmap,trees,fit,backgrounds,matrix,intersections,angles,quotes}
\usepackage{pgfplots}
\usepackage[version=4]{mhchem} % équations nucléaires \ce{^{235}_{92}U}
\usepackage[european,siunitx]{circuitikz} % schémas électriques
\pgfplotsset{compat=1.18}
\usepgfplotslibrary{fillbetween,groupplots} % aires entre courbes (calcul intégral)
\usepackage{enumitem}
\usepackage{fancyhdr}
% [most] charge toutes les bibliothèques tcolorbox (listings, minted, raster,
% documentation...) et gonfle inutilement la mémoire TeX sur un document long
% (~300 boîtes) : on ne charge que ce qui est réellement utilisé.
\usepackage[skins,breakable]{tcolorbox}
\usepackage{graphicx}
\usepackage{colortbl}
\usepackage{booktabs}
\usepackage{tabularx}
\usepackage{diagbox} % case d'en-tête diagonale des tableaux à double entrée
% Colonnes extensibles centrée / gauche / droite (C, L, R), souvent utilisées par les modèles
\newcolumntype{C}{>{\centering\arraybackslash}X}
\newcolumntype{L}{>{\raggedright\arraybackslash}X}
\newcolumntype{R}{>{\raggedleft\arraybackslash}X}
\usepackage{array}
\usepackage{multicol}
\usepackage{multirow}
\usepackage{siunitx}
% parse-numbers=false : accepte aussi \SI{2,5 \times 10^{-3}}{...}
\sisetup{locale=FR,output-decimal-marker={,},group-minimum-digits=4,parse-numbers=false}
\DeclareSIUnit\ohm{\ensuremath{\symup{\Omega}}} % U+2126 absent de la police
\DeclareSIUnit\division{div} % réglages d'oscilloscope (ms/div, V/div)
\DeclareSIUnit\tr{tr} % tours (vitesse de rotation, ex. \SI{1500}{\tr\per\minute})
\DeclareSIUnit\molar{M} % concentration molaire (ex. \SI{3}{\molar}, \SI{1}{\micro\molar})
\DeclareSIUnit\an{an} % année (le modèle confond parfois avec \year, un compteur TeX) — ex. \SI{5}{\kilo\meter\per\an}
\DeclareSIUnit\FCFA{FCFA} % franc CFA (ex. \SI{500}{\FCFA\per\kilo\gram})
\DeclareSIUnit\degreeBrix{\textdegree Brix} % teneur en sucre d'un sirop/jus (réfractométrie)
\DeclareSIUnit\month{mois} % le modèle utilise parfois \month, un compteur TeX, comme unité de durée
\usepackage{qrcode}
% \degree : commande fréquemment utilisée par les modèles pour un angle ou une
% température en dehors de \SI{}{} (non fournie par les paquets chargés).
\providecommand{\degree}{^{\circ}}
% \qed (fin de démonstration), utilisé par les modèles sans amsthm
\providecommand{\qed}{\hfill$\square$}
% \barre : double barre des valeurs interdites dans un tableau de variations (tkz-tab)
\providecommand{\barre}{\ensuremath{\|}}
% environnement proof (amsthm non chargé) utilisé par les modèles
\ifdefined\proof\else\newenvironment{proof}[1][Démonstration]{\par\noindent\textit{#1.}\ }{\qed\par}\fi
\usepackage{lastpage}
\usepackage{hyperref}
\hypersetup{hidelinks}

% ============================================================
% Palette « Pop » OnBuch+ — mêmes hex que la classe P (plus_ui.dart)
% ============================================================
\definecolor{poporange}{HTML}{F59321}
\definecolor{popdark}{HTML}{E07A0C}
\definecolor{popcream}{HTML}{FFF6E8}
\definecolor{popink}{HTML}{1C1714}
\definecolor{poppurple}{HTML}{7A5AE0}
\definecolor{popgreen}{HTML}{2E9E6A}
\definecolor{poppink}{HTML}{E0457B}
\definecolor{popblue}{HTML}{2F7DE1}
\definecolor{popgold}{HTML}{C98A12}
\definecolor{popmuted}{HTML}{8A7F76}
\definecolor{popline}{HTML}{EADFD2}
\definecolor{popgray}{HTML}{8A7F76}
\colorlet{popgrey}{popgray}
% Alias tolérants (le modèle nomme parfois les couleurs différemment)
\colorlet{popwhite}{white}
\colorlet{popblack}{popink}
\colorlet{popred}{poppink}
\colorlet{popyellow}{popgold}
% Teintes claires (fonds de tableaux/encadrés) — le modèle nomme parfois
% les couleurs avec un suffixe L (ex. popblueL) sans les avoir définies.
\colorlet{poporangeL}{poporange!20}
\colorlet{popdarkL}{popdark!20}
\colorlet{poppurpleL}{poppurple!20}
\colorlet{popgreenL}{popgreen!20}
\colorlet{poppinkL}{poppink!20}
\colorlet{popblueL}{popblue!20}
\colorlet{popgoldL}{popgold!20}
\colorlet{popredL}{popred!20}
\colorlet{popgrayL}{popgray!20}
% Teintes claires pour fonds de tableaux / zones
\colorlet{popblueL}{popblue!10!white}
\colorlet{poporangeL}{poporange!14!white}
\colorlet{popgreenL}{popgreen!12!white}
\colorlet{poppurpleL}{poppurple!12!white}
\colorlet{poppinkL}{poppink!12!white}
\colorlet{popgoldL}{popgold!14!white}

\pagecolor{popcream}

% ============================================================
% Typographie
% ============================================================
\defaultfontfeatures{Ligatures=TeX}
\setmainfont{PlusJakartaSans-Medium.ttf}[Path=../../../fonts/,BoldFont=PlusJakartaSans-Bold.ttf]
% Symboles, API et emojis absents de la police principale : repli sur DejaVu Sans (caractères actifs)
\newfontface\symfont{DejaVuSans.ttf}[Path=../../../fonts/]
\makeatletter
\newcommand\symactive[1]{\catcode#1=13 \begingroup\lccode`\~=#1 \lowercase{\endgroup\def~}{{\symfont\char#1}}}
\newcommand\symblank[1]{\catcode#1=13 \begingroup\lccode`\~=#1 \lowercase{\endgroup\def~}{}}
\newcommand\symhyph[1]{\catcode#1=13 \begingroup\lccode`\~=#1 \lowercase{\endgroup\def~}{\mbox{-}}}
\newcount\symc
\newcommand\symrange[2]{\symc=#1 \@whilenum\symc<#2 \do{\expandafter\symactive\expandafter{\the\symc}\advance\symc by 1 }}
\newcommand\symblankrange[2]{\symc=#1 \@whilenum\symc<#2 \do{\expandafter\symblank\expandafter{\the\symc}\advance\symc by 1 }}
\makeatother
\symrange{"0250}{"02FF}\symactive{"2713}\symactive{"2714}\symactive{"2717}\symactive{"2718}\symactive{"2610}\symactive{"2611}\symactive{"2612}
\symrange{"2600}{"27BF}
\catcode"2705=13 \begingroup\lccode`\~="2705 \lowercase{\endgroup\def~}{{\symfont\char"2713}}\symblank{"26BD}\symblank{"26A1}
\symrange{"2460}{"257F}\symactive{"223C}\symactive{"2248}\symactive{"2260}
\symactive{"01F8}\symactive{"01F9}\symactive{"1E3E}\symactive{"1E3F}
\symhyph{"2011}\symhyph{"2E1A}\symblankrange{"1F300}{"1FAFF}\symblank{"FE0F}
% Caractères chinois : WenQuanYi Zen Hei (sous-ensemble GB2312 + pinyin), gras/italique simulés
\setCJKmainfont{WenQuanYiZenHei-subset.ttf}[Path=../../../fonts/,AutoFakeBold=3,AutoFakeSlant=0.2]
\xeCJKsetup{CJKspace=false}
% Pinyin : voyelles à caron (ǎ ǐ ǒ ǔ ǖ ǘ ǚ ǜ) absentes de la police principale -> caractères actifs
\newfontface\pyfont{WenQuanYiZenHei-subset.ttf}[Path=../../../fonts/]
\newcommand\pyactive[1]{\catcode#1=13 \begingroup\lccode`\~=#1 \lowercase{\endgroup\def~}{{\pyfont\char#1}}}
\pyactive{"01CD}\pyactive{"01CE}\pyactive{"01CF}\pyactive{"01D0}\pyactive{"01D1}\pyactive{"01D2}\pyactive{"01D3}\pyactive{"01D4}\pyactive{"01D5}\pyactive{"01D6}\pyactive{"01D7}\pyactive{"01D8}\pyactive{"01D9}\pyactive{"01DA}\pyactive{"01DB}\pyactive{"01DC}
% Maths en sans-serif (Fira Math) avec chiffres et lettres droites de Plus
% Jakarta, pour que les formules se fondent dans le texte.
\usepackage[math-style=ISO,bold-style=ISO]{unicode-math}
\setmathfont{FiraMath-Regular.otf}[Path=../../../fonts/]
\setmathfont{PlusJakartaSans-Medium.ttf}[Path=../../../fonts/,range={up/num,up/{Latin,latin},\mathup}]
\usepackage{icomma} % virgule décimale sans espace en mode math
% Caractères mathématiques Unicode absents des polices de texte (Plus Jakarta,
% Archivo Black) : rendus par leur équivalent mathématique, en texte comme en
% formule — sinon ils s'affichent en carré vide (ex. « Arithmétique dans ℤ »).
\usepackage{newunicodechar}
\newunicodechar{ℝ}{\ensuremath{\mathbb{R}}}
\newunicodechar{ℤ}{\ensuremath{\mathbb{Z}}}
\newunicodechar{ℂ}{\ensuremath{\mathbb{C}}}
\newunicodechar{ℕ}{\ensuremath{\mathbb{N}}}
\newunicodechar{ℚ}{\ensuremath{\mathbb{Q}}}
\newunicodechar{𝔻}{\ensuremath{\mathbb{D}}}
\newunicodechar{α}{\ensuremath{\alpha}}
\newunicodechar{β}{\ensuremath{\beta}}
\newunicodechar{γ}{\ensuremath{\gamma}}
\newunicodechar{δ}{\ensuremath{\delta}}
\newunicodechar{ε}{\ensuremath{\varepsilon}}
\newunicodechar{θ}{\ensuremath{\theta}}
\newunicodechar{λ}{\ensuremath{\lambda}}
\newunicodechar{μ}{\ensuremath{\mu}}
\newunicodechar{σ}{\ensuremath{\sigma}}
\newunicodechar{τ}{\ensuremath{\tau}}
\newunicodechar{φ}{\ensuremath{\varphi}}
\newunicodechar{ψ}{\ensuremath{\psi}}
\newunicodechar{ω}{\ensuremath{\omega}}
\newunicodechar{Σ}{\ensuremath{\Sigma}}
% majuscules produites par \MakeUppercase dans les titres (α-aminés -> Α)
\newunicodechar{Α}{\ensuremath{\alpha}}
\newunicodechar{Β}{\ensuremath{\beta}}
\newunicodechar{Γ}{\ensuremath{\Gamma}}
\newunicodechar{Δ}{\ensuremath{\Delta}}
\newunicodechar{Ε}{\ensuremath{\varepsilon}}
\newunicodechar{Θ}{\ensuremath{\Theta}}
\newunicodechar{Λ}{\ensuremath{\Lambda}}
\newunicodechar{Μ}{\ensuremath{\mu}}
\newunicodechar{Τ}{\ensuremath{\tau}}
\newunicodechar{Φ}{\ensuremath{\Phi}}
\newunicodechar{Ψ}{\ensuremath{\Psi}}
\newunicodechar{Ω}{\ensuremath{\Omega}}
\newunicodechar{↦}{\ensuremath{\mapsto}}
\newunicodechar{∈}{\ensuremath{\in}}
\newunicodechar{∉}{\ensuremath{\notin}}
\newunicodechar{∧}{\ensuremath{\wedge}}
\newunicodechar{⇔}{\ensuremath{\Leftrightarrow}}
\newunicodechar{⟺}{\ensuremath{\Longleftrightarrow}}
\newunicodechar{⇒}{\ensuremath{\Rightarrow}}
\newunicodechar{⟹}{\ensuremath{\Longrightarrow}}
\newunicodechar{∘}{\ensuremath{\circ}}
\newunicodechar{∪}{\ensuremath{\cup}}
\newunicodechar{∩}{\ensuremath{\cap}}
\newunicodechar{≡}{\ensuremath{\equiv}}
\newunicodechar{⊂}{\ensuremath{\subset}}
\newunicodechar{⊥}{\ensuremath{\perp}}
\newunicodechar{∃}{\ensuremath{\exists}}
\newunicodechar{∀}{\ensuremath{\forall}}
\newunicodechar{∛}{\ensuremath{\sqrt[3]{\,}}}
\newunicodechar{∜}{\ensuremath{\sqrt[4]{\,}}}
\newunicodechar{⃗}{\ensuremath{{}^{\to}}}
\newunicodechar{ⁿ}{\ifmmode^{n}\else\textsuperscript{\ensuremath{n}}\fi}
\newunicodechar{ˣ}{\ifmmode^{x}\else\textsuperscript{\ensuremath{x}}\fi}
\newunicodechar{ᵇ}{\ifmmode^{b}\else\textsuperscript{\ensuremath{b}}\fi}
\newunicodechar{ʳ}{\ifmmode^{r}\else\textsuperscript{\ensuremath{r}}\fi}
\newunicodechar{ᵘ}{\ifmmode^{u}\else\textsuperscript{\ensuremath{u}}\fi}
\newunicodechar{ʸ}{\ifmmode^{y}\else\textsuperscript{\ensuremath{y}}\fi}
\newunicodechar{ᵃ}{\ifmmode^{a}\else\textsuperscript{\ensuremath{a}}\fi}
\newunicodechar{ᵉ}{\ifmmode^{e}\else\textsuperscript{\ensuremath{e}}\fi}
\newunicodechar{ᵏ}{\ifmmode^{k}\else\textsuperscript{\ensuremath{k}}\fi}
\newunicodechar{ᵖ}{\ifmmode^{p}\else\textsuperscript{\ensuremath{p}}\fi}
\newunicodechar{ᴺ}{\ifmmode^{N}\else\textsuperscript{\ensuremath{N}}\fi}
\newunicodechar{ᶜ}{\ifmmode^{c}\else\textsuperscript{\ensuremath{c}}\fi}
\newunicodechar{ʰ}{\ifmmode^{h}\else\textsuperscript{\ensuremath{h}}\fi}
\newunicodechar{ᵠ}{\ifmmode^{q}\else\textsuperscript{\ensuremath{q}}\fi}
\newunicodechar{ₙ}{\ifmmode_{n}\else\textsubscript{\ensuremath{n}}\fi}
\newunicodechar{ₐ}{\ifmmode_{a}\else\textsubscript{\ensuremath{a}}\fi}
\newunicodechar{ᵢ}{\ifmmode_{i}\else\textsubscript{\ensuremath{i}}\fi}
\newunicodechar{ᵣ}{\ifmmode_{r}\else\textsubscript{\ensuremath{r}}\fi}
\newunicodechar{ₖ}{\ifmmode_{k}\else\textsubscript{\ensuremath{k}}\fi}
\newunicodechar{ᵦ}{\ifmmode_{\beta}\else\textsubscript{\ensuremath{\beta}}\fi}
\newunicodechar{ₓ}{\ifmmode_{x}\else\textsubscript{\ensuremath{x}}\fi}
\newfontfamily\popdisplay{ArchivoBlack-Regular.ttf}[Path=../../../fonts/]
\newfontfamily\poplogo{SpaceGrotesk-Bold.ttf}[Path=../../../fonts/]
\newfontfamily\popsemi{PlusJakartaSans-SemiBold.ttf}[Path=../../../fonts/]
\newfontfamily\popxbold{PlusJakartaSans-ExtraBold.ttf}[Path=../../../fonts/]
\newfontfamily\popxboldls{PlusJakartaSans-ExtraBold.ttf}[Path=../../../fonts/,LetterSpace=3.0]

\setlist[enumerate]{leftmargin=1.7em,itemsep=5pt,topsep=4pt}
\setlist[itemize]{leftmargin=1.7em,itemsep=4pt,topsep=4pt,label={\color{poporange}\textbullet}}
\setlength{\parindent}{0pt}
\setlength{\parskip}{5pt}
\renewcommand{\arraystretch}{1.25}

% pgfplots : style Pop par défaut
\pgfplotsset{
  every axis/.append style={axis line style={popink,line width=1pt},tick style={popink},
    grid style={popline,line width=0.5pt},label style={font=\small},tick label style={font=\footnotesize}},
  popcurve/.style={poporange,line width=1.8pt,no markers,smooth},
}
% Même style disponible dans un \draw TikZ simple (hors pgfplots)
\tikzset{popcurve/.style={poporange,line width=1.8pt}}
\tikzset{popgrid/.style={popline,very thin}}
\colorlet{popcurve}{poporange} % le nom du style est parfois utilisé comme couleur

% ============================================================
% Pill / squiggle
% ============================================================
\newcommand{\poppill}[3]{%
  \tikz[baseline=-0.55ex]\node[draw=popink,line width=1.2pt,rounded corners=2.6pt,%
    fill=#2,inner xsep=6.5pt,inner ysep=3pt]{%
    \popxboldls\fontsize{7}{7}\selectfont\color{#3}\MakeUppercase{#1}};%
}
\newcommand{\popsquiggle}[1]{%
  \tikz[baseline=-0.4ex]{
    \draw[#1,line width=2.6pt,line cap=round]
      plot [smooth] coordinates {(0,0.09) (0.34,0.01) (0.68,0.21) (1.02,0.01) (1.36,0.10)};
  }%
}
% Mot-clé surligné (marqueur orange sous le texte)
\newcommand{\cle}[1]{{\popxbold\color{popdark}#1}}

% ============================================================
% En-tête / pied
% ============================================================
\pagestyle{fancy}
\fancyhf{}
\renewcommand{\headrulewidth}{0pt}
\renewcommand{\footrulewidth}{0pt}
\lhead{\raisebox{-0.15ex}{\poplogo\fontsize{13}{13}\selectfont\color{popink}OnBuch\color{poporange}+}}
\rhead{\poppill{\DOCMATIERE\ \textbullet\ \DOCNIVEAU\ \textbullet\ Leçon \DOCLECON}{white}{popink}}
\lfoot{\popxbold\fontsize{7.6}{7.6}\selectfont\color{popmuted}\MakeUppercase{Cours signé OnBuch+}\ \textbullet\ {\popsemi\DOCTITRE}}
\rfoot{\tikz[baseline=-0.6ex]\node[rounded corners=8.5pt,fill=popink,minimum height=17pt,minimum width=17pt,inner xsep=5pt,inner sep=0]{\popxbold\fontsize{8}{8}\selectfont\color{popcream}\thepage\,/\,\pageref*{LastPage}};}

% ============================================================
% Titres
% ============================================================
\newcommand{\chaptitre}[2]{%
  \begin{tcolorbox}[enhanced,colback=poporange,colframe=popink,boxrule=2.6pt,arc=11pt,
      drop shadow=popink,left=15pt,right=15pt,top=12pt,bottom=11pt,before skip=3pt,after skip=18pt]
    \poppill{#2}{popink}{popcream}\par\vspace{9pt}
    {\raggedright\hyphenpenalty=10000\exhyphenpenalty=10000\popdisplay\fontsize{22}{24}\selectfont\color{popcream}\MakeUppercase{#1}\par}
  \end{tcolorbox}
}
% Section numérotée : pastille orange avec le numéro + titre Archivo
\newcounter{popsec}
\newcommand{\coursec}[1]{%
  \stepcounter{popsec}\setcounter{popsub}{0}%
  \par\vspace{16pt}\noindent
  \begin{minipage}{\linewidth}
  \tikz[baseline=-0.7ex]\node[draw=popink,line width=1.6pt,fill=poporange,rounded corners=4pt,minimum size=22pt,inner sep=0]{\popdisplay\fontsize{12}{12}\selectfont\color{popcream}\thepopsec};%
  \hspace{8pt}{\popdisplay\fontsize{15}{17}\selectfont\color{popink}\MakeUppercase{#1}}\par
  \vspace{4pt}\noindent{\color{poporange}\rule{36pt}{3.6pt}}
  \end{minipage}\par\nopagebreak\vspace{8pt}
}
\newcounter{popsub}
\newcommand{\courssub}[1]{%
  \stepcounter{popsub}%
  \par\vspace{9pt}\noindent{\popxbold\fontsize{11.5}{13}\selectfont\color{popink}{\color{poporange}\thepopsec.\thepopsub}\hspace{6pt}#1}\par\nopagebreak\vspace{4pt}
}

% ============================================================
% Boîtes pédagogiques — même recette : fond blanc, bordure épaisse colorée,
% ombre dure encre, badge plein. Titre optionnel [#1].
% ============================================================
\tcbset{popbase/.style={breakable,enhanced,colback=white,boxrule=2.3pt,arc=9pt,
  shadow={3pt}{-3pt}{0pt}{popink},left=14pt,right=14pt,top=11pt,bottom=11pt,before skip=11pt,after skip=15pt,
  fontupper=\popsemi\fontsize{10.6}{14.5}\selectfont\color{popink},
  fontlower=\popsemi\fontsize{10.6}{14.5}\selectfont\color{popink}}}
% Coche dessinée (le glyphe ✓ n'existe pas dans les polices du document)
\newcommand{\popcheck}[1]{\tikz[baseline=-0.5ex]\draw[#1,line width=1.6pt,line cap=round,line join=round](-0.13,0.0)--(-0.04,-0.09)--(0.13,0.1);}
\providecommand{\coche}{\popcheck{popgreen}}
\providecommand{\resultat}[1]{\par\smallskip\noindent\fcolorbox{popgreen}{popgreenL}{\parbox{\dimexpr\linewidth-2\fboxsep-2\fboxrule\relax}{\textbf{Résultat.} #1}}\par\smallskip}
\newcommand{\popboxhead}[3]{\poppill{#1}{#2}{white}\ifx\relax#3\relax\else\hspace{6pt}{\popxbold\fontsize{10.6}{12}\selectfont\color{popink}#3}\fi\par\vspace{7pt}}

\newtcolorbox{aretenir}[1][]{popbase,colframe=popgreen,before upper={\popboxhead{À retenir}{popgreen}{#1}}}
% Texte en chinois (document de lecture, dialogue, extrait) : cadre
% d'étude. Un titre optionnel (source, auteur) s'affiche dans l'en-tête.
\newtcolorbox{textebox}[1][]{popbase,colframe=popblue,before upper={\popboxhead{文本 · Texte}{popblue}{#1}},after upper={}}
% Marques ✓ / ✗ (forme correcte / fautive) souvent utilisées par les modèles
\providecommand{\cross}{\textcolor{poppink}{\ensuremath{\times}}}
\providecommand{\xmark}{\cross}\providecommand{\crossmark}{\cross}\providecommand{\cmark}{\ensuremath{\checkmark}}
% Ligne de réponse / justification à compléter (inventée par les modèles)
\providecommand{\justification}[1]{\par\noindent\textit{Justification~:} \dotfill\par}
% \textipa (package tipa non chargé) : la police ne couvre pas l'API -> simple passage
\providecommand{\textipa}[1]{#1}
% Soulignement (ulem non chargé) : \uline, \ul
\providecommand{\uline}[1]{\underline{#1}}\providecommand{\ul}[1]{\underline{#1}}
% \glossaire{\item ...} inventée par les modèles : liste de glossaire
\providecommand{\glossaire}[1]{\begin{itemize}#1\end{itemize}}
% \alert (beamer) inventée par les modèles : texte mis en évidence
\providecommand{\alert}[1]{\textbf{\textcolor{poppink}{#1}}}
% variantes ASCII d'ordinaux écrites en macro
\providecommand{\iere}{\textsuperscript{re}}\providecommand{\ieme}{\textsuperscript{e}}\providecommand{\ere}{\textsuperscript{re}}\providecommand{\eme}{\textsuperscript{e}}\providecommand{\er}{\textsuperscript{er}}\providecommand{\ie}{\textsuperscript{e}}
% Ordinaux français écrits en macro par les modèles : 1\ière, 3\ième
\newcommand{\ière}{\textsuperscript{re}}\newcommand{\ième}{\textsuperscript{e}}
% \exemple (inventée par les modèles) : étiquette « Exemple : »
\providecommand{\exemple}{\textbf{Exemple~:}\ }
% \square utilisable aussi hors mode math (cases à cocher)
\AtBeginDocument{\let\squareMath\square\renewcommand{\square}{\ensuremath{\squareMath}}}
% Mot ou phrase en chinois à l'intérieur d'un paragraphe français
\newcommand{\zh}[1]{\ifmmode\text{#1}\else{#1}\fi}
\let\eng\zh \let\esp\zh \let\all\zh % alias tolérés
% flèches d'intonation inventées par les modèles
\providecommand{\textsearrow}{}\renewcommand{\textsearrow}{\ensuremath{\searrow}}\providecommand{\textnearrow}{}\renewcommand{\textnearrow}{\ensuremath{\nearrow}}
% \fr{..} : mot français (inventé par les modèles)
\providecommand{\fr}[1]{\textit{#1}}
% zhbox : encadré de texte chinois inventé par les modèles
\newenvironment{zhbox}{\begin{quote}}{\end{quote}}
% \courssubsub : titre de niveau 3 inventé par les modèles
\providecommand{\courssubsub}[1]{\par\medskip\noindent\textbf{#1}\par\smallskip}
% \zhbf : texte chinois en gras inventé par les modèles
\providecommand{\zhbf}[1]{\textbf{#1}}
% \vs : « versus » inventé par les modèles
\providecommand{\vs}{\textit{vs}\ }
% \blank : case à compléter inventée par les modèles
\providecommand{\blank}[1][]{\underline{\hspace{1.6em}}}
\newtcolorbox{exemplebox}[1][]{popbase,colframe=poporange,before upper={\popboxhead{Exemple}{poporange}{#1}}}
\newtcolorbox{definition}[1][]{popbase,colframe=popblue,before upper={\popboxhead{Définition}{popblue}{#1}}}
\newtcolorbox{propriete}[1][]{popbase,colframe=poppurple,before upper={\popboxhead{Propriété}{poppurple}{#1}}}
\newtcolorbox{methode}[1][]{popbase,colframe=popgold,before upper={\popboxhead{Méthode}{popgold}{#1}}}
\newtcolorbox{attention}[1][]{popbase,colframe=poppink,before upper={\popboxhead{Attention}{poppink}{#1}}}
\newtcolorbox{experience}[1][]{popbase,colframe=popblue,colback=popblueL,before upper={\popboxhead{Expérience}{popblue}{#1}}}
% « Le savais-tu ? » : carte violette pleine (contexte, culture, Cameroun)
\newtcolorbox{savaistu}[1][]{popbase,colback=poppurple,colframe=popink,
  fontupper=\popsemi\fontsize{10.4}{14.2}\selectfont\color{white},
  before upper={\poppill{Le savais-tu\,?}{white}{poppurple}\ifx\relax#1\relax\else\hspace{6pt}{\popxbold\color{white}#1}\fi\par\vspace{7pt}}}
% Exercice résolu : énoncé en haut, \tcblower puis la solution sur fond crème
\newcounter{popexo}\newcounter{popexr}
\newtcolorbox{exoresolu}[1][]{popbase,colframe=poporange,colbacklower=popcream,
  segmentation style={popink,line width=1.2pt,dashed},
  before upper={\stepcounter{popexr}\popboxhead{Exercice résolu \thepopexr}{poporange}{#1}},
  before lower={\poppill{Solution}{popink}{popcream}\par\vspace{6pt}}}
% Exercice (non résolu en place) — la correction est dans la partie Corrigés
\newtcolorbox{exercice}[1][]{popbase,colframe=popink,
  before upper={\stepcounter{popexo}\popboxhead{Exercice \thepopexo}{popink}{#1}}}
% Corrigé
\newtcolorbox{corrige}[1][]{popbase,colframe=popgreen,colback=popgreenL,
  before upper={\popboxhead{Corrigé}{popgreen}{#1}}}
% Objectifs / prérequis / bilan
\newtcolorbox{objectifs}{popbase,colframe=popink,colback=poporangeL,
  before upper={\popboxhead{Objectifs de la leçon}{popink}{}}}
\newtcolorbox{prerequis}{popbase,colframe=popink,colback=popblueL,
  before upper={\popboxhead{Prérequis}{popblue}{}}}
\newtcolorbox{bilan}{popbase,colframe=popink,colback=white,boxrule=2.8pt,
  before upper={\popboxhead{Fiche bilan}{poporange}{L'essentiel à connaître pour l'examen}}}
% Figure encadrée avec légende
\newtcolorbox{popfigure}[1][]{enhanced,breakable=false,colback=white,colframe=popline,boxrule=1.4pt,arc=8pt,
  left=10pt,right=10pt,top=10pt,bottom=8pt,before skip=10pt,after skip=12pt,halign=center,
  lower separated=false,halign lower=center,
  fontlower=\popsemi\fontsize{9}{11.5}\selectfont\color{popmuted},
  #1}
\newcommand{\legende}[1]{\par\vspace{6pt}{\popsemi\fontsize{9}{11.5}\selectfont\color{popmuted}#1}}

% ============================================================
% Signature OnBuch+
% ============================================================
\newcommand{\poplogoinline}[1]{{\poplogo\fontsize{#1}{#1}\selectfont\color{popink}OnBuch\color{poporange}+}}
\newcommand{\signatureonbuch}{%
  \begin{tcolorbox}[enhanced,colback=popink,colframe=popink,boxrule=2.2pt,arc=10pt,
      drop shadow=poporange,left=14pt,right=14pt,top=10pt,bottom=10pt,before skip=0pt,after skip=0pt]
    \tikz[baseline=-0.7ex]\node[circle,fill=popgreen,draw=popcream,line width=1.3pt,minimum size=22pt,inner sep=0]{\popcheck{white}};%
    \hspace{9pt}\begin{minipage}[c]{0.62\linewidth}
      {\popxbold\fontsize{12}{13}\selectfont\color{popcream}Signé\ \textbullet\ {\poplogo OnBuch\color{poporange}+}}\par\vspace{2pt}
      {\raggedright\popsemi\fontsize{8}{10.5}\selectfont\color{popline}\MakeUppercase{Équipe pédagogique OnBuch+}\\\MakeUppercase{Conforme au programme officiel MINESEC}\par}
    \end{minipage}\hfill
    {\poplogo\fontsize{22}{22}\selectfont\color{popcream}OnBuch\color{poporange}+}
  \end{tcolorbox}%
}

% ============================================================
% Page de garde
% #1 titre · #2 matière · #3 niveau · #4 module · #5 n° leçon · #6 durée ·
% #7 description / objectifs (texte ou itemize)
% ============================================================
\newcommand{\pagegarde}[7]{%
  \thispagestyle{empty}
  \newgeometry{a4paper,margin=1.7cm,top=1.6cm,bottom=1.4cm}%
  \begin{tikzpicture}[remember picture,overlay]
    \fill[poporange!18] ([xshift=-3.2cm,yshift=-3.6cm]current page.north east) circle (4.6cm);
    \fill[poppurple!14] ([xshift=2.4cm,yshift=7.6cm]current page.south west) circle (3.8cm);
  \end{tikzpicture}%
  {\poplogo\fontsize{30}{30}\selectfont\color{popink}OnBuch\color{poporange}+}\hfill
  \raisebox{0.6ex}{\poppill{Cours signé \textbullet\ Premium}{popink}{popcream}}\par
  \vspace{1pt}\popsquiggle{poppurple}\par
  \vspace{8pt}
  \begin{tcolorbox}[enhanced,colback=poporange,colframe=popink,boxrule=2.8pt,arc=13pt,
      drop shadow=popink,left=18pt,right=18pt,top=12pt,bottom=13pt,after skip=10pt]
    \poppill{#2 \textbullet\ #3}{popink}{popcream}\hspace{5pt}\poppill{#4}{white}{popink}\par\vspace{8pt}
    {\popdisplay\fontsize{14}{14}\selectfont\color{popink}LEÇON #5}\par\vspace{4pt}
    {\raggedright\hyphenpenalty=10000\exhyphenpenalty=10000\popdisplay\fontsize{25}{27}\selectfont\color{popcream}\MakeUppercase{#1}\par}
    \vspace{8pt}
    {\popxbold\fontsize{9.5}{9.5}\selectfont\color{popink}\MakeUppercase{Durée conseillée : #6}}
  \end{tcolorbox}
  \begin{tcolorbox}[enhanced,colback=white,colframe=popink,boxrule=2.2pt,arc=10pt,
      drop shadow=popink,left=16pt,right=16pt,top=10pt,bottom=10pt,after skip=8pt]
    \poppill{Dans cette leçon}{poppurple}{white}\par\vspace{5pt}
    \popsemi\fontsize{9.3}{12.6}\selectfont\color{popink}#7
  \end{tcolorbox}
  \noindent\begin{minipage}[t]{0.487\linewidth}
    \begin{tcolorbox}[enhanced,colback=popgreen,colframe=popink,boxrule=2.2pt,arc=10pt,
        drop shadow=popink,left=11pt,right=11pt,top=10pt,bottom=10pt,height=3.1cm,valign=center]
      \tcbox[colback=white,colframe=popink,boxrule=1.3pt,arc=5pt,boxsep=0pt,nobeforeafter,
        left=3pt,right=3pt,top=3pt,bottom=3pt,valign=center]{\qrcode[height=1.9cm]{https://play.google.com/store/apps/details?id=com.luvvix.onbuch&hl=fr}}%
      \hspace{8pt}\begin{minipage}[c]{0.5\linewidth}
        {\popdisplay\fontsize{11}{12.5}\selectfont\color{white}\MakeUppercase{Télécharge OnBuch}}\par\vspace{4pt}
        {\raggedright\popsemi\fontsize{8.4}{11}\selectfont\color{white}Gratuit \textbullet\ Google Play\\Tuteur IA, annales, concours, résultats\par}
      \end{minipage}
    \end{tcolorbox}
  \end{minipage}\hfill
  \begin{minipage}[t]{0.487\linewidth}
    \begin{tcolorbox}[enhanced,colback=poppurple,colframe=popink,boxrule=2.2pt,arc=10pt,
        drop shadow=popink,left=11pt,right=11pt,top=10pt,bottom=10pt,height=3.1cm,valign=center]
      \tikz[baseline=-0.7ex]\node[circle,fill=white,draw=popink,line width=1.3pt,minimum size=1.6cm,inner sep=0]{\popdisplay\fontsize{24}{24}\selectfont\color{poppurple}+};%
      \hspace{8pt}\begin{minipage}[c]{0.6\linewidth}
        {\popdisplay\fontsize{11}{12.5}\selectfont\color{white}\MakeUppercase{Passe à OnBuch+}}\par\vspace{4pt}
        {\raggedright\popsemi\fontsize{8.4}{11}\selectfont\color{white}Tous les cours signés, Tuteur IA illimité, fiches PDF — onglet OnBuch+ de l'app\par}
      \end{minipage}
    \end{tcolorbox}
  \end{minipage}
  \vspace{8pt}
  \popcontenu
  \vfill
  \signatureonbuch
  \restoregeometry
  \newpage
}

% Rangée « Ce cours contient » (page de garde)
\newcommand{\poptile}[3]{% #1 couleur #2 gros texte #3 légende
  \begin{tcolorbox}[enhanced,colback=#1,colframe=popink,boxrule=2pt,arc=9pt,shadow={2.5pt}{-2.5pt}{0pt}{popink},
      left=6pt,right=6pt,top=9pt,bottom=9pt,halign=center,valign=center,height=1.75cm,width=\linewidth,nobeforeafter]
    {\popdisplay\fontsize{12.5}{13}\selectfont\color{white}\MakeUppercase{#2}}\par\vspace{3pt}
    {\centering\popsemi\fontsize{7.8}{9.5}\selectfont\color{white}#3\par}
  \end{tcolorbox}}
\newcommand{\popcontenu}{%
  \par\noindent{\popxboldls\fontsize{7.5}{7.5}\selectfont\color{popmuted}CE COURS SIGNÉ CONTIENT}\par\vspace{7pt}
  \noindent\begin{minipage}[t]{0.235\linewidth}\poptile{popblue}{Cours}{complet, illustré, pas à pas}\end{minipage}\hfill
  \begin{minipage}[t]{0.235\linewidth}\poptile{poporange}{Méthodes}{et exercices résolus}\end{minipage}\hfill
  \begin{minipage}[t]{0.235\linewidth}\poptile{poppink}{Exercices}{type examen + corrigés}\end{minipage}\hfill
  \begin{minipage}[t]{0.235\linewidth}\poptile{popgreen}{Fiche bilan}{l'essentiel à retenir}\end{minipage}\par}

% Sommaire
\newtcolorbox{sommaire}{enhanced,colback=white,colframe=popink,boxrule=2.2pt,arc=10pt,
  drop shadow=popink,left=18pt,right=18pt,top=13pt,bottom=13pt,before skip=6pt,after skip=20pt,
  before upper={\poppill{Sommaire}{poppurple}{white}\par\vspace{9pt}\popsemi\fontsize{10.6}{14.5}\selectfont\color{popink}}}

% Signature de fin de cours — poussée tout en bas de la dernière page
\newcommand{\findecours}{%
  \par\vfill\nopagebreak
  \begin{center}\popsquiggle{poporange}\end{center}\vspace{4pt}
  \signatureonbuch
}
\AtBeginDocument{\def\llbracket{\mathopen{[\mkern-3.5mu[}}\def\rrbracket{\mathclose{]\mkern-3.5mu]}}} % intervalles d'entiers (glyphe ⟦ absent de la police)
% Glyphes absents de Fira Math : redéfinis à partir de glyphes disponibles
\AtBeginDocument{%
  \def\setminus{\mathbin{\backslash}}\let\smallsetminus\setminus
  \def\vdots{\mathord{\vbox{\baselineskip3.2pt\lineskiplimit0pt\kern4pt\hbox{.}\hbox{.}\hbox{.}}}}%
  \def\ddots{\mathinner{\mkern1mu\raise6.5pt\hbox{.}\mkern2mu\raise3.6pt\hbox{.}\mkern2mu\raise0.7pt\hbox{.}\mkern1mu}}%
  \def\triangle{\mathord{\scalebox{0.9}{$\symup{\Delta}$}}}%
  \def\nabla{\mathord{\raisebox{\depth}{\scalebox{1}[-1]{$\symup{\Delta}$}}}}%
  \def\checkmark{\popcheck{popdark}}%
  \def\star{\mathbin{\ast}}\def\bigstar{\mathord{\ast}}%
  \def\diamond{\mathbin{\scalebox{0.6}{$\Diamond$}}}%
  \def\bigcup{\mathop{\mathchoice{\raisebox{-0.3em}{\scalebox{1.7}{$\cup$}}}{\scalebox{1.25}{$\cup$}}{\cup}{\cup}}\displaylimits}%
  \def\bigcap{\mathop{\mathchoice{\raisebox{-0.3em}{\scalebox{1.7}{$\cap$}}}{\scalebox{1.25}{$\cap$}}{\cap}{\cap}}\displaylimits}%
  \def\lesseqqgtr{\mathrel{\lessgtr}}\def\bigoplus{\mathop{\oplus}\displaylimits}%
}
\providecommand{\ssi}{si et seulement si} % abréviation fréquente
\AtBeginDocument{\providecommand{\wideparen}{\overparen}} % arc de cercle

%% FIGFIX-BEGIN v5 — correctifs globaux des figures (docs/onbuch-generation/tools/figfix)
\usepackage{adjustbox}\usepackage{etoolbox}\tcbuselibrary{raster}
% toute figure TikZ/pgfplots plus large que la ligne est réduite à la largeur disponible
\BeforeBeginEnvironment{tikzpicture}{\begin{adjustbox}{max width=\linewidth}}
\AfterEndEnvironment{tikzpicture}{\end{adjustbox}}
% légendes pgfplots lisibles par-dessus les courbes
\pgfplotsset{every axis/.append style={legend style={fill=white,fill opacity=0.92,text opacity=1,draw=black!15,inner sep=3pt}}}
% étiquettes d'arêtes (syntaxe edge["..."]) avec fond blanc
\tikzset{every edge quotes/.append style={fill=white,inner sep=1.2pt}}
%% FIGFIX-END
\begin{document}

\pagegarde{Leçon 42 — Expression écrite : caractères complexes liés à la gestion des réseaux sociaux}{Chinois}{1ère A}{Module 5 : Mass médias et télécommunication}{ZHA42}{2 h}{Cette leçon d'expression écrite entraîne les élèves de 1ère A à écrire correctement les caractères complexes liés à la gestion des réseaux sociaux (publier, partager, commenter, bloquer, gérer un compte) et à produire ou traduire des textes simples sur ce thème, en français comme en chinois.
\vspace{4pt}
\begin{multicols}{2}\small
\begin{itemize}
\item À la fin de cette leçon, je sais écrire correctement les caractères clés du thème (发, 享, 评, 论, 关, 注, 管, 理, 账, 号, 屏, 蔽, 删, 除) en respectant radicaux et ordre des traits.
\item Je sais distinguer les caractères proches (己/已, 未/末, 辨/辩, 账/帐) et éviter les confusions.
\item Je sais rédiger un texte simple de 5 à 8 phrases sur la gestion d'un compte de réseau social.
\item Je sais utiliser les connecteurs 首先, 然后, 但是, 所以, 因为 pour structurer mon texte.
\item Je sais traduire des phrases simples du français vers le chinois (thème) sur ce sujet.
\item Je sais traduire un court texte du chinois vers le français (version).
\end{itemize}\end{multicols}}

\begin{sommaire}
\begin{multicols}{2}
\begin{enumerate}
\item Modèle de texte commenté
\item Lexique de la gestion des réseaux sociaux
\item Écriture correcte des caractères
\item Structure et connecteurs du texte
\item Thème : du français vers le chinois
\item Version : du chinois vers le français
\item Grille d'évaluation et production finale
\item Activité d'intégration
\item Méthodes et exercices résolus
\item Exercices
\item Corrigés des exercices
\item Fiche bilan
\end{enumerate}
\end{multicols}
\end{sommaire}

\begin{objectifs}
\begin{itemize}
\item À la fin de cette leçon, je sais écrire correctement les caractères clés du thème (发, 享, 评, 论, 关, 注, 管, 理, 账, 号, 屏, 蔽, 删, 除) en respectant radicaux et ordre des traits.
\item Je sais distinguer les caractères proches (己/已, 未/末, 辨/辩, 账/帐) et éviter les confusions.
\item Je sais rédiger un texte simple de 5 à 8 phrases sur la gestion d'un compte de réseau social.
\item Je sais utiliser les connecteurs 首先, 然后, 但是, 所以, 因为 pour structurer mon texte.
\item Je sais traduire des phrases simples du français vers le chinois (thème) sur ce sujet.
\item Je sais traduire un court texte du chinois vers le français (version).
\item Je sais repérer et corriger mes erreurs d'écriture grâce à une grille d'évaluation.
\item Je sais respecter la structure sujet-verbe-complément et la place des compléments de temps en chinois.
\end{itemize}
\end{objectifs}

\begin{prerequis}
\begin{itemize}
\item Vocabulaire de base des mass médias et d'Internet (手机, 网络, 微信, 信息) vu dans le Module 5
\item Radicaux fréquents appris en Seconde (讠parole, 扌main, 心cœur, 宀toit)
\item Ordre général des traits : de haut en bas, de gauche à droite, horizontal avant vertical
\item Structure de base de la phrase chinoise S-V-O et particule 了
\item Connecteurs simples déjà rencontrés : 和, 也, 很
\end{itemize}
\end{prerequis}

\newpage

\coursec{Modèle de texte commenté}

À Douala comme à Yaoundé, Amina gère la page Facebook de la boutique de sa mère : elle publie des photos des nouvelles robes, répond aux commentaires des clientes, supprime les spams et suit le nombre d'abonnés qui grandit chaque semaine. Un jour, une partenaire chinoise lui écrit : \zh{请管理一下我们的账号。} (\emph{Qǐng guǎnlǐ yīxià wǒmen de zhànghào.} — « Veuillez gérer notre compte. ») Amina réalise alors que derrière chaque bouton de son téléphone — publier, partager, bloquer, suivre — se cache un caractère chinois précis qu'elle doit savoir écrire sans faute. Comment écrire correctement ces caractères complexes et rédiger un petit texte sur la gestion des réseaux sociaux ? C'est tout l'enjeu de cette leçon d'expression écrite. Pour y répondre, nous commençons par lire, observer et décortiquer un \cle{texte modèle}.

\courssub{Lecture du texte modèle}

\begin{textebox}[我管理妈妈的网店账号 — Je gère le compte de la boutique en ligne de ma mère]
我帮妈妈管理网店的微信账号。\\
\emph{Wǒ bāng māma guǎnlǐ wǎngdiàn de Wēixìn zhànghào.}\\
« J'aide ma mère à gérer le compte WeChat de la boutique en ligne. »\\[0.4em]
首先，我每天发布新的照片。\\
\emph{Shǒuxiān, wǒ měitiān fābù xīn de zhàopiàn.}\\
« D'abord, je publie chaque jour de nouvelles photos. »\\[0.4em]
然后，我回复顾客的评论。\\
\emph{Ránhòu, wǒ huífù gùkè de pínglùn.}\\
« Ensuite, je réponds aux commentaires des clients. »\\[0.4em]
有时候有人发垃圾信息，我就删除或者屏蔽他们。\\
\emph{Yǒu shíhou yǒu rén fā lājī xìnxī, wǒ jiù shānchú huòzhě píngbì tāmen.}\\
« Parfois, quelqu'un envoie des spams ; je les supprime ou je les bloque. »\\[0.4em]
因为我很认真，所以我们的粉丝越来越多。\\
\emph{Yīnwèi wǒ hěn rènzhēn, suǒyǐ wǒmen de fěnsī yuè lái yuè duō.}\\
« Comme je suis très sérieuse, nos abonnés sont de plus en plus nombreux. »
\end{textebox}

\textbf{Glossaire du texte}\par

\begin{center}
\begin{tabularx}{\linewidth}{|X|X|X|X|}
\hline
\rowcolor{popblueL} Caractères & Pinyin & Nature & Traduction \\
\hline
管理 & guǎnlǐ & verbe & gérer, administrer \\
\hline
网店 & wǎngdiàn & nom & boutique en ligne \\
\hline
账号 & zhànghào & nom & compte (en ligne) \\
\hline
发布 & fābù & verbe & publier, poster \\
\hline
照片 & zhàopiàn & nom & photo \\
\hline
回复 & huífù & verbe & répondre (par écrit) \\
\hline
顾客 & gùkè & nom & client(e) \\
\hline
评论 & pínglùn & nom / verbe & commentaire ; commenter \\
\hline
垃圾信息 & lājī xìnxī & locution nominale & spam, message indésirable \\
\hline
删除 & shānchú & verbe & supprimer, effacer \\
\hline
屏蔽 & píngbì & verbe & bloquer, masquer \\
\hline
认真 & rènzhēn & adjectif & sérieux, appliqué \\
\hline
粉丝 & fěnsī & nom & abonnés, fans \\
\hline
\end{tabularx}
\end{center}

\begin{savaistu}[微信, le « super compte » chinois]
\zh{微信} (\emph{Wēixìn}, WeChat) compte plus d'un milliard d'utilisateurs dans le monde. En Chine, on y fait presque tout : discuter, payer au marché, commander un taxi, suivre une boutique. Au Cameroun, de nombreux commerçants chinois et camerounais l'utilisent aussi pour leurs affaires, par exemple pour commander des marchandises à Guangzhou. Le mot \zh{粉丝} (\emph{fěnsī}, « abonnés ») est un emprunt amusant : à l'origine, il désigne les vermicelles de haricot, mais sa prononciation ressemble à l'anglais \emph{fans} !
\end{savaistu}

\courssub{Traduction phrase par phrase et observation}

Observons maintenant le texte de près. Chaque phrase joue un rôle précis dans l'organisation du récit.

\begin{center}
\begin{tabularx}{\linewidth}{|X|X|X|}
\hline
\rowcolor{popblueL} Phrase chinoise + pinyin & Traduction & Fonction dans le texte \\
\hline
我帮妈妈管理网店的微信账号。 \emph{Wǒ bāng māma guǎnlǐ wǎngdiàn de Wēixìn zhànghào.} & J'aide ma mère à gérer le compte WeChat de la boutique en ligne. & Présentation : qui ? quoi ? \\
\hline
首先，我每天发布新的照片。 \emph{Shǒuxiān, wǒ měitiān fābù xīn de zhàopiàn.} & D'abord, je publie chaque jour de nouvelles photos. & Action quotidienne n°1 \\
\hline
然后，我回复顾客的评论。 \emph{Ránhòu, wǒ huífù gùkè de pínglùn.} & Ensuite, je réponds aux commentaires des clients. & Action quotidienne n°2 \\
\hline
有时候有人发垃圾信息，我就删除或者屏蔽他们。 \emph{Yǒu shíhou yǒu rén fā lājī xìnxī, wǒ jiù shānchú huòzhě píngbì tāmen.} & Parfois, quelqu'un envoie des spams ; je les supprime ou je les bloque. & Problème et solution \\
\hline
因为我很认真，所以我们的粉丝越来越多。 \emph{Yīnwèi wǒ hěn rènzhēn, suǒyǐ wǒmen de fěnsī yuè lái yuè duō.} & Comme je suis très sérieuse, nos abonnés sont de plus en plus nombreux. & Résultat, conclusion \\
\hline
\end{tabularx}
\end{center}

\begin{exemplebox}[Des exemples à retenir par cœur]
\zh{我回复顾客的评论。} \emph{Wǒ huífù gùkè de pínglùn.} = « Je réponds aux commentaires des clients. »\\[0.3em]
\zh{我每天发布新的照片。} \emph{Wǒ měitiān fābù xīn de zhàopiàn.} = « Je publie chaque jour de nouvelles photos. »\\[0.3em]
\zh{我就删除或者屏蔽他们。} \emph{Wǒ jiù shānchú huòzhě píngbì tāmen.} = « Je les supprime ou je les bloque. »
\end{exemplebox}

\courssub{Repérage : les verbes d'action de la gestion}

\begin{methode}[Analyser un texte modèle en trois gestes]
Avant d'imiter un texte, il faut le « démonter » comme un mécanicien démonte un moteur.
\begin{enumerate}
\item \textbf{Souligner les verbes d'action} : ce sont les moteurs du texte. Ici : \zh{管理} (gérer), \zh{发布} (publier), \zh{回复} (répondre), \zh{删除} (supprimer), \zh{屏蔽} (bloquer). Demande-toi : qui fait l'action ? sur quoi ?
\item \textbf{Encadrer les connecteurs} : ce sont les charnières qui relient les phrases. Ici : \zh{首先} (d'abord), \zh{然后} (ensuite), \zh{就} (alors), \zh{因为...所以...} (parce que... donc...).
\item \textbf{Numéroter les étapes du plan} : attribue à chaque phrase une fonction (présentation, action, problème-solution, résultat). Tu obtiens le squelette que tu pourras réutiliser avec ton propre contenu.
\end{enumerate}
\end{methode}

\begin{definition}[Les cinq verbes de la gestion des réseaux sociaux]
\cle{管理} (\emph{guǎnlǐ}) : gérer un compte, une page. — \cle{发布} (\emph{fābù}) : publier un contenu officiel (photo, annonce). — \cle{回复} (\emph{huífù}) : répondre par écrit à un message ou un commentaire. — \cle{删除} (\emph{shānchú}) : supprimer un contenu. — \cle{屏蔽} (\emph{píngbì}) : bloquer un utilisateur ou masquer un contenu.
\end{definition}

\begin{attention}[Ne confonds pas 回复 et 回答 !]
Les francophones traduisent souvent « répondre » par \zh{回答} (\emph{huídá}). Mais \zh{回答} s'emploie pour répondre à une \textbf{question} (à l'oral ou à l'écrit) : \zh{回答问题} (répondre à une question). Pour répondre à un \textbf{commentaire ou un message} sur les réseaux sociaux, on dit \zh{回复} : \zh{回复评论} (répondre à un commentaire), \zh{回复信息} (répondre à un message). Autre piège : \zh{删除} s'écrit avec le radical du couteau 刂 à droite de 删 — ne l'oublie pas, sinon le caractère n'existe pas !
\end{attention}

\courssub{Repérage : les connecteurs du texte}

\begin{propriete}[Quatre connecteurs pour construire un récit simple]
\begin{itemize}
\item \zh{首先} (\emph{shǒuxiān}, « d'abord ») et \zh{然后} (\emph{ránhòu}, « ensuite ») : placés en tête de phrase, suivis d'une virgule chinoise \zh{，}. Structure : \textbf{首先，Sujet + Verbe... 然后，Sujet + Verbe...}
\item \zh{就} (\emph{jiù}, « alors, aussitôt ») : placé \textbf{devant le verbe} de la seconde proposition, jamais en tête de phrase. Structure : \textbf{Proposition 1，Sujet + 就 + Verbe}. Exemple : \zh{有人发垃圾信息，我就删除。} (Quelqu'un envoie un spam, alors je supprime.)
\item \zh{因为...所以...} (\emph{yīnwèi... suǒyǐ...}, « parce que... donc... ») : contrairement au français, le chinois garde volontiers les \textbf{deux} mots dans la même phrase. Structure : \textbf{因为 + cause，所以 + conséquence}.
\end{itemize}
\end{propriete}

\begin{exemplebox}[Comparaison français / chinois]
\zh{因为我很认真，所以粉丝越来越多。} \emph{Yīnwèi wǒ hěn rènzhēn, suǒyǐ fěnsī yuè lái yuè duō.} = « Parce que je suis sérieuse, (donc) les abonnés sont de plus en plus nombreux. » En français, « parce que » et « donc » ne se combinent pas ; en chinois, la paire \zh{因为...所以...} est naturelle et même recommandée à l'écrit.\\[0.3em]
\zh{看到垃圾信息，我就屏蔽他们。} \emph{Kàndào lājī xìnxī, wǒ jiù píngbì tāmen.} = « Quand je vois un spam, je les bloque aussitôt. »
\end{exemplebox}

\courssub{Analyse de la structure du texte}

Le texte modèle suit un plan en quatre étapes, très facile à imiter : \textbf{présentation} (qui fait quoi) → \textbf{actions quotidiennes} (avec 首先 / 然后) → \textbf{problème et solution} (spam → supprimer / bloquer) → \textbf{résultat} (avec 因为...所以...). C'est exactement ce plan que tu réutiliseras dans ta production finale.

\begin{popfigure}
\begin{tikzpicture}[
  node distance=0.55cm,
  etape/.style={rectangle, rounded corners=4pt, draw=popblue, thick, fill=popblueL, align=center, minimum width=4.6cm, minimum height=1.1cm, font=\small},
  fleche/.style={-{Stealth[length=3mm]}, thick, popdark}
]
\node[etape, align=center] (intro){\textbf{Introduction}\\ 我帮妈妈管理网店的微信账号。};
\node[etape, below=of intro, fill=popgreenL, draw=popgreen, align=center] (actions){\textbf{Actions quotidiennes}\\ 首先，我发布照片。然后，我回复评论。};
\node[etape, below=of actions, fill=poporangeL, draw=poporange, align=center] (probleme){\textbf{Problème et solution}\\ 有人发垃圾信息，我就删除或者屏蔽。};
\node[etape, below=of probleme, fill=poppurpleL, draw=poppurple, align=center] (conclu){\textbf{Conclusion : résultat}\\ 因为我很认真，所以粉丝越来越多。};
\draw[fleche] (intro) -- (actions);
\draw[fleche] (actions) -- (probleme);
\draw[fleche] (probleme) -- (conclu);
\end{tikzpicture}
\legende{Organigramme du plan du texte modèle : quatre étapes enchaînées par les connecteurs.}
\end{popfigure}

\begin{exoresolu}[Retrouver le plan d'un mini-texte]
Énoncé : voici un texte en désordre. Remets les phrases dans l'ordre du plan (présentation → actions → problème-solution → résultat) : (a) \zh{所以我们的粉丝越来越多。} (b) \zh{我帮爸爸管理网店的账号。} (c) \zh{有人发垃圾信息，我就删除。} (d) \zh{首先，我发布照片。然后，我回复评论。}
\tcblower
Solution détaillée : la présentation vient toujours en premier : (b) « J'aide mon père à gérer le compte de la boutique. » Viennent ensuite les actions quotidiennes annoncées par \zh{首先} et \zh{然后} : (d). La phrase (c) décrit le problème (le spam) et la solution (\zh{删除}), introduite par \zh{就}. Enfin, (a) commence par \zh{所以} : c'est la conséquence, donc la conclusion. Ordre correct : \textbf{b → d → c → a}. Texte reconstitué : \zh{我帮爸爸管理网店的账号。首先，我发布照片。然后，我回复评论。有人发垃圾信息，我就删除。所以我们的粉丝越来越多。} (\emph{Wǒ bāng bàba guǎnlǐ wǎngdiàn de zhànghào. Shǒuxiān, wǒ fābù zhàopiàn. Ránhòu, wǒ huífù pínglùn. Yǒu rén fā lājī xìnxī, wǒ jiù shānchú. Suǒyǐ wǒmen de fěnsī yuè lái yuè duō.}) — « J'aide mon père à gérer le compte de la boutique. D'abord, je publie des photos. Ensuite, je réponds aux commentaires. Quelqu'un envoie un spam, je le supprime. Donc nos abonnés sont de plus en plus nombreux. »
\end{exoresolu}

\begin{exercice}[À ton tour : observe et repère]
Dans le texte modèle \zh{我管理妈妈的网店账号} : 1) souligne les cinq verbes d'action et nomme leur complément (on gère \emph{quoi} ? on publie \emph{quoi} ?) ; 2) encadre les quatre connecteurs et précise leur place dans la phrase ; 3) recopie la phrase de conclusion et traduis-la mot à mot. Le corrigé sera construit en classe à partir du tableau de traduction phrase par phrase.
\end{exercice}

\begin{aretenir}[L'essentiel de cette section]
Un bon texte simple en chinois repose sur trois piliers : des \cle{verbes d'action} précis (\zh{管理、发布、回复、删除、屏蔽}), des \cle{connecteurs} bien placés (\zh{首先、然后} en tête de phrase, \zh{就} devant le verbe, la paire \zh{因为...所以...}), et un \cle{plan en quatre étapes} : présentation → actions quotidiennes → problème et solution → résultat. Garde ce squelette en tête : il te servira pour toutes tes productions écrites sur les réseaux sociaux.
\end{aretenir}

\coursec{Lexique de la gestion des réseaux sociaux}

Dans la section précédente, nous avons lu et commenté un modèle de texte : un court article dans lequel un élève de Première raconte comment il gère son compte sur les réseaux sociaux. Vous avez remarqué que ce texte reposait sur un petit nombre de verbes très fréquents : \zh{发布}, \zh{分享}, \zh{评论}, \zh{回复}, \zh{管理}, \zh{删除}… C'est exactement ce vocabulaire que nous allons maintenant étudier en profondeur, mot par mot, afin que vous puissiez l'employer sans faute dans vos propres productions écrites. Objectif de la section : connaître le sens exact de chaque mot, savoir avec quels noms il s'emploie, et pouvoir le réutiliser dans une phrase correcte.

\courssub{Observation : le vocabulaire en contexte}

Avant toute explication, observons ces mots dans une situation authentique. Lisez le court texte suivant, puis repérez les verbes de la leçon.

\begin{textebox}[Un dimanche sur mon téléphone]
今天是星期天。早上，我打开手机，看到我的账号有很多新消息。\\
Jīntiān shì xīngqītiān. Zǎoshang, wǒ dǎkāi shǒujī, kàndào wǒ de zhànghào yǒu hěn duō xīn xiāoxī.\\
« Aujourd'hui, c'est dimanche. Le matin, j'ouvre mon téléphone et je vois que mon compte a beaucoup de nouveaux messages. »\\[4pt]
昨天我发布了一张照片，很多朋友点赞，还有人评论：「真好看！」\\
Zuótiān wǒ fābù le yì zhāng zhàopiàn, hěn duō péngyou diǎnzàn, hái yǒu rén pínglùn : « Zhēn hǎokàn ! »\\
« Hier, j'ai publié une photo ; beaucoup d'amis ont liké, et quelqu'un a commenté : "C'est vraiment beau !" »\\[4pt]
我马上回复了他们的评论，也分享了朋友的视频。\\
Wǒ mǎshàng huífù le tāmen de pínglùn, yě fēnxiǎng le péngyou de shìpín.\\
« J'ai tout de suite répondu à leurs commentaires, et j'ai aussi partagé la vidéo d'un ami. »\\[4pt]
但是有一个账号总是发垃圾信息，我删除了那些信息，还屏蔽了那个账号。\\
Dànshì yǒu yí ge zhànghào zǒngshì fā lājī xìnxī, wǒ shānchú le nàxiē xìnxī, hái píngbì le nàge zhànghào.\\
« Mais un compte envoyait toujours des spams ; j'ai supprimé ces messages et j'ai bloqué ce compte. »\\[4pt]
现在，我的账号有两百个粉丝。管理账号不容易，但是很有意思！\\
Xiànzài, wǒ de zhànghào yǒu liǎng bǎi ge fěnsī. Guǎnlǐ zhànghào bù róngyì, dànshì hěn yǒu yìsi !\\
« Maintenant, mon compte a deux cents abonnés. Gérer un compte n'est pas facile, mais c'est très intéressant ! »
\end{textebox}

\textbf{Questions de compréhension.}
\begin{enumerate}
\item Que fait le narrateur quand il ouvre son téléphone le matin ?
\item Quelles réactions sa photo a-t-elle reçues ?
\item Comment le narrateur traite-t-il le compte qui envoie des spams ?
\item Combien d'abonnés le narrateur a-t-il ?
\end{enumerate}

\textbf{Glossaire rapide.} \zh{消息} xiāoxī (message, nouvelle) ; \zh{照片} zhàopiàn (photo) ; \zh{视频} shìpín (vidéo) ; \zh{总是} zǒngshì (toujours) ; \zh{容易} róngyì (facile) ; \zh{打开} dǎkāi (ouvrir, allumer).

\courssub{Tableau récapitulatif du lexique}

Voici l'ensemble du vocabulaire de la leçon, classé par fonction. Apprenez-le par groupes, pas dans le désordre : la mémoire retient mieux les familles que les listes mélangées.

\begin{center}
\begin{tabularx}{\linewidth}{|X|X|X|X|}
\hline
\rowcolor{popblueL} Caractères & Pinyin & Nature & Traduction \\
\hline
发布 & fābù & verbe & publier, diffuser \\
\hline
分享 & fēnxiǎng & verbe & partager \\
\hline
转发 & zhuǎnfā & verbe & republier, transférer \\
\hline
点赞 & diǎnzàn & verbe & liker, mettre un « j'aime » \\
\hline
评论 & pínglùn & verbe / nom & commenter / commentaire \\
\hline
回复 & huífù & verbe & répondre (à un message) \\
\hline
关注 & guānzhù & verbe & suivre (un compte), s'abonner \\
\hline
粉丝 & fěnsī & nom & abonné, fan \\
\hline
管理 & guǎnlǐ & verbe & gérer, administrer \\
\hline
账号 & zhànghào & nom & compte (d'utilisateur) \\
\hline
删除 & shānchú & verbe & supprimer, effacer \\
\hline
屏蔽 & píngbì & verbe & bloquer, masquer \\
\hline
垃圾信息 & lājī xìnxī & nom & spam, message indésirable \\
\hline
\end{tabularx}
\end{center}

\begin{definition}[屏蔽 píngbì]
\zh{屏蔽} píngbì est un verbe qui signifie \cle{bloquer un utilisateur} ou \cle{masquer un contenu} sur un réseau social : la personne bloquée ne peut plus voir vos publications ni vous envoyer de messages. Le caractère \zh{屏} évoque un « écran, une cloison » et \zh{蔽} signifie « cacher » : littéralement, « cacher derrière un écran ».
\end{definition}

\courssub{Étude détaillée des verbes : emplois et collocations}

\textbf{1. Le groupe « publier » : 发布, 分享, 转发.}

\zh{发布} fābù est le verbe officiel de la publication : on publie une photo, un texte, une annonce. \zh{分享} fēnxiǎng insiste sur le fait de \emph{partager avec les autres} (une vidéo, un lien, une expérience) ; on peut préciser le destinataire avec \zh{和} : \zh{和大家分享} (partager avec tout le monde). \zh{转发} zhuǎnfā, composé de \zh{转} (tourner, transmettre) et \zh{发} (envoyer), signifie « republier » le contenu de quelqu'un d'autre sur son propre compte.

\begin{exemplebox}[Trois verbes, trois actions]
\zh{我每天早上发布一张照片。} \emph{Wǒ měitiān zǎoshang fābù yì zhāng zhàopiàn.} « Je publie une photo chaque matin. »\\
\zh{她想和大家分享她的旅行。} \emph{Tā xiǎng hé dàjiā fēnxiǎng tā de lǚxíng.} « Elle veut partager son voyage avec tout le monde. »\\
\zh{请转发这个消息，谢谢！} \emph{Qǐng zhuǎnfā zhège xiāoxī, xièxie !} « Merci de transférer ce message ! »
\end{exemplebox}

\textbf{2. Le groupe « interagir » : 点赞, 评论, 回复.}

\zh{点赞} diǎnzàn : \zh{点} = « cliquer, pointer », \zh{赞} = « louer » ; c'est le « like ». \zh{评论} pínglùn peut être un verbe (commenter) ou un nom (un commentaire) — attention à son rôle dans la phrase : dans \zh{有人评论}, c'est un verbe ; dans \zh{回复评论}, c'est un nom, complément du verbe. \zh{回复} huífù signifie « répondre » et se construit directement avec un complément d'objet : \zh{回复评论} (répondre aux commentaires), \zh{回复消息} (répondre à un message).

\textbf{3. Le groupe « gérer » : 管理, 删除, 屏蔽.}

\zh{管理} guǎnlǐ (gérer) s'emploie pour un compte, une page, un groupe : \zh{管理账号}. \zh{删除} shānchú (supprimer) s'emploie pour un contenu : \zh{删除信息}, \zh{删除照片}. \zh{屏蔽} píngbì (bloquer) s'emploie pour une personne ou un compte : \zh{屏蔽那个账号}. Retenez ce partage des rôles : on \emph{gère} un compte, on \emph{supprime} un contenu, on \emph{bloque} une personne.

\textbf{4. Le groupe « communauté » : 关注, 粉丝.}

\zh{关注} guānzhù signifie « suivre, s'abonner à » un compte ; la personne qui vous suit est votre \zh{粉丝} fěnsī.

\begin{exemplebox}[Exemples traduits]
\zh{请关注我的账号！} \emph{Qǐng guānzhù wǒ de zhànghào !} « Suivez mon compte, s'il vous plaît ! »\\
\zh{她的账号有五千个粉丝。} \emph{Tā de zhànghào yǒu wǔ qiān ge fěnsī.} « Son compte a cinq mille abonnés. »\\
\zh{老师管理我们学校的中文网页。} \emph{Lǎoshī guǎnlǐ wǒmen xuéxiào de Zhōngwén wǎngyè.} « Le professeur gère la page web en chinois de notre école. »\\
\zh{我删除了那条垃圾信息。} \emph{Wǒ shānchú le nà tiáo lājī xìnxī.} « J'ai supprimé ce spam. »
\end{exemplebox}

\begin{attention}[粉丝 ne se traduit pas littéralement !]
\zh{粉丝} fěnsī désigne à l'origine les « vermicelles (de haricot mungo) » ! Dans le sens qui nous occupe, c'est en réalité un \cle{emprunt phonétique} à l'anglais \emph{fans} : les syllabes fěn-sī imitent le son anglais. Ne traduisez donc jamais \zh{我有很多粉丝} par « J'ai beaucoup de vermicelles » : cela signifie « J'ai beaucoup d'abonnés ». De même, \zh{垃圾信息} se compose de \zh{垃圾} (ordures) + \zh{信息} (information) : littéralement « informations-ordures », c'est-à-dire le spam.
\end{attention}

\courssub{Les collocations à connaître par cœur}

En chinois comme en français, certains verbes « vont avec » certains noms. Ce sont les \cle{collocations}. Les quatre collocations officielles de cette leçon :

\begin{center}
\begin{tabularx}{\linewidth}{|X|X|X|}
\hline
\rowcolor{popblueL} Collocation & Pinyin & Traduction \\
\hline
发布照片 & fābù zhàopiàn & publier des photos \\
\hline
回复评论 & huífù pínglùn & répondre aux commentaires \\
\hline
管理账号 & guǎnlǐ zhànghào & gérer un compte \\
\hline
删除信息 & shānchú xìnxī & supprimer un message \\
\hline
\end{tabularx}
\end{center}

\begin{propriete}[Structure : verbe + complément d'objet]
La structure de base est : \textbf{Sujet + Verbe + (了) + Complément}. Le verbe se place toujours \emph{avant} son complément, comme en français, mais sans préposition : on dit \zh{回复评论} (répondre + commentaires), jamais \zh{回复对评论}. Le chinois n'utilise pas de préposition « à » devant le complément d'objet direct. La particule \zh{了}, placée juste après le verbe, marque une action accomplie : \zh{我删除了那条信息} « J'ai supprimé ce message ». Elle ne s'emploie pas pour les actions habituelles : \zh{我每天早上发布照片} (sans \zh{了}).
\end{propriete}

\begin{attention}[Erreurs fréquentes des francophones]
\begin{itemize}
\item Ne dites pas \zh{回答评论} pour « répondre aux commentaires » : \zh{回答} huídá sert à répondre à une \emph{question} ; pour un message ou un commentaire, utilisez \zh{回复} huífù.
\item \zh{看} kàn (regarder) ne remplace pas \zh{关注} guānzhù : « suivre un compte » ne se dit jamais \zh{看账号}.
\item N'oubliez pas le classificateur : \zh{一张照片} (une photo), \zh{一条信息} (un message), \zh{一个账号} (un compte), \zh{一个粉丝} (un abonné).
\item Attention aux tons : \zh{fābù} (1\textsuperscript{er} puis 4\textsuperscript{e} ton), \zh{huífù} (2\textsuperscript{e} puis 4\textsuperscript{e} ton), \zh{guǎnlǐ} (deux 3\textsuperscript{e} tons, le premier se prononçant alors comme un 2\textsuperscript{e} ton à l'oral).
\end{itemize}
\end{attention}

\courssub{Carte mentale du lexique}

Pour mémoriser durablement, organisez les mots autour du mot central \zh{账号} (le compte), en quatre familles d'actions.

\begin{popfigure}
\begin{tikzpicture}[
  mindmap,
  grow cyclic,
  every node/.style={concept, font=\small, align=center},
  concept color=popblue!60,
  level 1 concept/.append style={level distance=3.4cm, sibling angle=90, font=\small},
  level 2 concept/.append style={level distance=2.2cm, font=\footnotesize},
]
\node[concept color=poporange!70, align=center]{账号\\ zhànghào\\ le compte}
  child[concept color=popgreen!60] { node[align=center]{发布 publier\\ 分享 partager\\ 转发 republier} }
  child[concept color=poppurple!60] { node[align=center]{点赞 liker\\ 评论 commenter\\ 回复 répondre} }
  child[concept color=poppink!60] { node[align=center]{管理 gérer\\ 删除 supprimer\\ 屏蔽 bloquer} }
  child[concept color=popgold!70] { node[align=center]{关注 suivre\\ 粉丝 abonné} };
\end{tikzpicture}
\legende{Carte mentale du lexique : quatre familles d'actions autour du compte 账号.}
\end{popfigure}

\begin{methode}[Mémoriser avec les icônes du téléphone]
Pour retenir chaque verbe, associez-le mentalement à l'icône correspondante sur votre téléphone :
\begin{enumerate}
\item Le pouce levé (like) = \zh{点赞} diǎnzàn : imaginez que vous « pointez » (\zh{点}) le pouce.
\item La corbeille = \zh{删除} shānchú : on « efface » (\zh{删}) en jetant à la poubelle.
\item La flèche de partage = \zh{分享} fēnxiǎng ou \zh{转发} zhuǎnfā selon que l'on partage ou republie.
\item Le crayon de publication = \zh{发布} fābù.
\item Le symbole « interdit » = \zh{屏蔽} píngbì : on bloque l'utilisateur indésirable.
\end{enumerate}
Répétez chaque paire icône-verbe à voix haute trois fois, puis testez-vous sans regarder la liste.
\end{methode}

\courssub{Exercice résolu et entraînement}

\begin{exoresolu}[Compléter avec le bon verbe]
Énoncé. Complétez chaque phrase avec le verbe qui convient : \zh{发布}, \zh{回复}, \zh{删除}, \zh{关注}, \zh{屏蔽}.\\
1. 请\_\_\_\_我的账号！（Suivez mon compte !）\\
2. 他昨天\_\_\_\_了一张照片。（Il a publié une photo hier.）\\
3. 我马上\_\_\_\_了你的评论。（J'ai tout de suite répondu à ton commentaire.）\\
4. 这条信息没有用，请\_\_\_\_它。（Ce message est inutile, supprime-le.）\\
5. 他总是发垃圾信息，我要\_\_\_\_他。（Il envoie toujours des spams, je vais le bloquer.）
\tcblower
Solution détaillée.\\
1. \zh{请关注我的账号！} — « suivre un compte » = \zh{关注}.\\
2. \zh{他昨天发布了一张照片。} — publier une photo = \zh{发布照片} ; notez le classificateur \zh{张} et la marque de l'action accomplie \zh{了} après le verbe.\\
3. \zh{我马上回复了你的评论。} — répondre à un commentaire = \zh{回复评论} (et non \zh{回答}).\\
4. \zh{这条信息没有用，请删除它。} — supprimer un message = \zh{删除信息}.\\
5. \zh{他总是发垃圾信息，我要屏蔽他。} — bloquer une personne = \zh{屏蔽}.
\end{exoresolu}

\begin{exercice}[À vous de jouer]
Traduisez en chinois, avec le pinyin : 1. « Elle partage souvent des vidéos. » 2. « Mon compte a cent abonnés. » 3. « Le professeur gère la page de l'école. » 4. « Ne supprime pas ce message ! »
\end{exercice}

\begin{corrige}[Exercice « À vous de jouer »]
1. \zh{她常常分享视频。} \emph{Tā chángcháng fēnxiǎng shìpín.}\\
2. \zh{我的账号有一百个粉丝。} \emph{Wǒ de zhànghào yǒu yì bǎi ge fěnsī.}\\
3. \zh{老师管理学校的网页。} \emph{Lǎoshī guǎnlǐ xuéxiào de wǎngyè.}\\
4. \zh{别删除这条信息！} \emph{Bié shānchú zhè tiáo xìnxī !} (pour une interdiction, on utilise \zh{别} bié, et non \zh{不}.)
\end{corrige}

\begin{aretenir}[L'essentiel de la section]
\begin{itemize}
\item Quatre familles de verbes autour de \zh{账号} : publier (\zh{发布}, \zh{分享}, \zh{转发}), interagir (\zh{点赞}, \zh{评论}, \zh{回复}), gérer (\zh{管理}, \zh{删除}, \zh{屏蔽}), communauté (\zh{关注}, \zh{粉丝}).
\item Collocations obligatoires : \zh{发布照片}, \zh{回复评论}, \zh{管理账号}, \zh{删除信息}.
\item Structure : Sujet + Verbe + (了) + Complément, sans préposition ; \zh{了} seulement pour une action accomplie.
\item Pièges : \zh{粉丝} = abonnés (pas vermicelles) ; \zh{回复} pour les messages, \zh{回答} pour les questions ; toujours un classificateur devant le nom (\zh{张}, \zh{条}, \zh{个}).
\end{itemize}
\end{aretenir}

Ce lexique est maintenant votre boîte à outils. Dans la section suivante, nous passerons à l'écriture correcte de ces caractères : traits, radicaux et caractères proches à ne pas confondre.

\coursec{Écriture correcte des caractères}

Dans la section précédente, nous avons découvert le lexique de la gestion des réseaux sociaux : \zh{发布} (fābù, publier), \zh{分享} (fēnxiǎng, partager), \zh{评论} (pínglùn, commenter), \zh{管理} (guǎnlǐ, gérer), \zh{账号} (zhànghào, compte), \zh{删除} (shānchú, supprimer), \zh{屏蔽} (píngbì, bloquer/masquer). Savoir prononcer et comprendre ces mots ne suffit pas : pour réussir l'expression écrite, il faut pouvoir les \cle{écrire} correctement, trait par trait. C'est l'objet de cette section : observer la structure de chaque caractère, appliquer les règles d'ordre des traits, et éviter les confusions entre caractères proches.

\courssub{Observer avant d'écrire : la logique de construction}

\begin{exemplebox}[Observer la structure des caractères du thème]
Regardez attentivement ces mots et repérez les parties qui se répètent :\\
\zh{评论} pínglùn « commentaire » — \zh{管理} guǎnlǐ « gérer » — \zh{账号} zhànghào « compte » — \zh{删除} shānchú « supprimer »\\
On remarque que la plupart des caractères sont composés de \cle{deux parties} : une partie à gauche et une partie à droite (评, 理, 账, 删), ou une partie en haut et une partie en bas (管, 蔽, 享). La partie récurrente est souvent le \cle{radical} (clé), qui donne un indice sur le sens ou la prononciation.
\end{exemplebox}

\begin{propriete}[Règle générale de l'ordre des traits]
L'écriture chinoise suit des règles strictes, valables pour tous les caractères :
\begin{enumerate}
\item On écrit de \cle{haut en bas} et de \cle{gauche à droite}.
\item Pour un caractère gauche-droite, on écrit \cle{toujours la partie de gauche (souvent le radical) avant la partie de droite} : 评 = 讠 puis 平 ; 理 = 王 puis 里 ; 账 = 贝 puis 长 ; 删 = 册 puis 刂.
\item Pour un caractère haut-bas, on écrit \cle{la partie haute avant la partie basse} : 管 = ⺮ puis 官 ; 蔽 = 艹 puis 敝 ; 享 = 亠 puis 口 puis 子.
\item Le trait horizontal (一) s'écrit avant le trait vertical (丨) quand ils se croisent ; on « ferme » un cadre (comme 口, 日, 田) par son dernier trait horizontal du bas.
\item Le point (丶) en haut à gauche s'écrit en premier ; le point en bas à droite s'écrit en dernier (ex. : 发, 我).
\end{enumerate}
\end{propriete}

\courssub{Étude caractère par caractère}

\textbf{1. 发 (fā) — envoyer, publier : 5 traits, radical 又}\par
Ordre des traits : ① 丿 (piě), ② ㇇ (héng zhé gōu), ③ 丶 (diǎn), ④ 丿 (piě), ⑤ 丶 (diǎn). Attention : le caractère se termine par un \cle{point final} (丶) en bas à droite, que les élèves oublient souvent. La version traditionnelle est 發 ; au Cameroun comme en Chine continentale, on utilise uniquement la forme simplifiée 发.
Exemple : \zh{我今天发了一张照片。} \emph{Wǒ jīntiān fā le yì zhāng zhàopiàn.} « J'ai publié une photo aujourd'hui. »

\textbf{2. 享 (xiǎng) — jouir, partager (分享) : 8 traits, radical 亠}\par
Structure haut-bas en trois étages. On écrit de haut en bas : d'abord 亠 (2 traits : 一, 亅), puis 口 (3 traits : 丨, ㇆, 一), puis 子 (3 traits : 一, ㇆, 一). Trois étages bien empilés, sans les faire déborder.
Exemple : \zh{她喜欢和朋友分享视频。} \emph{Tā xǐhuan hé péngyou fēnxiǎng shìpín.} « Elle aime partager des vidéos avec ses amis. »

\textbf{3. 评 (píng) — commenter : 7 traits, radical 讠 (parole)}\par
Structure gauche-droite. Le radical 讠, à gauche, s'écrit en \cle{deux traits} (le point 丶 puis le trait brisé ㇊) ; ensuite seulement on écrit 平 à droite (5 traits : 一, 丨, 一, 一, 一). Le radical « parole » rappelle que commenter, c'est parler.
Exemple : \zh{请删除这个评论。} \emph{Qǐng shānchú zhège pínglùn.} « Supprime ce commentaire, s'il te plaît. »

\textbf{4. 管 (guǎn) — gérer, administrer : 14 traits, radical ⺮ (bambou)}\par
Caractère complexe haut-bas. En haut, la clé du bambou ⺮ (6 traits, écrits en deux groupes de trois : 丿 丶 丿 / 丿 丶 丿) ; en bas, 官 (8 traits : 宀 3 traits, puis ㇆, 一, 丨, 一, 丨). On écrit d'abord tout le haut (⺮), puis tout le bas (官).

\textbf{5. 理 (lǐ) — raison, ordre (管理) : 11 traits, radical 王 (jade)}\par
Structure gauche-droite. À gauche, la clé du jade 王 (4 traits : 三 horizontaux, puis le trait montant final ㇀) ; à droite, 里 (7 traits : 日 4 traits — 丨, ㇆, 一, 一 — puis 土 3 traits : 一, 丨, 一). Attention : dans 里, le trait vertical central de 土 s'écrit \emph{après} le 日 et \emph{avant} les deux horizontaux du bas.

\textbf{6. 账 (zhàng) — compte, comptabilité : 8 traits, radical 贝 (coquillage/argent)}\par
Structure gauche-droite. À gauche 贝 (4 traits : 丿, ㇆, 丶, 丶), la clé de l'argent — autrefois les coquillages servaient de monnaie ; à droite 长 (4 traits : 丿, 一, 丨, 乀). Ne pas confondre avec \zh{帐} (zhàng, tente, moustiquaire), qui prend la clé de l'étoffe 巾. Un \zh{账号} (compte en ligne) s'écrit avec 贝, car il concerne l'argent et les registres.

\textbf{7. 删 (shān) — supprimer : 7 traits, radical 刂 (couteau)}\par
Structure gauche-droite. À gauche 册 (5 traits : 丨, ㇆, 一, 一, 一 — comme un registre de tablettes liées), à droite le couteau 刂 (2 traits : 丿, 丨), qui « coupe » le texte. On écrit 册 d'abord, puis le couteau.
Exemple : \zh{老师删掉了不好的评论。} \emph{Lǎoshī shāndiào le bù hǎo de pínglùn.} « Le professeur a supprimé les mauvais commentaires. »

\textbf{8. 蔽 (bì) — masquer, cacher (屏蔽) : 14 traits, radical 艹 (herbe)}\par
Structure haut-bas. On écrit de haut en bas : d'abord l'herbe 艹 (3 traits : 一, 丨, 一), puis la partie 敝 en dessous (11 traits, de gauche à droite : ㇆, 丶, 丿, 丶, 丨, ㇆, 一, 一, 一, 丨, 一). L'image est belle : l'herbe qui \cle{recouvre} et cache ce qui est dessous.

\begin{popfigure}
\begin{tikzpicture}[node distance=0.55cm, every node/.style={font=\small, align=center}]
\node[draw=popblue, fill=popblueL, rounded corners, minimum width=1.6cm, minimum height=0.9cm] (guan) {管 guǎn};
\node[draw=popgreen, fill=popgreenL, rounded corners, below left=0.55cm and 0.2cm of guan, minimum width=1.2cm, align=center] (zhu){⺮ (haut)\\6 traits};
\node[draw=popgreen, fill=popgreenL, rounded corners, below right=0.55cm and 0.2cm of guan, minimum width=1.2cm, align=center] (guan2){官 (bas)\\8 traits};
\draw[-{Stealth}, thick, popdark] (guan) -- (zhu);
\draw[-{Stealth}, thick, popdark] (guan) -- (guan2);
\node[draw=poporange, fill=poporangeL, rounded corners, minimum width=1.6cm, minimum height=0.9cm, right=2.2cm of guan] (ping) {评 píng};
\node[draw=poppurple, fill=poppurpleL, rounded corners, below left=0.55cm and 0.2cm of ping, minimum width=1.2cm, align=center] (yan){讠 (gauche)\\2 traits};
\node[draw=poppurple, fill=poppurpleL, rounded corners, below right=0.55cm and 0.2cm of ping, minimum width=1.2cm, align=center] (ping2){平 (droite)\\5 traits};
\draw[-{Stealth}, thick, popdark] (ping) -- (yan);
\draw[-{Stealth}, thick, popdark] (ping) -- (ping2);
\node[draw=poppink, fill=poppinkL, rounded corners, minimum width=1.6cm, minimum height=0.9cm, right=2.2cm of ping] (zhang) {账 zhàng};
\node[draw=popgold, fill=popgoldL, rounded corners, below left=0.55cm and 0.2cm of zhang, minimum width=1.2cm, align=center] (bei){贝 (gauche)\\4 traits};
\node[draw=popgold, fill=popgoldL, rounded corners, below right=0.55cm and 0.2cm of zhang, minimum width=1.2cm, align=center] (chang){长 (droite)\\4 traits};
\draw[-{Stealth}, thick, popdark] (zhang) -- (bei);
\draw[-{Stealth}, thick, popdark] (zhang) -- (chang);
\end{tikzpicture}
\legende{Décomposition des caractères : flèche vers la partie que l'on écrit en premier.}
\end{popfigure}

\courssub{Tableau récapitulatif des caractères du thème}

\begin{center}
\begin{tabularx}{\linewidth}{|X|X|X|X|}
\hline
\rowcolor{popblueL} Caractère & Pinyin / traits & Radical & Traduction et mot-clé \\
\hline
发 & fā, 5 traits & 又 & publier : 发布 fābù \\
\hline
享 & xiǎng, 8 traits & 亠 & partager : 分享 fēnxiǎng \\
\hline
评 & píng, 7 traits & 讠 (parole) & commenter : 评论 pínglùn \\
\hline
管 & guǎn, 14 traits & ⺮ (bambou) & gérer : 管理 guǎnlǐ \\
\hline
理 & lǐ, 11 traits & 王 (jade) & ordre, raison : 管理 \\
\hline
账 & zhàng, 8 traits & 贝 (argent) & compte : 账号 zhànghào \\
\hline
删 & shān, 7 traits & 刂 (couteau) & supprimer : 删除 shānchú \\
\hline
蔽 & bì, 14 traits & 艹 (herbe) & masquer : 屏蔽 píngbì \\
\hline
\end{tabularx}
\end{center}

\courssub{Caractères proches à ne pas confondre}

\begin{center}
\begin{tabularx}{\linewidth}{|X|X|X|X|}
\hline
\rowcolor{popblueL} Paires & Pinyin & Différence graphique & Traduction \\
\hline
己 / 已 / 巳 & jǐ / yǐ / sì & 3\textsuperscript{e} trait : ouvert / à moitié fermé / fermé & soi-même / déjà / (6\textsuperscript{e} branche terrestre) \\
\hline
未 / 末 & wèi / mò & 2\textsuperscript{e} trait horizontal : court / long & pas encore / fin, bout \\
\hline
辩 / 辨 & biàn / biàn & Radical : 讠 (gauche) / 刂 (droite) & débattre, argumenter / distinguer, discerner \\
\hline
账 / 帐 & zhàng / zhàng & Radical : 贝 (argent) / 巾 (étoffe) & compte, facture / tente, rideau \\
\hline
\end{tabularx}
\end{center}

\begin{attention}[Erreurs fréquentes des francophones]
\begin{itemize}
\item Oublier le \cle{point final} (5\textsuperscript{e} trait) de 发 : sans lui, le caractère est incomplet et considéré comme faux à l'examen.
\item Confondre \zh{未} (wèi, « pas encore », deuxième trait horizontal \emph{court}, ne dépasse pas le premier) et \zh{末} (mò, « fin », deuxième trait horizontal \emph{long}, dépasse le premier). Moyen mnémotechnique : « pas encore » (未) → l'avenir est court devant soi ; « la fin » (末) → tout s'allonge vers la fin.
\item Écrire \zh{账号} avec \zh{帐} : le compte concerne l'argent/la comptabilité, donc la clé \zh{贝} (coquillage/monnaie), pas \zh{巾} (étoffe).
\item Écrire la partie droite avant le radical de gauche dans 评, 理, 账, 删 : c'est une faute d'ordre des traits (règle « gauche avant droite »), sanctionnée à l'examen.
\item Confondre \zh{辩} (débattre, radical 讠) et \zh{辨} (distinguer le vrai du faux, radical 刂) : sur les réseaux sociaux, on \zh{辨别} (biànbié) les vraies et fausses infos.
\end{itemize}
\end{attention}

\begin{methode}[Comment apprendre à écrire 管 (14 traits) et 蔽 (14 traits)]
\begin{enumerate}
\item \textbf{Décompose} : 管 = ⺮ (haut) + 官 (bas) ; 蔽 = 艹 (haut) + 敝 (bas).
\item \textbf{Trace dans l'air} avec le doigt en comptant les traits à voix haute : pour 管, « 1-2-3-4-5-6 » (bambou), « 7-8-9-10-11-12-13-14 » (officiel) ; pour 蔽, « 1-2-3 » (herbe), « 4 à 14 » (bas).
\item \textbf{Vérifie} le nombre de traits par composant (6+8=14 ; 3+11=14).
\item \textbf{Écris sur le cahier}, lentement, trois fois de suite, en comptant toujours.
\item \textbf{Contextualise} : recopie le mot entier 管理 et 屏蔽, puis une phrase complète : \zh{我管理学校的网站，屏蔽了不良账号。} \emph{Wǒ guǎnlǐ xuéxiào de wǎngzhàn, píngbì le bùliáng zhànghào.} « Je gère le site web de l'école et j'ai bloqué les comptes indésirables. »
\end{enumerate}
\end{methode}

\begin{savaistu}[Pourquoi un couteau dans 删 « supprimer » ?]
Avant l'invention du papier, les Chinois écrivaient sur des tablettes de bambou ou de bois liées ensemble par des cordons — c'est le caractère \zh{册} (cè), qu'on retrouve à gauche de 删. Pour corriger une erreur, le scribe \cle{grattait} la surface de la tablette avec un petit couteau : c'est le radical \zh{刂} (lìdāo, « couteau debout ») à droite. « Supprimer », c'était littéralement « couper/gratter le texte » ! Aujourd'hui, sur ton téléphone à Yaoundé ou à Douala, tu appuies sur un bouton « supprimer », mais le caractère garde la mémoire du couteau des scribes.
\end{savaistu}

\begin{exoresolu}[Remettre les traits dans l'ordre et choisir le bon caractère]
Énoncé : 1) Indique l'ordre d'écriture des parties de 评 et de 管. 2) Choisis le bon caractère : mon (账/帐) de réseau social ; (未/末) « pas encore » dans \zh{未来} « le futur » ; on doit (辩/辨) les rumeurs sur Internet.
\tcblower
Solution : 1) 评 : on écrit d'abord le radical de parole 讠 à gauche (2 traits), puis 平 à droite (5 traits) — règle « gauche avant droite ». 管 : on écrit d'abord la clé du bambou ⺮ en haut (6 traits), puis 官 en bas (8 traits) — règle « haut avant bas ». 2) Un compte se dit \zh{账号} avec 贝 (argent) ; « le futur » se dit \zh{未来} (wèilái) avec 未 (deuxième trait court) ; distinguer/démêler les rumeurs utilise \zh{辨别} (biànbié) avec le radical 刂 (couteau qui sépare le vrai du faux).
\end{exoresolu}

\begin{exercice}[À toi d'écrire]
1) Écris cinq fois les caractères 发, 享, 评 en comptant les traits à voix haute.\\
2) Décompose 理, 账, 删, 蔽 en radical + partie restante, et indique l'ordre d'écriture (gauche puis droite / haut puis bas).\\
3) Entoure le bon caractère : (未/末)来 ; 网(账/帐)号 ; 分(辩/辨) les vraies et fausses informations.\\
4) Recopie la phrase : \zh{请删除这个评论，屏蔽那个账号。} \emph{Qǐng shānchú zhège pínglùn, píngbì nàge zhànghào.} « Supprime ce commentaire et bloque ce compte, s'il te plaît. »
\end{exercice}

\begin{aretenir}[L'essentiel de la section]
Pour écrire correctement : décompose chaque caractère, identifie le radical (il donne le sens), respecte l'ordre « gauche avant droite, haut avant bas », compte les traits, et méfie-toi des paires 未/末, 账/帐, 己/已/巳, 辩/辨. Un caractère bien écrit est un caractère bien \cle{compris}.
\end{aretenir}

\coursec{Structure et connecteurs du texte}

Dans la section précédente, nous avons appris à écrire correctement les caractères complexes du thème des réseaux sociaux (\zh{管理}, \zh{发布}, \zh{粉丝}, \zh{删除}...). Mais écrire de beaux caractères ne suffit pas : un bon texte chinois, comme un bon texte français, doit être \cle{structuré} et \cle{articulé}. C'est ce que nous allons découvrir maintenant : le plan type d'un texte de gestion de réseaux sociaux et les connecteurs qui relient les idées entre elles.

\courssub{Observer d'abord : un texte bien structuré}

Lisons d'abord le texte de Marlyse, une élève de Première à Douala qui gère la page du club de football de son lycée. Observons comment les idées s'enchaînent.

\begin{textebox}[我帮学校管理网页 — Wǒ bāng xuéxiào guǎnlǐ wǎngyè]
我帮我们学校管理网页。我非常喜欢这个工作。\\
Wǒ bāng wǒmen xuéxiào guǎnlǐ wǎngyè. Wǒ fēicháng xǐhuan zhège gōngzuò.\\
J'aide notre école à gérer sa page web. J'aime beaucoup ce travail.\\[4pt]
首先，我每天发布照片和新闻；然后，我回复朋友和粉丝的评论；最后，我检查网页的内容。\\
Shǒuxiān, wǒ měitiān fābù zhàopiàn hé xīnwén ; ránhòu, wǒ huífù péngyou hé fěnsī de pínglùn ; zuìhòu, wǒ jiǎnchá wǎngyè de nèiróng.\\
D'abord, je publie chaque jour des photos et des nouvelles ; ensuite, je réponds aux commentaires des amis et des abonnés ; enfin, je vérifie le contenu de la page.\\[4pt]
有时候有人发垃圾信息，我就马上删除它们。\\
Yǒu shíhou yǒu rén fā lājī xìnxī, wǒ jiù mǎshàng shānchú tāmen.\\
Parfois, quelqu'un envoie des messages indésirables, alors je les supprime aussitôt.\\[4pt]
我很喜欢网络，但是我要小心，因为网络上有不好的人。\\
Wǒ hěn xǐhuan wǎngluò, dànshì wǒ yào xiǎoxīn, yīnwèi wǎngluò shang yǒu bù hǎo de rén.\\
J'aime beaucoup Internet, mais je dois faire attention, parce qu'il y a de mauvaises personnes sur Internet.\\[4pt]
因为我很认真，所以我们的粉丝越来越多。大家都说这个网页很好，我很高兴。\\
Yīnwèi wǒ hěn rènzhēn, suǒyǐ wǒmen de fěnsī yuè lái yuè duō. Dàjiā dōu shuō zhège wǎngyè hěn hǎo, wǒ hěn gāoxìng.\\
Parce que je suis sérieuse, nos abonnés sont de plus en plus nombreux. Tout le monde dit que cette page est très bien, je suis contente.
\end{textebox}

\textbf{Glossaire rapide :} \zh{垃圾信息} \emph{lājī xìnxī} « messages indésirables, spam » ; \zh{马上} \emph{mǎshàng} « aussitôt » ; \zh{检查} \emph{jiǎnchá} « vérifier » ; \zh{认真} \emph{rènzhēn} « sérieux, appliqué » ; \zh{越来越} \emph{yuè lái yuè} « de plus en plus » ; \zh{它们} \emph{tāmen} « ils/elles » (pour les choses, voir la mise en garde plus bas).

\textbf{Questions de compréhension :}
\begin{enumerate}
\item Que fait Marlyse chaque jour ? (她每天做什么？\emph{Tā měitiān zuò shénme ?})
\item Que fait-elle quand quelqu'un envoie des messages indésirables ?
\item Pourquoi le nombre d'abonnés augmente-t-il ?
\end{enumerate}

\begin{attention}[他们 ou 它们 ?]
Les deux se prononcent \emph{tāmen}, mais ne s'écrivent pas pareil : \zh{他们} (avec le radical de l'homme \zh{亻}) désigne des \textbf{personnes} ; \zh{它们} (avec le radical du toit \zh{宀}) désigne des \textbf{choses ou des animaux}. Les messages indésirables sont des choses : on écrit donc \zh{删除它们} « les supprimer », et non *\zh{删除他们}. À l'écrit, l'examen sanctionne cette confusion !
\end{attention}

\courssub{Le plan type d'un texte de gestion}

\begin{definition}[Plan type en quatre étapes]
Un texte simple sur la gestion des réseaux sociaux suit presque toujours le même plan :
\begin{enumerate}
\item \cle{Qui / quoi} : je me présente et je dis ce que je gère → \zh{我帮...管理...} \emph{Wǒ bāng... guǎnlǐ...} « j'aide... à gérer... »
\item \cle{Les actions} : je décris mes tâches dans l'ordre → \zh{首先...然后...最后...}
\item \cle{Difficulté et solution} : je raconte un problème et ma réaction → \zh{有时候...我就...}
\item \cle{Le bilan} : je conclus avec une cause et un résultat → \zh{因为...所以...}
\end{enumerate}
\end{definition}

\begin{methode}[Comment structurer son texte en quatre étapes]
\begin{enumerate}
\item Écris d'abord ton plan \textbf{en français} : (1) qui je suis et ce que je gère ; (2) mes actions ; (3) une difficulté et ma solution ; (4) mon bilan.
\item Traduis \textbf{étape par étape}, avec \textbf{un connecteur par étape} : \zh{首先}, \zh{然后}, \zh{最后} pour les actions ; \zh{就} pour la solution ; \zh{因为...所以...} pour le bilan.
\item Vérifie la place de chaque complément de temps : \textbf{avant le verbe}.
\item Relis en comptant tes connecteurs : un texte de 8 à 10 phrases doit en contenir au moins 4 ou 5.
\end{enumerate}
\end{methode}

\courssub{Les connecteurs de séquence : 首先, 然后, 最后}

Ces trois connecteurs ordonnent les actions dans le temps, comme « d'abord, ensuite, enfin » en français. Ils se placent \textbf{en tête de phrase} (ou juste après le sujet), suivis d'une virgule chinoise \zh{，}.

\begin{exemplebox}[La journée de travail d'un gestionnaire de page]
\zh{首先，我发布照片；然后，我回复评论；最后，我看看粉丝的数量。}\\
\emph{Shǒuxiān, wǒ fābù zhàopiàn ; ránhòu, wǒ huífù pínglùn ; zuìhòu, wǒ kànkan fěnsī de shùliàng.}\\
« D'abord, je publie des photos ; ensuite, je réponds aux commentaires ; enfin, je regarde le nombre d'abonnés. »
\end{exemplebox}

Remarque : dans \zh{看看} \emph{kànkan} « jeter un coup d'œil », le verbe est \cle{redoublé} pour adoucir l'action ; la deuxième syllabe perd son ton.

\begin{attention}[首先 n'est pas 第一]
\zh{首先} \emph{shǒuxiān} signifie « d'abord, en premier lieu » pour ordonner des \textbf{actions}. \zh{第一} \emph{dì-yī} signifie « premier / numéro un » (un classement). Ne dis pas *\zh{第一，我发布照片} pour « d'abord, je publie des photos » ; dis \zh{首先，我发布照片。}
\end{attention}

\courssub{Le connecteur de conséquence immédiate : 就}

\zh{就} \emph{jiù} est un petit mot très fréquent. Placé \textbf{devant le verbe} de la deuxième partie de la phrase, il marque la \cle{conséquence immédiate} : « alors, aussitôt, du coup ».

\begin{propriete}[Structure : situation + 就 + réaction]
\textbf{Schéma :} Proposition 1 (situation) ，Sujet + \zh{就} + Verbe (réaction).\\
\zh{有人发垃圾信息，我就删除它们。}\\
\emph{Yǒu rén fā lājī xìnxī, wǒ jiù shānchú tāmen.}\\
« Quand quelqu'un envoie des messages indésirables, je les supprime aussitôt. »\\[3pt]
\zh{粉丝有问题，我就回答。}\\
\emph{Fěnsī yǒu wèntí, wǒ jiù huídá.}\\
« Si les abonnés ont une question, je réponds aussitôt. »
\end{propriete}

\begin{attention}[就 se place devant le verbe, jamais en tête de phrase]
En français, « alors » commence la phrase. En chinois, \zh{就} se met \textbf{après le sujet et avant le verbe} : \zh{我就删除} « moi, aussitôt je supprime ». Ne dis jamais *\zh{就我删除它们}.
\end{attention}

\courssub{Cause et conséquence : 因为...所以...}

\begin{propriete}[因为 et 所以 peuvent coexister dans la même phrase]
\textbf{Schéma :} \zh{因为} + cause ，\zh{所以} + conséquence.\\
En chinois, contrairement au français, on peut (et on préfère souvent) utiliser \textbf{les deux connecteurs ensemble} :\\
\zh{因为我很认真，所以粉丝越来越多。}\\
\emph{Yīnwèi wǒ hěn rènzhēn, suǒyǐ fěnsī yuè lái yuè duō.}\\
« Parce que je suis sérieux, les abonnés sont de plus en plus nombreux. »\\[3pt]
On peut aussi n'utiliser que \zh{因为} (cause seule) ou que \zh{所以} (conséquence seule) :\\
\zh{我很小心，因为网络上有不好的人。} \emph{Wǒ hěn xiǎoxīn, yīnwèi wǎngluò shang yǒu bù hǎo de rén.} « Je fais attention parce qu'il y a de mauvaises personnes sur Internet. »\\
\zh{网络很有用，所以我很喜欢它。} \emph{Wǎngluò hěn yǒuyòng, suǒyǐ wǒ hěn xǐhuan tā.} « Internet est très utile, donc je l'aime beaucoup. »
\end{propriete}

\begin{attention}[Piège de traduction : « parce que... donc »]
En français, dire « \emph{parce que} je suis sérieux, \emph{donc} les abonnés augmentent » est incorrect : on choisit l'un ou l'autre. En chinois, c'est l'inverse : la combinaison \zh{因为...所以...} est \textbf{normale et naturelle}. Quand tu traduis du chinois vers le français, supprime l'un des deux ; quand tu traduis du français vers le chinois, tu peux mettre les deux.
\end{attention}

\courssub{Opposition : 但是 ; fréquence : 每天, 有时候}

\zh{但是} \emph{dànshì} « mais » marque l'opposition, comme en français. Il se place en tête de la deuxième proposition.

\begin{exemplebox}[Opposition et prudence]
\zh{我很喜欢网络，但是我要小心。}\\
\emph{Wǒ hěn xǐhuan wǎngluò, dànshì wǒ yào xiǎoxīn.}\\
« J'aime beaucoup Internet, mais je dois faire attention. »\\[3pt]
\zh{管理网页很累，但是很有意思。}\\
\emph{Guǎnlǐ wǎngyè hěn lèi, dànshì hěn yǒu yìsi.}\\
« Gérer une page web est fatigant, mais c'est très intéressant. »
\end{exemplebox}

Les expressions de fréquence \zh{每天} \emph{měitiān} « chaque jour » et \zh{有时候} \emph{yǒu shíhou} « parfois » sont des \cle{compléments de temps}. Or, en chinois, la règle d'or est la suivante :

\begin{propriete}[Le complément de temps se place AVANT le verbe]
\textbf{Schéma :} Sujet + Temps + Verbe + Objet\\
\zh{我每天发布照片。} \emph{Wǒ měitiān fābù zhàopiàn.} « Je publie des photos chaque jour. »\\
\zh{他有时候回复评论。} \emph{Tā yǒu shíhou huífù pínglùn.} « Il répond parfois aux commentaires. »
\end{propriete}

\begin{attention}[Erreur n° 1 des francophones]
En français, « je publie des photos \emph{chaque jour} » met le temps à la fin. En chinois, c'est impossible : ne dis jamais *\zh{我发布照片每天}. La seule place correcte est \zh{我每天发布照片} : le temps se met \textbf{entre le sujet et le verbe} (ou en tout début de phrase : \zh{每天我发布照片。}).
\end{attention}

\begin{popfigure}
\begin{tikzpicture}[font=\small]
\node[draw=popblue, fill=popblueL, rounded corners, minimum width=2.2cm, minimum height=0.9cm, align=center] (s) at (0,0) {\textbf{Sujet}\\ 我 wǒ};
\node[draw=poporange, fill=poporangeL, rounded corners, minimum width=2.2cm, minimum height=0.9cm, align=center] (t) at (3.4,0) {\textbf{Temps}\\ 每天 měitiān};
\node[draw=popgreen, fill=popgreenL, rounded corners, minimum width=2.2cm, minimum height=0.9cm, align=center] (v) at (6.8,0) {\textbf{Verbe}\\ 发布 fābù};
\node[draw=poppurple, fill=poppurpleL, rounded corners, minimum width=2.2cm, minimum height=0.9cm, align=center] (o) at (10.2,0) {\textbf{Objet}\\ 照片 zhàopiàn};
\draw[-{Stealth}, thick, popdark] (s) -- (t);
\draw[-{Stealth}, thick, popdark] (t) -- (v);
\draw[-{Stealth}, thick, popdark] (v) -- (o);
\draw[-{Stealth}, very thick, poppink, bend right=35] (t.south) to node[below, poppink, align=center]{le temps passe\\ avant le verbe} (v.south);
\end{tikzpicture}
\legende{Structure de la phrase chinoise : Sujet + Temps + Verbe + Objet. Le complément de temps se place toujours avant le verbe.}
\end{popfigure}

\courssub{Tableau récapitulatif des connecteurs}

\begin{center}
\begin{tabularx}{\linewidth}{|X|X|X|X|}
\hline
\rowcolor{popblueL} Caractères & Pinyin & Fonction & Traduction \\
\hline
首先 & shǒuxiān & séquence (début) & d'abord \\
\hline
然后 & ránhòu & séquence (milieu) & ensuite \\
\hline
最后 & zuìhòu & séquence (fin) & enfin \\
\hline
就 & jiù & conséquence immédiate & alors, aussitôt \\
\hline
因为...所以... & yīnwèi...suǒyǐ... & cause + conséquence & parce que... donc... \\
\hline
但是 & dànshì & opposition & mais \\
\hline
每天 & měitiān & fréquence (temps) & chaque jour \\
\hline
有时候 & yǒu shíhou & fréquence (temps) & parfois \\
\hline
\end{tabularx}
\end{center}

\begin{center}
\begin{tabularx}{\linewidth}{|X|X|X|}
\hline
\rowcolor{popblueL} Structure & Exemple en chinois + pinyin & Traduction \\
\hline
S + 时间 + V + O & 我每天发布照片。Wǒ měitiān fābù zhàopiàn. & Je publie des photos chaque jour. \\
\hline
首先...然后...最后... & 首先我发布照片，然后我回复评论。Shǒuxiān wǒ fābù zhàopiàn, ránhòu wǒ huífù pínglùn. & D'abord je publie des photos, ensuite je réponds aux commentaires. \\
\hline
Situation + 就 + réaction & 有人发垃圾信息，我就删除它们。Yǒu rén fā lājī xìnxī, wǒ jiù shānchú tāmen. & Si quelqu'un envoie du spam, je le supprime aussitôt. \\
\hline
因为 + cause, 所以 + résultat & 因为我很认真，所以粉丝越来越多。Yīnwèi wǒ hěn rènzhēn, suǒyǐ fěnsī yuè lái yuè duō. & Parce que je suis sérieux, les abonnés augmentent. \\
\hline
...但是... & 我很喜欢网络，但是我要小心。Wǒ hěn xǐhuan wǎngluò, dànshì wǒ yào xiǎoxīn. & J'aime Internet, mais je fais attention. \\
\hline
\end{tabularx}
\end{center}

\courssub{Exercice résolu et entraînement}

\begin{exoresolu}[Remettre le texte dans l'ordre et traduire]
\textbf{Énoncé.} Voici quatre phrases mélangées. Remets-les dans l'ordre logique du plan type (présentation → actions → difficulté → bilan), puis traduis-les en français.\\
a) \zh{有时候有人发不好的评论，我就删除它们。}\\
b) \zh{因为我很努力，所以我们的粉丝很多。}\\
c) \zh{我帮我们学校管理网页。}\\
d) \zh{首先，我每天发布新闻；然后，我回复评论。}
\tcblower
\textbf{Solution détaillée.}\\
Ordre logique : \textbf{c → d → a → b}.\\
(1) Présentation : \zh{我帮我们学校管理网页。} \emph{Wǒ bāng wǒmen xuéxiào guǎnlǐ wǎngyè.} « J'aide notre école à gérer sa page web. »\\
(2) Actions : \zh{首先，我每天发布新闻；然后，我回复评论。} \emph{Shǒuxiān, wǒ měitiān fābù xīnwén ; ránhòu, wǒ huífù pínglùn.} « D'abord, je publie des nouvelles chaque jour ; ensuite, je réponds aux commentaires. » (Note : \zh{每天} est bien placé avant le verbe \zh{发布}.)\\
(3) Difficulté et solution : \zh{有时候有人发不好的评论，我就删除它们。} \emph{Yǒu shíhou yǒu rén fā bù hǎo de pínglùn, wǒ jiù shānchú tāmen.} « Parfois quelqu'un publie de mauvais commentaires, alors je les supprime aussitôt. » (Note : \zh{它们} et non \zh{他们}, car ce sont des choses.)\\
(4) Bilan : \zh{因为我很努力，所以我们的粉丝很多。} \emph{Yīnwèi wǒ hěn nǔlì, suǒyǐ wǒmen de fěnsī hěn duō.} « Parce que je travaille dur, nous avons beaucoup d'abonnés. »
\end{exoresolu}

\begin{exercice}[À toi de structurer]
Traduis en chinois le texte suivant en respectant le plan type et en utilisant au moins quatre connecteurs : « J'aide mon club de football à gérer sa page. D'abord, je publie des photos chaque jour ; ensuite, je réponds aux commentaires. Parfois, quelqu'un envoie des messages indésirables, alors je les supprime aussitôt. J'aime beaucoup ce travail, mais c'est fatigant. Parce que je suis sérieux, nos abonnés sont de plus en plus nombreux. »
\end{exercice}

\begin{corrige}[Exercice « À toi de structurer »]
\zh{我帮我们足球俱乐部管理网页。首先，我每天发布照片；然后，我回复评论。有时候有人发垃圾信息，我就马上删除它们。我很喜欢这个工作，但是很累。因为我很认真，所以我们的粉丝越来越多。}\\
\emph{Wǒ bāng wǒmen zúqiú jùlèbù guǎnlǐ wǎngyè. Shǒuxiān, wǒ měitiān fābù zhàopiàn ; ránhòu, wǒ huífù pínglùn. Yǒu shíhou yǒu rén fā lājī xìnxī, wǒ jiù mǎshàng shānchú tāmen. Wǒ hěn xǐhuan zhège gōngzuò, dànshì hěn lèi. Yīnwèi wǒ hěn rènzhēn, suǒyǐ wǒmen de fěnsī yuè lái yuè duō.}\\
Points de vigilance : \zh{每天} avant \zh{发布} ; \zh{就} après le sujet \zh{我} ; \zh{它们} pour les messages (choses) ; les deux connecteurs \zh{因为...所以...} conservés ensemble.
\end{corrige}

\begin{aretenir}[L'essentiel de la section]
\begin{itemize}
\item Un texte de gestion suit 4 étapes : présentation → actions → difficulté/solution → bilan.
\item \zh{首先...然后...最后...} ordonnent les actions ; \zh{就} (devant le verbe) marque la réaction immédiate ; \zh{但是} marque l'opposition.
\item \zh{因为} et \zh{所以} peuvent coexister dans la même phrase chinoise, contrairement au français.
\item Le complément de temps (\zh{每天}, \zh{有时候}) se place \textbf{avant le verbe} : \zh{我每天发布照片。}
\item Pour les choses, le pronom pluriel s'écrit \zh{它们} ; pour les personnes, \zh{他们} (même prononciation \emph{tāmen}).
\end{itemize}
\end{aretenir}

\coursec{Thème : du français vers le chinois}

Dans la section précédente, nous avons étudié la structure d'un texte simple et ses connecteurs (\zh{因为……所以……}, \zh{每天}, \zh{有时候}). Nous savons donc maintenant comment un texte chinois est organisé. Il est temps de passer à la pratique : le \cle{thème}, c'est-à-dire la traduction du \textbf{français vers le chinois}. C'est l'exercice le plus exigeant de l'expression écrite, car il mobilise à la fois le vocabulaire, la grammaire et l'écriture des caractères. Cette section vous donne une méthode sûre en quatre étapes, appliquée à quatre phrases modèles sur la gestion des réseaux sociaux.

\courssub{Observer avant de traduire : quatre phrases modèles}

Avant d'énoncer la méthode, observons quatre phrases déjà traduites. Lisez-les attentivement et cherchez ce qui change entre le français et le chinois.

\begin{textebox}[Quatre phrases modèles — gestion des réseaux sociaux]
1. 我管理妈妈的账号。\\
\emph{Wǒ guǎnlǐ māma de zhànghào.}\\
« Je gère le compte de ma mère. »\\[4pt]
2. 我每天发布新的照片。\\
\emph{Wǒ měitiān fābù xīn de zhàopiàn.}\\
« Chaque jour, je publie de nouvelles photos. »\\[4pt]
3. 有时候，我删除垃圾信息。\\
\emph{Yǒu shíhou, wǒ shānchú lājī xìnxī.}\\
« Parfois, je supprime les spams. »\\[4pt]
4. 因为评论很多，所以我晚上回复。\\
\emph{Yīnwèi pínglùn hěn duō, suǒyǐ wǒ wǎnshang huífù.}\\
« Parce qu'il y a beaucoup de commentaires, je réponds le soir. »
\end{textebox}

Observations à noter dans votre cahier :

\begin{itemize}
\item Dans la phrase 1, « de ma mère » devient \zh{妈妈的} : la particule \cle{的} relie les deux noms, alors que le français utilise la préposition « de ».
\item Dans la phrase 2, « chaque jour » (\zh{每天}) se place \textbf{entre le sujet et le verbe}, pas au début comme en français.
\item Dans la phrase 3, « parfois » (\zh{有时候}) peut commencer la phrase, suivi d'une virgule, puis le sujet.
\item Dans la phrase 4, la cause et la conséquence s'expriment avec le couple \zh{因为……所以……}, alors que le français n'utilise qu'un seul mot (« parce que »).
\end{itemize}

\courssub{La méthode du thème en quatre étapes}

\begin{methode}[Réussir un thème : Lire — Analyser — Traduire — Vérifier]
\begin{enumerate}
\item \textbf{Lire} la phrase française entièrement, sans commencer à écrire. Repérer le sens global : qui fait quoi, quand, pourquoi ?
\item \textbf{Analyser} la phrase : identifier le \cle{sujet} (S), le \cle{verbe} (V) et le \cle{complément d'objet} (O). Repérer aussi les compléments de temps et de cause.
\item \textbf{Traduire} mot à mot d'abord sur un brouillon, puis \textbf{réorganiser} selon l'ordre chinois : Sujet + temps + Verbe + Objet. Choisir le bon caractère pour chaque mot (attention aux caractères proches).
\item \textbf{Vérifier} : relire en contrôlant les traits des caractères, la place des compléments de temps, la présence de \zh{的} entre deux noms, et la ponctuation chinoise pleine chasse (\zh{，} \zh{。}).
\end{enumerate}
\end{methode}

\begin{popfigure}
\begin{tikzpicture}[
  node distance=0.55cm,
  etape/.style={rectangle, rounded corners=4pt, draw=popblue, fill=popblueL, thick, minimum width=4.6cm, minimum height=0.95cm, align=center, font=\small},
  detail/.style={rectangle, draw=popmuted, fill=popcream, rounded corners=3pt, minimum width=5.4cm, minimum height=0.95cm, align=center, font=\small},
  fleche/.style={-{Stealth[length=3mm]}, thick, popdark}
]
\node[etape, align=center] (lire){\textbf{1. Lire}\\ Comprendre le sens global};
\node[detail, right=1.1cm of lire] (d1) {Qui ? Fait quoi ? Quand ?};
\node[etape, below=of lire, align=center] (analyser){\textbf{2. Analyser}\\ Sujet / Verbe / Objet};
\node[detail, right=1.1cm of analyser] (d2) {Repérer temps, cause, possession};
\node[etape, below=of analyser, align=center] (traduire){\textbf{3. Traduire}\\ Mot à mot puis réorganiser};
\node[detail, right=1.1cm of traduire] (d3) {Ordre chinois : S + temps + V + O};
\node[etape, below=of traduire, align=center] (verifier){\textbf{4. Vérifier}\\ Relire et corriger};
\node[detail, right=1.1cm of verifier] (d4) {Traits, \zh{的}, place du temps, tons};
\draw[fleche] (lire) -- (analyser);
\draw[fleche] (analyser) -- (traduire);
\draw[fleche] (traduire) -- (verifier);
\draw[fleche, dashed, poporange] (verifier.west) to[bend left=25] node[left, font=\small, text=poporange]{si erreur} (traduire.west);
\end{tikzpicture}
\legende{Organigramme de la méthode du thème en quatre étapes}
\end{popfigure}

\begin{aretenir}[L'ordre des mots en chinois]
La phrase chinoise de base suit l'ordre : \textbf{Sujet + (temps / lieu) + Verbe + Objet}.\\
Contrairement au français, le complément de temps ne se met presque jamais à la fin : il se place \textbf{après le sujet et avant le verbe}.\\
Français : « Je publie des photos \emph{chaque jour}. »\\
Chinois : \zh{我每天发布照片。} \emph{Wǒ měitiān fābù zhàopiàn.} (litt. « je chaque-jour publie photos »)
\end{aretenir}

\courssub{Application guidée : les quatre phrases pas à pas}

\textbf{Phrase 1 : « Je gère le compte de ma mère. »}

\begin{itemize}
\item Étape 2 (analyse) : S = je ; V = gère ; O = le compte de ma mère.
\item Étape 3 (traduction) : je → \zh{我} ; gérer → \zh{管理} ; compte → \zh{账号} ; de ma mère → \zh{妈妈的}.
\item Résultat : \zh{我管理妈妈的账号。} \emph{Wǒ guǎnlǐ māma de zhànghào.}
\end{itemize}

\begin{attention}[La particule \zh{的} entre deux noms]
« De ma mère » se traduit obligatoirement par \zh{妈妈的} : la particule \cle{的} marque la possession ou le lien entre deux noms. \textbf{Ne l'oubliez jamais.}\\
\zh{妈妈账号} est une erreur ; il faut \zh{妈妈的账号} (le compte \emph{de} maman).\\
Autres exemples : \zh{我的手机} \emph{wǒ de shǒujī} « mon téléphone » ; \zh{老师的照片} \emph{lǎoshī de zhàopiàn} « les photos du professeur ».
\end{attention}

\textbf{Phrase 2 : « Chaque jour, je publie de nouvelles photos. »}

\begin{itemize}
\item Analyse : S = je ; V = publie ; O = de nouvelles photos ; temps = chaque jour.
\item Traduction mot à mot : \zh{每天} (chaque jour) + \zh{我} + \zh{发布} (publier) + \zh{新的照片} (de nouvelles photos).
\item Réorganisation : le temps passe après le sujet → \zh{我每天发布新的照片。} \emph{Wǒ měitiān fābù xīn de zhàopiàn.}
\end{itemize}

\begin{attention}[La place de \zh{每天}]
Le francophone place naturellement « chaque jour » au début ou à la fin. En chinois, \zh{每天} se met \textbf{entre le sujet et le verbe} : \zh{我每天发布}, jamais \zh{我发布照片每天}. Notez aussi le \zh{的} de \zh{新的照片} : entre un adjectif de plus d'une syllabe et le nom, on garde \zh{的}.
\end{attention}

\textbf{Phrase 3 : « Parfois, je supprime les spams. »}

\begin{itemize}
\item Analyse : S = je ; V = supprime ; O = les spams ; temps = parfois.
\item Traduction : parfois → \zh{有时候} ; supprimer → \zh{删除} ; spam → \zh{垃圾信息} (litt. « information-poubelle »).
\item Résultat : \zh{有时候，我删除垃圾信息。} \emph{Yǒu shíhou, wǒ shānchú lājī xìnxī.} Ici, \zh{有时候} peut ouvrir la phrase, suivi d'une virgule chinoise \zh{，} ; on peut aussi dire \zh{我有时候删除垃圾信息。}
\end{itemize}

\textbf{Phrase 4 : « Parce qu'il y a beaucoup de commentaires, je réponds le soir. »}

\begin{itemize}
\item Analyse : cause = il y a beaucoup de commentaires ; S = je ; V = réponds ; temps = le soir.
\item Traduction : parce que → \zh{因为} ; commentaire → \zh{评论} ; beaucoup → \zh{很多} ; donc → \zh{所以} ; le soir → \zh{晚上} ; répondre → \zh{回复}.
\item Résultat : \zh{因为评论很多，所以我晚上回复。} \emph{Yīnwèi pínglùn hěn duō, suǒyǐ wǒ wǎnshang huífù.}
\end{itemize}

\begin{propriete}[Le couple \zh{因为……所以……}]
En chinois, on emploie souvent \textbf{les deux} corrélatifs : \zh{因为} (cause) en tête de la première proposition et \zh{所以} (conséquence) en tête de la seconde, séparées par une virgule. Le français, lui, n'utilise que « parce que ». Structure : \zh{因为} + cause \zh{，} \zh{所以} + Sujet + temps + Verbe. Le temps (\zh{晚上}) reste avant le verbe (\zh{回复}).
\end{propriete}

\begin{exemplebox}[Un exemple supplémentaire traduit]
« Il bloque les spams. » → \zh{他屏蔽垃圾信息。} \emph{Tā píngbì lājī xìnxī.}\\
Analyse : S = \zh{他} (il) ; V = \zh{屏蔽} (bloquer, filtrer) ; O = \zh{垃圾信息} (les spams). L'ordre S-V-O est identique au français : phrase facile pour commencer un thème.
\end{exemplebox}

\courssub{Tableau récapitulatif : correspondances français / chinois}

\begin{center}
\begin{tabularx}{\linewidth}{|X|X|X|}
\hline
\rowcolor{popblueL} \textbf{Français} & \textbf{Chinois (caractères + pinyin)} & \textbf{Point de vigilance} \\
\hline
je gère & 我管理 \emph{wǒ guǎnlǐ} & S + V, ordre identique \\
\hline
le compte de ma mère & 妈妈的账号 \emph{māma de zhànghào} & \zh{的} obligatoire entre les deux noms \\
\hline
chaque jour & 每天 \emph{měitiān} & placé après le sujet, avant le verbe \\
\hline
de nouvelles photos & 新的照片 \emph{xīn de zhàopiàn} & \zh{的} après l'adjectif \zh{新} \\
\hline
parfois & 有时候 \emph{yǒu shíhou} & en tête de phrase + virgule, ou après le sujet \\
\hline
les spams & 垃圾信息 \emph{lājī xìnxī} & mot composé : « déchet + information » \\
\hline
parce que … donc … & 因为……所以…… \emph{yīnwèi … suǒyǐ …} & les deux corrélatifs en chinois \\
\hline
le soir & 晚上 \emph{wǎnshang} & avant le verbe : \zh{我晚上回复} \\
\hline
\end{tabularx}
\end{center}

\courssub{Le piège des caractères proches : \zh{已} et \zh{己}}

\begin{attention}[\zh{已} / \zh{己} : un trait qui change tout]
Ces deux caractères se ressemblent beaucoup : \zh{已} \emph{yǐ} (« déjà », dans \zh{已经}) et \zh{己} \emph{jǐ} (« soi-même », dans \zh{自己}).\\
Moyen mnémotechnique : pour \zh{已} (déjà), la barre verticale \textbf{dépasse} à moitié — « déjà à moitié sorti » ; pour \zh{己} (soi-même), elle \textbf{ne dépasse pas} — « on reste chez soi ».\\
Dans nos phrases de thème, méfiez-vous aussi de \zh{账} (compte, clé \zh{贝} « argent ») qu'il ne faut pas confondre avec \zh{帐} ; et de \zh{照} (photo) dont la clé est \zh{灬} (le feu, en bas).
\end{attention}

\courssub{Exercice résolu et entraînement}

\begin{exoresolu}[Thème guidé]
Énoncé : traduisez en chinois « Ma mère publie des photos chaque jour. »\\
Indice : publier = \zh{发布} ; photo = \zh{照片} ; chaque jour = \zh{每天}.
\tcblower
Solution détaillée :
\begin{enumerate}
\item Lecture : qui ? ma mère ; quoi ? publie des photos ; quand ? chaque jour.
\item Analyse : S = ma mère ; V = publie ; O = des photos ; temps = chaque jour.
\item Traduction mot à mot : ma mère → \zh{我的妈妈} ou simplement \zh{妈妈} ; publier → \zh{发布} ; photos → \zh{照片} ; chaque jour → \zh{每天}.
\item Réorganisation selon l'ordre chinois S + temps + V + O : \zh{妈妈每天发布照片。} \emph{Māma měitiān fābù zhàopiàn.}
\item Vérification : \zh{每天} est bien avant le verbe ; pas de \zh{的} nécessaire entre \zh{妈妈} et le verbe ; ponctuation chinoise \zh{。} en fin de phrase.
\end{enumerate}
\end{exoresolu}

\begin{exercice}[À vous de jouer — collectif puis individuel]
Traduisez en chinois, en suivant les quatre étapes de la méthode. La première phrase se fait collectivement au tableau avec le professeur ; les trois suivantes individuellement sur le cahier.
\begin{enumerate}
\item « Je gère le compte de mon père. » (père = \zh{爸爸})
\item « Parfois, ma mère supprime les spams. »
\item « Il publie de nouvelles photos le soir. »
\item « Parce qu'il y a beaucoup de spams, je bloque les comptes. » (bloquer = \zh{屏蔽})
\end{enumerate}
\end{exercice}

\begin{corrige}[Exercice « À vous de jouer »]
\begin{enumerate}
\item \zh{我管理爸爸的账号。} \emph{Wǒ guǎnlǐ bàba de zhànghào.} — \zh{的} obligatoire entre \zh{爸爸} et \zh{账号}.
\item \zh{有时候，我妈妈删除垃圾信息。} \emph{Yǒu shíhou, wǒ māma shānchú lājī xìnxī.} — « ma mère » = \zh{我妈妈} (le \zh{的} est facultatif après un pronom possessif devant un lien de famille).
\item \zh{他晚上发布新的照片。} \emph{Tā wǎnshang fābù xīn de zhàopiàn.} — \zh{晚上} avant le verbe, \zh{的} après \zh{新}.
\item \zh{因为垃圾信息很多，所以我屏蔽账号。} \emph{Yīnwèi lājī xìnxī hěn duō, suǒyǐ wǒ píngbì zhànghào.} — couple \zh{因为……所以……} complet.
\end{enumerate}
\end{corrige}

\begin{aretenir}[Les cinq réflexes du thème réussi]
\begin{enumerate}
\item Je lis toute la phrase avant d'écrire.
\item Je repère S, V, O et les compléments de temps.
\item Je place le temps \textbf{avant} le verbe (\zh{我每天发布}, \zh{我晚上回复}).
\item Je n'oublie jamais \zh{的} entre deux noms (\zh{妈妈的账号}).
\item Je relève les caractères proches (\zh{已}/\zh{己}) et je relis ma phrase à voix haute.
\end{enumerate}
\end{aretenir}

\coursec{Version : du chinois vers le français}

Dans la section précédente, nous avons appris à passer du français au chinois (le \cle{thème}) : il fallait construire des phrases chinoises correctes, choisir les bons caractères et respecter l'ordre des mots. Nous allons maintenant faire le chemin inverse : partir d'un texte chinois et le rendre en français clair et fidèle. C'est la \cle{version}. La version est un exercice royal au baccalauréat : elle vérifie à la fois votre compréhension de l'écrit chinois et votre maîtrise du français. Restons sur notre thème du module : la gestion des réseaux sociaux et d'une boutique en ligne.

\courssub{Le texte de version : lecture et observation}

\begin{textebox}[Un grand frère qui gère sa boutique en ligne]
我的哥哥管理一个网店。首先，他发布商品的照片；然后，他回复顾客的评论。有时候有人发垃圾信息，他就屏蔽他们。因为他的服务很好，所以粉丝越来越多。\\
\emph{Wǒ de gēge guǎnlǐ yí ge wǎngdiàn. Shǒuxiān, tā fābù shāngpǐn de zhàopiàn ; ránhòu, tā huífù gùkè de pínglùn. Yǒu shíhou yǒu rén fā lājī xìnxī, tā jiù píngbì tāmen. Yīnwèi tā de fúwù hěn hǎo, suǒyǐ fěnsī yuè lái yuè duō.}\\
« Mon grand frère gère une boutique en ligne. D'abord, il publie les photos des produits ; ensuite, il répond aux commentaires des clients. Parfois, des gens envoient des spams, alors il les bloque. Comme son service est bon, ses abonnés sont de plus en plus nombreux. »
\end{textebox}

L'enseignant lit le texte à voix haute, phrase par phrase, en insistant sur les tons : \zh{管理} \emph{guǎnlǐ} (3\textsuperscript{e} + 3\textsuperscript{e} ton, le premier devient un 2\textsuperscript{e} ton à l'oral), \zh{商品} \emph{shāngpǐn}, \zh{顾客} \emph{gùkè}, \zh{屏蔽} \emph{píngbì}, \zh{粉丝} \emph{fěnsī}, \zh{因为} \emph{yīnwèi}, \zh{所以} \emph{suǒyǐ} (3\textsuperscript{e} + 3\textsuperscript{e} ton, le premier devient un 2\textsuperscript{e} ton à l'oral). Les élèves répètent en chœur, puis lisent individuellement.

\textbf{Glossaire du texte}

\begin{center}
\begin{tabularx}{\linewidth}{|X|X|X|X|}
\hline
\rowcolor{popblueL} Caractères & Pinyin & Nature & Traduction \\
\hline
哥哥 & gēge & nom & grand frère \\
\hline
管理 & guǎnlǐ & verbe & gérer, administrer \\
\hline
网店 & wǎngdiàn & nom & boutique en ligne \\
\hline
首先 & shǒuxiān & adverbe & d'abord, premièrement \\
\hline
发布 & fābù & verbe & publier, poster \\
\hline
商品 & shāngpǐn & nom & produits, marchandises \\
\hline
照片 & zhàopiàn & nom & photo \\
\hline
然后 & ránhòu & adverbe & ensuite, puis \\
\hline
回复 & huífù & verbe & répondre (à) \\
\hline
顾客 & gùkè & nom & client \\
\hline
评论 & pínglùn & nom & commentaire, avis \\
\hline
有时候 & yǒu shíhou & locution & parfois \\
\hline
垃圾信息 & lājī xìnxī & nom & spam, message indésirable \\
\hline
屏蔽 & píngbì & verbe & bloquer, masquer \\
\hline
因为……所以…… & yīnwèi... suǒyǐ... & conjonctions & parce que... donc... \\
\hline
服务 & fúwù & nom & service \\
\hline
粉丝 & fěnsī & nom & abonnés, fans \\
\hline
越来越多 & yuè lái yuè duō & locution & de plus en plus nombreux \\
\hline
\end{tabularx}
\end{center}

\begin{propriete}[Analyse des caractères complexes du thème]
Le thème « gestion des réseaux sociaux / boutique en ligne » mobilise des caractères à structure composée (souvent clé phonétique + clé sémantique). Les identifier aide à la mémorisation et évite les confusions visuelles.

\begin{center}
\begin{tabularx}{\linewidth}{|X|X|X|X|X|}
\hline
\rowcolor{popblueL} Caractère & Pinyin & Clé (Radical) & Structure / Ordre des traits principaux & Sens de la clé \\
\hline
管 & guǎn & \zh{竹} (zhú, bambou) & Haut \zh{竹} (6 traits) → Bas \zh{官} (8 traits) & Bambou : tubes, tuyaux → gérer, contrôler \\
\hline
发 & fā & \zh{又} (yòu, main droite) & Simplifié de \zh{髮} : \zh{友} (ami) → \zh{又} & Main qui agit → envoyer, émettre \\
\hline
布 & bù & \zh{巾} (jīn, tissu) & \zh{巾} (3 traits) + \zh{父} (4 traits) & Tissu → étaler, annoncer, publier \\
\hline
屏 & píng & \zh{尸} (shī, corps/maison) & \zh{尸} (3 traits) + \zh{并} (6 traits) & Abri/écran → protéger, bloquer \\
\hline
蔽 & bì & \zh{艹} (cǎo, herbe) & \zh{艹} (3 traits) + \zh{敝} (bì, phonétique) & Herbe dense qui cache → masquer, cacher \\
\hline
粉 & fěn & \zh{米} (mǐ, riz) & \zh{米} (6 traits) + \zh{分} (fēn, phonétique) & Riz pilé → poudre,粉丝 (fans/vermicelles) \\
\hline
丝 & sī & \zh{纟} (sī, soie) & Clé \zh{纟} (3 traits) + \zh{㐌} & Fil de soie → fil, détail, 粉丝 (fans) \\
\hline
评 & píng & \zh{讠} (yán, parole) & \zh{讠} (2 traits) + \zh{平} (píng, phonétique) & Parole 平 → commenter, évaluer \\
\hline
论 & lùn & \zh{讠} (yán, parole) & \zh{讠} (2 traits) + \zh{仑} (lún, phonétique) & Parole 仑 → discuter, théorie \\
\hline
顾 & gù & \zh{页} (yè, tête/page) & \zh{页} (6 traits) + \zh{雇} (gù, phonétique) & Se tourner vers → regarder, client (顾客) \\
\hline
客 & kè & \zh{宀} (mián, toit) & \zh{宀} (3 traits) + \zh{各} (gè, phonétique) & Sous le toit → invité, client \\
\hline
服 & fú & \zh{月} (yuè, chair/lune) & \zh{月} (4 traits) + \zh{又} (yòu) + \zh{寸} (cùn) & Vêtement sur le corps → vêtement, servir \\
\hline
务 & wù & \zh{力} (lì, force) & \zh{力} (2 traits) + \zh{务} (phonétique simplifié) & Force appliquée → affaire, service, devoir \\
\hline
\end{tabularx}
\end{center}
\emph{Conseil d'écriture :} Respectez l'ordre « haut → bas, gauche → droite, horizontal avant vertical, extérieur avant intérieur ». Pour \zh{蔽} et \zh{论}, la clé sémantique (\zh{艹}, \zh{讠}) se écrit en premier (haut ou gauche).
\end{propriete}

\begin{definition}[Version]
La \cle{version} est la traduction d'un texte de la langue étrangère (ici le chinois) vers la langue maternelle ou d'enseignement (ici le français). Une bonne version est à la fois \textbf{fidèle} (elle ne trahit pas le sens du texte chinois, elle n'ajoute rien et n'oublie rien) et \textbf{correcte} (elle est écrite dans un français naturel, grammatical et élégant).
\end{definition}

\courssub{Méthode de la version en cinq étapes}

\begin{methode}[Comment réussir une version]
\begin{enumerate}
\item \textbf{Repérer les mots connus.} Lis tout le texte une première fois sans traduire. Souligne les mots que tu connais déjà : ici \zh{哥哥}, \zh{网店}, \zh{照片}, \zh{评论}, \zh{粉丝}. Ils te donnent le thème général (un frère, une boutique en ligne).
\item \textbf{Découper en segments Sujet -- Verbe -- Objet.} Le chinois suit presque toujours l'ordre S + V + O, comme le français. Découpe chaque phrase en petits blocs : \zh{他 / 发布 / 商品的照片} = il / publie / les photos des produits.
\item \textbf{Traduire littéralement au brouillon.} Écris une traduction mot à mot, même si elle sonne mal en français : « parce que son service très bon, donc fans de plus en plus nombreux ».
\item \textbf{Rédiger en bon français.} Reformule chaque segment dans un français naturel : « Comme son service est très bon, ses abonnés sont de plus en plus nombreux. » C'est ici que tu gagnes des points de style.
\item \textbf{Relire et vérifier.} Compte les phrases : chaque phrase chinoise doit correspondre à exactement une phrase française. Vérifie qu'aucun mot chinois n'a été oublié et qu'aucune idée n'a été ajoutée.
\end{enumerate}
\end{methode}

\begin{popfigure}
\begin{tikzpicture}[node distance=0.35cm, every node/.style={align=center}]
\node[draw=popblue, fill=popblueL, rounded corners, minimum width=2.6cm, minimum height=1cm, align=center] (a){1. Mots\\connus};
\node[draw=popgreen, fill=popgreenL, rounded corners, minimum width=2.6cm, minimum height=1cm, right=of a, align=center] (b){2. Segments\\S--V--O};
\node[draw=poporange, fill=poporangeL, rounded corners, minimum width=2.6cm, minimum height=1cm, right=of b, align=center] (c){3. Traduction\\littérale};
\node[draw=poppurple, fill=poppurpleL, rounded corners, minimum width=2.6cm, minimum height=1cm, right=of c, align=center] (d){4. Français\\correct};
\node[draw=poppink, fill=poppinkL, rounded corners, minimum width=2.6cm, minimum height=1cm, right=of d] (e) {5. Relecture};
\draw[-{Stealth}, thick, popdark] (a) -- (b);
\draw[-{Stealth}, thick, popdark] (b) -- (c);
\draw[-{Stealth}, thick, popdark] (c) -- (d);
\draw[-{Stealth}, thick, popdark] (d) -- (e);
\end{tikzpicture}
\legende{Les cinq étapes de la version : du texte chinois au français naturel.}
\end{popfigure}

\courssub{Application guidée : segmentons le texte ensemble}

Appliquons la méthode phrase par phrase.

\textbf{Phrase 1 :} \zh{我的哥哥管理一个网店。}\\
Segments : \zh{我的哥哥} (mon grand frère, sujet) + \zh{管理} (gère, verbe) + \zh{一个网店} (une boutique en ligne, objet).\\
Traduction littérale : « mon grand frère gère une boutique en ligne ». Français final : « Mon grand frère gère une boutique en ligne. » Aucune difficulté : l'ordre chinois et l'ordre français coïncident.

\textbf{Phrase 2 :} \zh{首先，他发布商品的照片；然后，他回复顾客的评论。}\\
Segments : \zh{首先} (d'abord) + \zh{他} (il) + \zh{发布} (publie) + \zh{商品的照片} (les photos des produits) ; \zh{然后} (ensuite) + \zh{他} (il) + \zh{回复} (répond à) + \zh{顾客的评论} (les commentaires des clients).\\
Remarque le mot-outil \zh{的} : \zh{商品的照片} se construit comme « les photos \textbf{des} produits » et \zh{顾客的评论} comme « les commentaires \textbf{des} clients ». Le \zh{的} chinois devient très souvent « de » en français.\\
Français final : « D'abord, il publie les photos des produits ; ensuite, il répond aux commentaires des clients. »

\textbf{Phrase 3 :} \zh{有时候有人发垃圾信息，他就屏蔽他们。}\\
Segments : \zh{有时候} (parfois) + \zh{有人} (il y a des gens / des gens) + \zh{发垃圾信息} (envoient des spams) ; \zh{他} (il) + \zh{就} (alors, aussitôt) + \zh{屏蔽他们} (les bloque).\\
Le petit mot \zh{就} marque ici la conséquence logique : « alors ». Ne l'oublie pas dans la traduction.\\
Français final : « Parfois, des gens envoient des spams, alors il les bloque. »

\textbf{Phrase 4 :} \zh{因为他的服务很好，所以粉丝越来越多。}\\
Segments : \zh{因为} (parce que) + \zh{他的服务很好} (son service est très bon) + \zh{所以} (donc) + \zh{粉丝} (les abonnés) + \zh{越来越多} (de plus en plus nombreux).\\
Français final : « Comme son service est bon, ses abonnés sont de plus en plus nombreux. »

\begin{propriete}[La structure 因为……所以……]
Structure : \textbf{因为 + cause，所以 + conséquence}. En chinois, on peut utiliser les deux conjonctions ensemble dans la même phrase, ce qui est impossible en français correct. En français, on choisit \emph{une seule} formulation : « parce que... » ou « donc », ou bien « comme..., ... ».\\
\zh{因为他很忙，所以他不回复。} \emph{Yīnwèi tā hěn máng, suǒyǐ tā bù huífù.} « Comme il est très occupé, il ne répond pas. » (et non « Parce qu'il est occupé, donc il ne répond pas », qui est une faute de français).
\end{propriete}

\begin{exemplebox}[Exemples traduits autour du mot 服务]
\zh{他的服务很好。} \emph{Tā de fúwù hěn hǎo.} « Son service est très bon. »\\
\zh{这个网店的服务怎么样？} \emph{Zhège wǎngdiàn de fúwù zěnmeyàng ?} « Comment est le service de cette boutique en ligne ? »\\
\zh{好的服务能带来更多顾客。} \emph{Hǎo de fúwù néng dàilái gèng duō gùkè.} « Un bon service peut attirer plus de clients. »
\end{exemplebox}

\courssub{Les points délicats du texte}

\begin{attention}[越来越多 : la nuance de progression]
\zh{越来越 + adjectif} signifie « de plus en plus... » : il exprime une \textbf{progression}, un changement qui continue. Ne traduis jamais \zh{粉丝越来越多} par « les abonnés sont nombreux » : tu perdrais l'idée d'augmentation. Dis : « les abonnés sont de plus en plus nombreux » ou « le nombre d'abonnés augmente ».\\
Autres exemples : \zh{天气越来越热。} \emph{Tiānqì yuè lái yuè rè.} « Il fait de plus en plus chaud. » ; \zh{他的中文越来越好。} \emph{Tā de Zhōngwén yuè lái yuè hǎo.} « Son chinois devient de mieux en mieux. »
\end{attention}

\begin{attention}[商品 et 服务 : deux mots à ne pas confondre]
\zh{商品} \emph{shāngpǐn} désigne les \textbf{produits, marchandises} (ce qu'on vend) ; \zh{服务} \emph{fúwù} désigne le \textbf{service} (la manière dont on traite les clients). Dans le texte, le frère publie les photos des \emph{produits}, mais c'est la qualité de son \emph{service} qui attire les abonnés. Inverser les deux mots est une faute de fidélité. Attention aussi à ne pas confondre \zh{商品} avec \zh{商店} \emph{shāngdiàn} (magasin) : les caractères \zh{品} et \zh{店} se ressemblent par leur silhouette mais n'ont ni le même sens ni la même prononciation.
\end{attention}

\begin{propriete}[Connecteurs logiques et structure du texte]
Le texte modèle utilise une chaîne de connecteurs pour organiser la narration et l'argumentation. Maîtriser ces mots-outils est indispensable pour l'expression écrite (thème) et la compréhension (version).

\begin{center}
\begin{tabularx}{\linewidth}{|X|X|X|}
\hline
\rowcolor{popblueL} Connecteur chinois & Pinyin & Emploi / Traduction contextuelle \\
\hline
首先 & shǒuxiān & Marque la première étape : « D'abord », « Premièrement ». \\
\hline
然后 & ránhòu & Marque l'étape suivante : « Ensuite », « Puis ». \\
\hline
有时候 & yǒu shíhou & Introduit une situation occasionnelle : « Parfois », « Il arrive que ». \\
\hline
就 & jiù & Marque la conséquence immédiate ou la suite logique : « alors », « du coup ». \\
\hline
因为……所以…… & yīnwèi... suǒyǐ... & Relation cause-effet : « Comme..., ... » / « Puisque..., ... ». \\
\hline
越来越 + adj. & yuè lái yuè... & Progression continue : « De plus en plus... ». \\
\hline
\end{tabularx}
\end{center}
\end{propriete}

\begin{exoresolu}[Version guidée d'une phrase nouvelle]
Énoncé : traduis en français la phrase \zh{因为我的妹妹每天回复评论，所以她的网店越来越有名。}
\tcblower
Solution détaillée. Étape 1, mots connus : \zh{妹妹} (petite sœur), \zh{每天} (chaque jour), \zh{回复} (répondre), \zh{评论} (commentaires), \zh{网店} (boutique en ligne), \zh{有名} (célèbre). Étape 2, segments : \zh{因为} + \zh{我的妹妹} + \zh{每天} + \zh{回复评论} ; \zh{所以} + \zh{她的网店} + \zh{越来越有名}. Étape 3, traduction littérale : « parce que ma petite sœur chaque jour répond aux commentaires, donc sa boutique en ligne de plus en plus célèbre ». Étape 4, français naturel : « Comme ma petite sœur répond chaque jour aux commentaires, sa boutique en ligne devient de plus en plus célèbre. » Étape 5, vérification : une seule phrase chinoise, une seule phrase française ; \zh{每天} est bien placé avant le verbe en français (« chaque jour ») ; la nuance de progression de \zh{越来越} est rendue par « de plus en plus ».
\end{exoresolu}

\courssub{La vérification croisée : le contrôle final}

Avant de rendre ta copie, fais toujours la \cle{vérification croisée} : mets en regard chaque phrase chinoise et sa phrase française. Notre texte compte quatre phrases chinoises (repérées grâce aux points 。), ta version doit donc compter exactement quatre phrases françaises. Vérifie ensuite que chaque mot important du chinois a un correspondant : \zh{首先} → « d'abord », \zh{然后} → « ensuite », \zh{有时候} → « parfois », \zh{就} → « alors », \zh{因为……所以……} → « comme... », \zh{越来越} → « de plus en plus ». Si un connecteur a disparu, ta version perd en fidélité.

\begin{exercice}[À toi de jouer]
Traduis en français le petit texte suivant, en appliquant les cinq étapes de la méthode :\\
\zh{我的姐姐管理两个网店。首先，她发布新商品；然后，她回复顾客的评论。有时候有人发垃圾信息，她就屏蔽他们。因为她的服务很好，所以她的粉丝越来越多。}\\
\emph{(Wǒ de jiějie guǎnlǐ liǎng ge wǎngdiàn. Shǒuxiān, tā fābù xīn shāngpǐn ; ránhòu, tā huífù gùkè de pínglùn. Yǒu shíhou yǒu rén fā lājī xìnxī, tā jiù píngbì tāmen. Yīnwèi tā de fúwù hěn hǎo, suǒyǐ tā de fěnsī yuè lái yuè duō.)}
\end{exercice}

\begin{corrige}[Exercice : version du texte sur la grande sœur]
« Ma grande sœur gère deux boutiques en ligne. D'abord, elle publie les nouveaux produits ; ensuite, elle répond aux commentaires des clients. Parfois, des gens envoient des spams, alors elle les bloque. Comme son service est très bon, ses abonnés sont de plus en plus nombreux. »\\
Points de vigilance : \zh{两个} exige le classificateur \zh{个} et se traduit « deux » ; \zh{新商品} = « les nouveaux produits » (pluriel implicite) ; \zh{她的粉丝} = « ses abonnés » ; ne pas oublier « de plus en plus » pour \zh{越来越多}. La version compte quatre phrases, comme le texte chinois.
\end{corrige}

\begin{savaistu}[La version au baccalauréat]
Au baccalauréat camerounais, la version est notée sur \textbf{deux critères} qui comptent chacun : la \textbf{fidélité} (as-tu bien compris et restitué le sens du texte chinois, sans omission ni ajout ?) et la \textbf{qualité du français} (orthographe, grammaire, style naturel). Une traduction exacte mais écrite en mauvais français perd des points, et un beau français qui trahit le texte chinois aussi ! D'où l'importance de l'étape 4 (reformulation) et de l'étape 5 (relecture) : les deux comptent.
\end{savaistu}

\begin{aretenir}[L'essentiel de la version]
\begin{itemize}
\item La version se fait en cinq étapes : mots connus, segments S--V--O, traduction littérale au brouillon, reformulation en français naturel, relecture.
\item Le mot-outil \zh{的} devient souvent « de » : \zh{商品的照片} = « les photos des produits ».
\item \zh{因为……所以……} ne se traduit jamais par « parce que... donc... » : en français, une seule marque de relation suffit.
\item \zh{越来越} = « de plus en plus » : toujours rendre la nuance de progression.
\item Vérification croisée obligatoire : une phrase chinoise = exactement une phrase française.
\item Les caractères complexes du thème (\zh{管, 发, 布, 屏, 蔽, 粉, 丝, 评, 论, 顾, 客, 服, 务}) se décomposent en clé sémantique + phonétique : les identifier sécurise l'écriture et la lecture.
\end{itemize}
\end{aretenir}

\coursec{Grille d'évaluation et production finale}

Après avoir traduit du chinois vers le français (la version), vous savez désormais lire, comprendre et restituer un texte sur la gestion des réseaux sociaux. Il reste l'étape ultime : \cle{produire votre propre texte} et apprendre à l'évaluer comme le ferait un correcteur d'examen. Cette dernière section vous donne la grille officielle de notation, une méthode d'auto-correction en trois relectures, et une séance de correction collective avant la production finale.

\courssub{Observer d'abord : une phrase modèle parfaite}

Avant de parler de barème, observons ce qu'un correcteur attend. Voici une phrase qui obtiendrait tous les points, car elle réunit caractères exacts, ordre des mots correct, connecteur logique et lexique du thème.

\begin{exemplebox}[La phrase modèle à imiter]
\zh{因为我很认真，所以我们的粉丝越来越多。}\\[2pt]
\emph{Yīnwèi wǒ hěn rènzhēn, suǒyǐ wǒmen de fěnsī yuè lái yuè duō.}\\[2pt]
« Parce que je suis très sérieux, nos abonnés sont de plus en plus nombreux. »
\end{exemplebox}

Pourquoi cette phrase est-elle « parfaite » ? Analysons-la point par point :

\begin{itemize}
\item \textbf{Caractères exacts} : \zh{认真} (rènzhēn, « sérieux ») est écrit avec tous ses traits ; \zh{粉丝} (fěnsī, « abonnés ») ne confond pas \zh{丝} (sī, fil de soie) avec \zh{系} (xì, système/attacher).
\item \textbf{Connecteur complet} : la paire \zh{因为……所以……} (yīnwèi... suǒyǐ..., « parce que... donc... ») est entière, ce qui est naturel en chinois alors qu'en français on évite de répéter « parce que... donc ».
\item \textbf{Ordre des mots} : Sujet (\zh{我们的粉丝}) + adverbe de degré \zh{越来越} + adjectif \zh{多}, structure standard « de plus en plus... ».
\item \textbf{Lexique du thème} : \zh{粉丝} appartient au vocabulaire des réseaux sociaux étudié dans cette leçon.
\end{itemize}

\begin{aretenir}[Ce que le correcteur récompense]
Une phrase courte, correcte et bien connectée vaut toujours mieux qu'une phrase longue et bancale. Visez la \cle{précision}, pas la performance.
\end{aretenir}

\courssub{Le barème type : 20 points en 5 critères}

\begin{propriete}[Barème type de la production écrite]
La production finale est notée sur \textbf{20 points}, répartis en \textbf{5 critères de 4 points chacun} :
\begin{enumerate}
\item \textbf{Exactitude des caractères} (4 pts) : traits complets, radicaux corrects, aucun caractère fantaisiste inventé.
\item \textbf{Ordre des mots et grammaire} (4 pts) : schéma Sujet + (temps/lieu) + Verbe + Objet, emploi correct de \zh{的}, connecteurs bien placés.
\item \textbf{Respect du plan} (4 pts) : introduction, actions menées, problème et solution, bilan.
\item \textbf{Richesse du lexique} (4 pts) : au moins \textbf{6 mots} de la leçon (réseaux sociaux, gestion de compte).
\item \textbf{Présentation et lisibilité} (4 pts) : écriture soignée, caractères de taille régulière, copie propre.
\end{enumerate}
\end{propriete}

\begin{center}
\begin{tabularx}{\linewidth}{|X|X|X|}\hline
\rowcolor{popblueL} Critère & Points & Ce que le correcteur vérifie \\ \hline
1. Caractères exacts & 4 & traits, radicaux, aucun caractère inventé \\ \hline
2. Grammaire et ordre des mots & 4 & S-temps-V-O, \zh{的}, connecteurs \\ \hline
3. Respect du plan & 4 & intro, actions, problème-solution, bilan \\ \hline
4. Lexique du thème & 4 & au moins 6 mots de la leçon \\ \hline
5. Présentation & 4 & lisibilité, soin, régularité \\ \hline
\end{tabularx}
\end{center}

\begin{popfigure}
\begin{tikzpicture}[
crit/.style={rectangle, draw=popblue, fill=popblueL, rounded corners=2pt, minimum height=0.9cm, align=left, text width=5.6cm, font=\small},
pts/.style={rectangle, draw=poporange, fill=poporangeL, rounded corners=2pt, minimum height=0.9cm, align=center, text width=1.6cm, font=\small},
chk/.style={rectangle, draw=popgreen, fill=popgreenL, rounded corners=2pt, minimum height=0.9cm, align=center, text width=3.4cm, font=\small},
head/.style={rectangle, draw=popdark, fill=popdark, text=white, rounded corners=2pt, minimum height=0.9cm, align=center, font=\small\bfseries}
]
\node[head, text width=5.6cm] (h1) at (0,0) {Critère};
\node[head, text width=1.6cm, right=0.15cm of h1] (h2) {Points};
\node[head, text width=3.4cm, right=0.15cm of h2] (h3) {Auto-contrôle};
\node[crit, below=0.12cm of h1] (c1) {1. Exactitude des caractères};
\node[pts, below=0.12cm of h2] (p1) {/4};
\node[chk, below=0.12cm of h3] (k1) {$\square$ oui \quad $\square$ non};
\node[crit, below=0.12cm of c1] (c2) {2. Ordre des mots et grammaire};
\node[pts, below=0.12cm of p1] (p2) {/4};
\node[chk, below=0.12cm of k1] (k2) {$\square$ oui \quad $\square$ non};
\node[crit, below=0.12cm of c2] (c3) {3. Respect du plan en 4 parties};
\node[pts, below=0.12cm of p2] (p3) {/4};
\node[chk, below=0.12cm of k2] (k3) {$\square$ oui \quad $\square$ non};
\node[crit, below=0.12cm of c3] (c4) {4. Lexique du thème (6 mots min.)};
\node[pts, below=0.12cm of p3] (p4) {/4};
\node[chk, below=0.12cm of k3] (k4) {$\square$ oui \quad $\square$ non};
\node[crit, below=0.12cm of c4] (c5) {5. Présentation et lisibilité};
\node[pts, below=0.12cm of p4] (p5) {/4};
\node[chk, below=0.12cm of k4] (k5) {$\square$ oui \quad $\square$ non};
\node[below=0.25cm of c5, font=\small\itshape, text=popmuted] {Total : /20 — je ne rends ma copie que si les 5 cases « oui » sont cochées.};
\end{tikzpicture}
\legende{Grille d'évaluation avec colonne d'auto-contrôle : l'élève coche chaque critère avant de rendre sa copie.}
\end{popfigure}

\courssub{L'auto-évaluation : trois relectures, pas une seule}

L'erreur classique de l'élève pressé est de relire « en général », c'est-à-dire de ne rien relire du tout. La relecture efficace est \cle{ciblée} : chaque passage sur le texte ne cherche qu'un seul type de faute.

\begin{methode}[S'auto-corriger en trois relectures]
\begin{enumerate}
\item \textbf{Première relecture — les caractères.} Vérifiez chaque caractère un à un : aucun trait manquant (\zh{我} sans le dernier trait est faux), aucun radical confondu (\zh{请} avec la clé de la parole \zh{讠}, pas \zh{清} avec la clé de l'eau \zh{氵}), aucun caractère « inventé » de mémoire approximative.
\item \textbf{Deuxième relecture — l'ordre des mots.} Soulignez le sujet de chaque phrase, puis vérifiez : le complément de temps ou de lieu est-il \emph{avant} le verbe ? Le \zh{的} est-il présent entre le possesseur et le nom (\zh{我们的账号}, « notre compte ») ?
\item \textbf{Troisième relecture — les connecteurs.} Vérifiez que chaque idée est reliée : \zh{首先} (d'abord), \zh{然后} (ensuite), \zh{但是} (mais), \zh{因为……所以……} (parce que... donc...), \zh{最后} (enfin). Un texte sans connecteur est une liste, pas un texte.
\end{enumerate}
\end{methode}

\begin{attention}[Un trait manquant = un caractère faux]
À l'examen, un caractère mal écrit — même d'un seul trait — est compté \textbf{faux}, exactement comme une faute d'orthographe en français. \zh{认真} écrit avec un trait en moins ne vaut rien, même si l'idée est bonne. La \cle{précision prime sur la vitesse} : écrivez moins de caractères, mais écrivez-les parfaitement. Méfiez-vous aussi des paires proches : \zh{未} (wèi, pas encore) et \zh{末} (mò, fin), \zh{己} (jǐ, soi-même) et \zh{已} (yǐ, déjà), \zh{买} (mǎi, acheter) et \zh{卖} (mài, vendre).
\end{attention}

\courssub{Correction collective d'une production anonyme}

Le professeur projette au tableau une production anonyme d'élève. La classe la corrige ensemble, critère par critère, avec la grille. Voici le texte proposé :

\begin{textebox}[Production anonyme à corriger en classe]
\zh{我管理我们学校的中文账号。先首我每天发照片，然后回答粉丝的问题。上个月有一个问题：一个视频太长了，粉丝不喜欢。所以我改了，现在视频很短。因为我认真，所以粉丝越来越多。我很高兴。}\\[3pt]
\emph{Wǒ guǎnlǐ wǒmen xuéxiào de Zhōngwén zhànghào. Xiān shǒu wǒ měitiān fā zhàopiàn, ránhòu huídá fěnsī de wèntí. Shàng gè yuè yǒu yí gè wèntí : yí gè shìpín tài cháng le, fěnsī bù xǐhuan. Suǒyǐ wǒ gǎi le, xiànzài shìpín hěn duǎn. Yīnwèi wǒ rènzhēn, suǒyǐ fěnsī yuè lái yuè duō. Wǒ hěn gāoxìng.}\\[3pt]
« Je gère le compte de chinois de notre école. D'abord je publie des photos chaque jour, ensuite je réponds aux questions des abonnés. Le mois dernier, il y a eu un problème : une vidéo était trop longue, les abonnés n'ont pas aimé. Alors je l'ai modifiée, maintenant les vidéos sont courtes. Parce que je suis sérieux, les abonnés sont de plus en plus nombreux. Je suis content. »
\end{textebox}

\begin{exoresolu}[Corriger la production anonyme avec la grille]
Énoncé : attribuez une note sur 20 à cette production, critère par critère, et justifiez.
\tcblower
\textbf{Solution détaillée.}
\begin{itemize}
\item \textbf{Critère 1 — Caractères (3/4)} : tous les caractères sont exacts sauf \zh{先首}, inversion fautive de \zh{首先} (shǒuxiān, « d'abord »). Un mot mal écrit : on retire 1 point. Notez aussi l'absence de classificateur pour \zh{视频} (devrait être \zh{一条视频}) et \zh{问题} (devrait être \zh{一个问题} ou \zh{一条问题}), toléré ici mais à corriger.
\item \textbf{Critère 2 — Grammaire (4/4)} : ordre Sujet + temps + Verbe + Objet respecté (\zh{我每天发照片}) ; \zh{的} correctement employé (\zh{我们学校的中文账号}, \zh{粉丝的问题}) ; \zh{了} bien placé dans \zh{我改了} et \zh{太长了}.
\item \textbf{Critère 3 — Plan (4/4)} : introduction (\zh{我管理……账号}), actions (photos, réponses), problème-solution (vidéo trop longue → modifiée), bilan (\zh{我很高兴}). Les quatre parties sont présentes.
\item \textbf{Critère 4 — Lexique (4/4)} : au moins 7 mots du thème : \zh{管理}, \zh{账号}, \zh{发}, \zh{粉丝}, \zh{回答}, \zh{视频}, \zh{问题}. Objectif dépassé.
\item \textbf{Critère 5 — Présentation (4/4)} : copie lisible, ponctuation chinoise correcte (\zh{。} et \zh{：}).
\end{itemize}
\textbf{Total : 19/20.} Remarque collective : la correction de \zh{先首} en \zh{首先} et l'ajout des classificateurs (\zh{一条视频}) auraient valu 20/20 — d'où l'importance de la première relecture consacrée aux caractères et de la deuxième à la structure.
\end{exoresolu}

\courssub{Production finale : consignes et rappel pour l'examen}

\begin{exercice}[Production finale — à rendre]
Rédigez un texte de \textbf{6 à 8 phrases} en chinois sur le sujet suivant : « Je gère le compte d'un club de mon lycée (sport, musique ou chinois) sur les réseaux sociaux. » Votre texte doit contenir : une introduction (qui, quel compte), deux actions régulières, un problème rencontré et sa solution, un bilan. Utilisez au moins 6 mots du lexique de la leçon et au moins 3 connecteurs. Avant de rendre, cochez les 5 cases de la grille d'auto-contrôle.
\end{exercice}

\begin{corrige}[Production finale — éléments de correction]
Barème : grille des 5 critères (20 points). Exemple de texte attendu :\\[2pt]
\zh{我管理我们学校足球俱乐部的账号。首先我每天发照片和视频，然后回答粉丝的问题。上个月有个问题：我发的消息太少了，粉丝不高兴。所以我决定每天发一条消息。因为我很认真，所以我们的粉丝越来越多。我觉得这份工作很有意思。}\\[2pt]
\emph{Wǒ guǎnlǐ wǒmen xuéxiào zúqiú jùlèbù de zhànghào. Shǒuxiān wǒ měitiān fā zhàopiàn hé shìpín, ránhòu huídá fěnsī de wèntí. Shàng gè yuè yǒu gè wèntí : wǒ fā de xiāoxi tài shǎo le, fěnsī bú gāoxìng. Suǒyǐ wǒ juédìng měitiān fā yī tiáo xiāoxi. Yīnwèi wǒ hěn rènzhēn, suǒyǐ wǒmen de fěnsī yuè lái yuè duō. Wǒ juéde zhè fèn gōngzuò hěn yǒu yìsi.}\\[2pt]
« Je gère le compte du club de football de notre école. D'abord, je publie chaque jour des photos et des vidéos, ensuite je réponds aux questions des abonnés. Le mois dernier, il y a eu un problème : je publiais trop peu de messages, les abonnés n'étaient pas contents. J'ai donc décidé de publier un message par jour. Parce que je suis très sérieux, nos abonnés sont de plus en plus nombreux. Je trouve ce travail très intéressant. »\\[2pt]
Points clés à vérifier : \zh{首先……然后……所以……} (3 connecteurs) ; lexique du thème : \zh{管理、账号、发、视频、粉丝、消息} (6 mots) ; plan en 4 parties respecté ; classificateurs corrects (\zh{一条视频、一条消息、这份工作}) ; tons sandhi appliqués (\zh{yí gè, yī tiáo, bú gāoxìng}).
\end{corrige}

\begin{attention}[Complément utile pour l'examen]
Au Baccalauréat, un texte de \textbf{5 à 8 phrases bien structuré} vaut toujours mieux qu'un long texte plein de fautes. Le correcteur note la maîtrise, pas la quantité : dix phrases avec cinq caractères faux rapportent moins que six phrases impeccables. Écrivez court, écrivez juste, écrivez propre.
\end{attention}

\begin{savaistu}[Le mot \zh{粉丝} vient de l'anglais]
Le mot chinois \zh{粉丝} (fěnsī, « abonnés, fans ») est un emprunt phonétique à l'anglais \emph{fans} ! Ses caractères désignent à l'origine les vermicelles de fécule (aussi appelés \zh{粉丝} ou \zh{冬粉}), très populaires en cuisine chinoise. C'est un bel exemple de la façon dont le chinois moderne absorbe les mots étrangers du numérique par transcription phonétique — un phénomène que l'on observe aussi au Cameroun, où le français et l'anglais empruntent aux langues locales (ex. \emph{ndolé}, \emph{bikutsi}) et inversement.
\end{savaistu}

\newpage

\coursec{Activité d'intégration}

\coursec{Leçon 42 — Expression écrite : caractères complexes liés à la gestion des réseaux sociaux}

\courssub{Objectifs de l'activité}
\par
Cette activité de synthèse vise à mobiliser l'ensemble des acquis de la leçon 42 dans une tâche d'expression écrite authentique et contextualisée. À l'issue de cette séance, l'élève doit être capable de :
\begin{itemize}
    \item \textbf{Réinvestir} le lexique spécialisé de la gestion des réseaux sociaux (noms d'actions, termes techniques, indicateurs) dans une production écrite cohérente.
    \item \textbf{Écrire correctement} les caractères complexes étudiés (ordre des traits, radicaux, structure) sans hésitation ni confusion avec des caractères proches.
    \item \textbf{Structurer} un texte court (6--8 phrases) à l'aide de connecteurs logiques (chronologie, cause-conséquence, concession) pour organiser la narration d'une journée type.
    \item \textbf{Traduire} mentalement du français vers le chinois (thème) en respectant l'ordre des mots chinois (Sujet -- Temps -- Lieu -- Verbe -- Complément) et la place des compléments du verbe, notamment la structure \zh{把} pour l'objet défini.
    \item \textbf{S'auto-évaluer et évaluer un pair} à l'aide d'une grille critériée (orthographe des caractères, justesse grammaticale, richesse lexicale, cohérence textuelle).
\end{itemize}

\courssub{Situation}
\par
Vous êtes stagiaire \textbf{community manager} pour le compte officiel WeChat (\zh{微信公众号}) de la boutique \emph{«~Savoir-Faire de l'Ouest~»}, située au quartier \emph{Mvog-Ada} à Yaoundé. Cette boutique vend des pagnes tissés de Bafoussam, du café arabica des hauteurs de Dschang et de l'artisanat en bronze de Foumban. Le gérant, M. \textsc{Fotso}, vous confie la gestion de la page pour la journée du \emph{Samedi 15 Juin}, jour de forte affluence car c'est la veille de la fête du \emph{Ngon} (fête traditionnelle Bamiléké).
\begin{textebox}[Brief du Community Manager -- Contexte authentique]
\zh{情况说明：} (Qíngkuàng shuōmíng) -- Situation : \\
\zh{账号名称：} \zh{「西部匠心」} (Xībù Jiàngxīn) -- «~Savoir-Faire de l'Ouest~» \\
\zh{日期：} \zh{六月十五日，星期六} (Liùyuè shíwǔ rì, Xīngqīliù) -- 15 juin, samedi \\
\zh{任务：} \zh{作为实习社群运营，你需要完成一天的标准运营流程：} (En tant que stagiaire community manager, tu dois accomplir le processus standard d'une journée :) \\
\zh{1. 早上发布新品预告（新品预告 xīnpǐn yùgào）。} -- Matin : publier l'avant-première nouveau produit. \\
\zh{2. 中午回复粉丝私信和评论（回复 私信 评论）。} -- Midi : répondre aux MP et commentaires des fans. \\
\zh{3. 下午处理垃圾广告：删除、屏蔽恶意账号（删除 垃圾广告，屏蔽 恶意账号）。} -- Après-midi : traiter le spam : supprimer, bloquer comptes malveillants. \\
\zh{4. 晚上看数据分析，写简短总结（数据分析，总结）。} -- Soir : consulter l'analyse de données, écrire un bref bilan. \\
\zh{要求：} \zh{用中文写一篇 6--8 句的「工作日志」，发在团队群里。} (Écrire en chinois un « journal de bord » de 6-8 phrases, le poster dans le groupe d'équipe.)
\end{textebox}
\par
Votre mission : rédiger le \textbf{journal de bord} (\zh{工作日志}) de cette journée en chinois simplifié, à destination de l'équipe de direction (M. Fotso et la gérante adjointe, Mme \textsc{Ngouo}). Ce texte servira de rapport d'activité officiel.

\courssub{Consignes}
\par
Travaillez \textbf{par binômes}. Chaque binôme produit \textbf{un seul texte manuscrit} (sur feuille A4, caractères soignés) ou \textbf{numérique} (traitement de texte avec saisie pinyin). Suivez scrupuleusement les étapes ci-dessous :

\begin{enumerate}
    \item \textbf{Analyse du brief (5 min) :} Identifiez les 4 moments clés de la journée (Matin, Midi, Après-midi, Soir) et les verbes d'action imposés (\zh{发布}, \zh{回复}, \zh{删除}, \zh{屏蔽}, \zh{分析}, \zh{总结}). Soulignez les mots de vocabulaire obligatoires à caser (minimum 6).
    \item \textbf{Planification et connecteurs (5 min) :} Choisissez 3 connecteurs logiques parmi la liste de la leçon pour relier les 4 moments. Exemple : \zh{首先}... \zh{接着}... \zh{因为}... \zh{所以}... \zh{最后}... Notez la structure de chaque phrase sur brouillon (Sujet + Temps + Lieu + Verbe + Objet).
    \item \textbf{Rédaction du texte (15 min) :} Rédigez le \zh{工作日志} (6 à 8 phrases). \textbf{Contraintes strictes :}
        \begin{itemize}
            \item Utiliser \textbf{au moins 6 mots} du tableau de vocabulaire de la leçon (voir Ressources).
            \item Utiliser \textbf{au moins 3 connecteurs} différents.
            \item Respecter l'ordre des mots : \cle{Sujet + Temps + Lieu + Verbe + Objet}.
            \item Écrire les caractères complexes (\zh{发布}, \zh{删除}, \zh{屏蔽}, \zh{分析}, \zh{垃圾}, \zh{账号}, \zh{粉丝}, \zh{互动}) avec un trait soigné (ordre des traits respecté).
        \end{itemize}
    \item \textbf{Correction croisée (10 min) :} Échangez votre feuille avec un autre binôme. Utilisez la \textbf{Grille d'évaluation} (fournie ci-dessous) pour noter le texte de vos camarades : cochez les cases, entourez les erreurs de caractères au stylo rouge, proposez une correction en marge. Signez en bas : «~Corrigé par [Vos Noms]~».
    \item \textbf{Rédaction finale et affichage (5 min) :} Récupérez votre texte, corrigez les erreurs signalées. Les 3 meilleurs textes (note $\ge$ 16/20) sont photocopiés et affichés au «~Mur des Publications~» (\zh{优秀作品展示墙}) comme «~Page Officielle du Jour~».
\end{enumerate}

\courssub{Ressources à mobiliser}
\par
Voici le rappel synthétique des outils linguistiques vus dans les séquences 1 à 4 de cette leçon. Gardez cette page ouverte pendant la rédaction.

\begin{definition}[Lexique clé -- Gestion des réseaux sociaux (HSK 2-3)]
\begin{center}
\begin{tabularx}{\linewidth}{|X|X|X|X|}
\hline
\rowcolor{popblueL} \textbf{Caractères} & \textbf{Pinyin} & \textbf{Nature} & \textbf{Traduction} \\ \hline
账号 & zhànghào & n. & compte, identifiant \\ \hline
发布 & fābù & v. & publier, diffuser \\ \hline
内容 & nèiróng & n. & contenu \\ \hline
粉丝 & fěnsī & n. & fans, abonnés \\ \hline
互动 & hùdòng & v./n. & interagir / interaction \\ \hline
回复 & huífù & v. & répondre, répliquer \\ \hline
评论 & pínglùn & v./n. & commenter / commentaire \\ \hline
私信 & sīxìn & n. & message privé (MP) \\ \hline
删除 & shānchú & v. & supprimer, effacer \\ \hline
屏蔽 & bǐngbì & v. & bloquer, masquer (utilisateur/contenu) \\ \hline
垃圾信息 & lājī xìnxī & n. & spam, pourriel (litt. «~infos poubelle~») \\ \hline
恶意账号 & èyì zhànghào & n. & compte malveillant \\ \hline
点赞 & diǎnzàn & v. & liker, mettre un «~j'aime~» \\ \hline
转发 & zhuǎnfā & v. & partager, retweeter \\ \hline
数据分析 & shùjù fēnxī & n. & analyse de données \\ \hline
总结 & zǒngjié & v./n. & résumer / résumé, bilan \\ \hline
推广 & tuīguǎng & v. & promouvoir, faire la promo \\ \hline
关注 & guānzhù & v. & suivre, s'abonner \\ \hline
设置 & shèzhì & v. & configurer, paramétrer \\ \hline
管理员 & guǎnlǐyuán & n. & administrateur, modérateur \\ \hline
后台 & hòutái & n. & back-office, interface d'admin \\ \hline
\end{tabularx}
\end{center}
\end{definition}

\begin{propriete}[Connecteurs logiques pour structurer le récit]
\begin{center}
\begin{tabularx}{\linewidth}{|X|X|X|}
\hline
\rowcolor{popgreenL} \textbf{Fonction} & \textbf{Connecteur (Caractères | Pinyin)} & \textbf{Place dans la phrase} \\ \hline
Chronologie (début) & \zh{首先} (shǒuxiān) & Début de phrase, suivi de \zh{，} \\ \hline
Chronologie (suite) & \zh{接着} (jiēzhe) / \zh{然后} (ránhòu) & Début de phrase ou après sujet \\ \hline
Chronologie (fin) & \zh{最后} (zuìhòu) & Début de phrase ou après sujet \\ \hline
Cause -- Conséquence & \zh{因为}... \zh{所以}... (yīnwèi... suǒyǐ...) & \zh{因为} clause 1, \zh{所以} clause 2 \\ \hline
Concession & \zh{虽然}... \zh{但是}... (suīrán... dànshì...) & \zh{虽然} clause 1, \zh{但是} clause 2 \\ \hline
Progressif & \zh{不但}... \zh{而且}... (búdàn... érqiě...) & \zh{不但} clause 1, \zh{而且} clause 2 \\ \hline
Condition & \zh{如果}... \zh{就}... (rúguǒ... jiù...) & \zh{如果} clause 1, \zh{就} clause 2 \\ \hline
\end{tabularx}
\end{center}
\end{propriete}

\begin{propriete}[Rappel : Ordre des mots de base (SVO + Compléments) et structure \zh{把}]
\par
En chinois, l'ordre canonique est immuable : 
\zh{\textbf{Sujet} + \textbf{Temps} + \textbf{Lieu} + \textbf{Verbe} + \textbf{Objet}}.
\end{propriete}

\begin{exemplebox}[Exemples d'application]
\zh{我 早上 在公司 发布 内容。} (Wǒ zǎoshang zài gōngsī fābù nèiróng.) -- Moi, matin, à l'entreprise, publier contenu. \\
\zh{因为 垃圾信息 多， 所以 我 删除了 它们。} (Yīnwèi lājī xìnxī duō, suǒyǐ wǒ shānchúle tāmen.) -- Parce que spams nombreux, donc moi les ai supprimés.
\end{exemplebox}
\par
\cle{Attention :} Le complément de lieu (\zh{在微信上}, \zh{在后台}) se place \textbf{avant} le verbe. Le particle \zh{了} (action accomplie) se place \textbf{après} le verbe (ou en fin de phrase si objet long).
\par
\cle{Structure \zh{把} (objet défini + action aboutie) :} \zh{Sujet + 把 + Objet + Verbe + 了 / 结果/方向}. Exemple : \zh{我 把 垃圾信息 删除 了。} (J'ai supprimé les spams). Obligatoire si l'objet est défini (\zh{它, 它们, 这个}) et le verbe marque une manipulation/résultat.


\begin{attention}[Pièges fréquents pour francophones -- Caractères complexes]
\begin{itemize}
    \item \zh{发布} vs \zh{发} (seul) : \zh{发布} = publier officiellement (article, vidéo, communiqué) ; \zh{发} = envoyer (message, colis, WeChat). Ne pas écrire \zh{发布了一个微信} (incorrect), dire \zh{发了一个微信} ou \zh{发布了一篇文章}.
    \item \zh{删除} vs \zh{屏蔽} : \zh{删除} = supprimer le contenu (effacer un message) ; \zh{屏蔽} = bloquer la source (utilisateur, mot-clé, IP) pour qu'il ne réapparaisse pas. Pinyin : \zh{屏蔽} se prononce \textbf{bǐngbì} (3e + 4e ton), pas *bìngbì.
    \item \zh{垃圾信息} : \zh{垃圾} (poubelle) + \zh{信息} (info). Ne pas confondre \zh{圾} (jí, radical 土) et \zh{基} (jī, base, radical 土) ni \zh{信} (xìn, message, radical 亻) et \zh{伸} (shēn, étendre, radical 亻).
    \item \zh{粉丝} : \zh{粉} (poudre/farine, radical 米 en haut) + \zh{丝} (fil/soie, 6 traits : piě, zhé, zhé, diǎn, piě, nà). Ordre : haut → bas, gauche → droite.
    \item \zh{互动} : \zh{互} (mutuel, radical \zh{彳}, 4 traits) + \zh{动} (bouger, radical \zh{力}). Ne pas écrire \zh{动} à la place de \zh{互}.
    \item \zh{分析} : \zh{分} (partager, 8 traits, radical 刀) + \zh{析} (analyser, radical \zh{木}, 8 traits). Symétrie visuelle : les deux font 8 traits.
\end{itemize}
\end{attention}

\courssub{Proposition de production et corrigé commenté}
\par
Voici une production modèle respectant toutes les contraintes (7 phrases, 9 mots de vocabulaire, 5 connecteurs). Elle sert de référence pour l'évaluation.

\begin{exoresolu}[Journal de bord modèle -- \zh{「西部匠心」} 运营日志]
\textbf{Énoncé :} Rédiger le \zh{工作日志} de la journée du 15 juin pour le compte \zh{「西部匠心」}.

\tcblower

\textbf{Production modèle (Chinois + Pinyin + Traduction)} \\

\zh{1. 首先，我 早上 十点 在 后台 发布 了 新品 预告。} \\
\emph{Shǒuxiān, wǒ zǎoshang shí diǎn zài hòutái fābùle xīnpǐn yùgào.} \\
«~Premièrement, moi matin 10h au back-office ai publié [accompli] nouveau-produit avant-première.~» \\

\zh{2. 接着，我 回复 了 所有 粉丝 的 私信 和 评论。} \\
\emph{Jiēzhe, wǒ huífùle suǒyǒu fěnsī de sīxìn hé pínglùn.} \\
«~Ensuite, moi ai répondu [accompli] tous fans DE messages-privés et commentaires.~» \\

\zh{3. 中午，互动 量 很高，不但 有 点赞，而且 有 很多 转发。} \\
\emph{Zhōngwǔ, hùdòngliàng hěn gāo, búdàn yǒu diǎnzàn, érqiě yǒu hěnduō zhuǎnfā.} \\
«~Midi, quantité-interaction très-haute, non-seulement il-y-a likes, mais-aussi il-y-a beaucoup partages.~» \\

\zh{4. 但是，我 下午 发现 了 很多 垃圾 信息 和 恶意 账号。} \\
\emph{Dànshì, wǒ xiàwǔ fāxiànle hěnduō lājī xìnxī hé èyì zhànghào.} \\
«~Mais, moi après-midi ai découvert [accompli] beaucoup spams et malveillants comptes.~» \\

\zh{5. 因为 这些 账号 影响 体验，所以 我 把 它们 都 删除 了，还 屏蔽 了 相关 关键词。} \\
\emph{Yīnwèi zhèxiē zhànghào yǐngxiǎng tǐyàn, suǒyǐ wǒ bǎ tāmen dōu shānchúle, hái bǐngbìle xiāngguān guānjiàncí.} \\
«~Parce que ces comptes affectent expérience, donc moi (BA) elles toutes ai-supprimés, aussi ai-bloqué [accompli] relatifs mots-clés.~» \\

\zh{6. 最后，我 看 了 数据 分析，关注 量 增加 了 两百 人。} \\
\emph{Zuìhòu, wǒ kànle shùjù fēnxī, guānzhùliáng zēngjiāle liǎngbǎi rén.} \\
«~Finalement, moi ai-regardé [accompli] données analyse, quantité-abonnés augmentée [accompli] deux-cents personnes.~» \\

\zh{7. 总之，今天 的 运营 很 顺利，明天 继续 推广 咖啡 和 青铜 工艺品。} \\
\emph{Zòngguī, jīntiān de yùnyíng hěn shùnlì, míngtiān jìxù tuīguǎng kāfēi hé qīngtóng gōngyǐnpǐn.} \\
«~En-bref, aujourd'hui DE gestion très fluide, demain continuer promouvoir café et bronze artisanat.~» \\

\textbf{Commentaire détaillé des choix de langue (pour l'enseignant)} \\
\begin{enumerate}
    \item \textbf{Respect des contraintes quantitatives :} 7 phrases. Vocabulaire utilisé (10) : \zh{发布, 粉丝, 私信, 评论, 互动, 点赞, 转发, 垃圾信息, 恶意账号, 删除, 屏蔽, 数据分析, 关注, 推广, 后台} ($\ge$ 6). Connecteurs utilisés (5) : \zh{首先, 接着, 不但...而且..., 因为...所以..., 最后} ($\ge$ 3).
    \item \textbf{Structure \zh{把} (phrase 5) :} \zh{我把它们都删除了} est la structure correcte pour un objet défini (\zh{它们}) avec un verbe d'action aboutie (\zh{删除}) + \zh{了}. C'est un point de grammaire HSK 3 essentiel à valider.
    \item \textbf{Structure \zh{不但...而且...} (phrase 3) :} Permet de lister deux aspects positifs (\zh{点赞} et \zh{转发}) de l'\zh{互动量}. Notez l'absence de sujet répété devant \zh{而且}.
    \item \textbf{Gestion du temps / aspect :} \zh{了} marque l'accompli sur \zh{发布, 回复, 发现, 删除, 屏蔽, 看, 增加}. Pas de \zh{了} sur \zh{继续} (futur) ni sur \zh{影响} (verbe d'état dans clause \zh{因为}).
    \item \textbf{Contexte camerounais :} \zh{咖啡} (café de l'Ouest), \zh{青铜工艺品} (bronzes de Foumban) ancrent le texte. \zh{后台} (back-office) est le terme technique standard pour l'interface d'administration WeChat.
    \item \textbf{Caractères cibles :} \zh{发布, 删除, 屏蔽, 垃圾, 互动, 分析, 粉丝, 账号} sont tous présents et doivent être écrits parfaitement.
\end{enumerate}
\end{exoresolu}

\courssub{Bilan de l'activité}
\par
Cette activité de «~Community Manager d'un jour~» constitue l'évaluation sommative de la leçon 42. Elle permet de vérifier le transfert des acquis lexicaux, graphiques et grammaticaux dans une situation de communication professionnelle simulée, ancrée dans le réel économique et culturel camerounais.

\begin{methode}[Grille d'évaluation critériée -- Sur 20 points -- Pour correction croisée]
\par
Remettez cette grille à chaque binôme correcteur. Chaque critère est noté sur 4 points (0 = absent/incorrect, 1 = insuffisant, 2 = acceptable, 3 = bon, 4 = excellent).
\begin{center}
\begin{tabularx}{\linewidth}{|X|c|c|c|c|c|}
\hline
\rowcolor{poporangeL} \textbf{Critère} & \textbf{0} & \textbf{1} & \textbf{2} & \textbf{3} & \textbf{4} \\ \hline
\textbf{1. Caractères cibles} (8 mots imposés : \zh{发布 删除 屏蔽 垃圾 互动 分析 粉丝 账号}) & $>3$ erreurs & 3 erreurs & 2 erreurs & 1 erreur & 0 erreur, traits parfaits \\ \hline
\textbf{2. Richesse lexicale} (Total mots leçon utilisés) & $<4$ & 4--5 & 6 & 7 & $\ge 8$ \\ \hline
\textbf{3. Connecteurs logiques} (Variété et justesse) & 0 ou mal utilisés & 1 seul & 2 corrects & 3 corrects & $\ge 3$ variés \& justes \\ \hline
\textbf{4. Grammaire / Ordre mots} (SVO, Lieu avant V, \zh{了}, \zh{把}, \zh{因为...所以...}) & $>4$ fautes graves & 3--4 fautes & 2 fautes mineures & 1 faute mineure & Parfait \\ \hline
\textbf{5. Cohérence / Situation} (Respect brief, 4 moments, ton pro, contexte CM) & Hors sujet & Incomplet (2 moments) & 3 moments, ton moyen & 4 moments, ton pro & 4 moments, ton pro, style fluide \\ \hline
\end{tabularx}
\end{center}
\textbf{Total :} \_\_\_ / 20 \quad \textbf{Seuil affichage «~Page Officielle~» :} $\ge$ 16/20. \quad \textbf{Signature correcteurs :} \underline{\hspace{5cm}}
\end{methode}

\begin{savaistu}[Le saviez-vous ? -- WeChat (\zh{微信}) en Chine vs WhatsApp/Facebook au Cameroun]
\par
En Chine, \zh{微信} (Wēixìn, «~micro-message~») n'est pas qu'une messagerie : c'est un \textbf{super-application} (\zh{超级应用}). On y paie (\zh{微信支付}), on prend le métro, on réserve l'hôpital, on déclare ses impôts, et les entreprises y gèrent leur CRM via les \zh{公众号} (Comptes Officiels -- Service ou Abonnement) et les \zh{小程序} (Mini-programmes).
\par
Au Cameroun, l'écosystème est fragmenté : WhatsApp (messagerie + Business + Paiement naissant), Facebook (visibilité + Marketplace), Mobile Money (MTN MoMo / Orange Money) pour le paiement. Le \zh{社群运营} (Community Management) au Cameroun demande donc de \textbf{multi-canaux}, alors qu'en Chine, tout se concentre souvent dans l'écosystème WeChat/Tencent. C'est une différence structurelle majeure pour un gestionnaire de réseaux sociaux biculturel !
\end{savaistu}

\begin{experience}[Prolongement oral -- Débat : «~Gérer une crise de \zh{垃圾信息}~»]
\par
\emph{Consigne :} En grand groupe, sur la base des textes affichés. M. Fotso (le prof) annonce : «~Hier soir, un concurrent a lancé une attaque de bots : 500 commentaires publicitaires pour des paris sportifs en 10 minutes. Mon \zh{后台} a planté. Que faites-vous \textbf{maintenant} ?~»
\par
Les binômes doivent proposer une \textbf{procédure d'urgence en 3 étapes} en chinois oral (vocabulaire : \zh{紧急处理}, \zh{一键清理}, \zh{设置关键词拦截}, \zh{向平台投诉}, \zh{安抚粉丝}). Durée : 10 min. Évaluation : réactivité, justesse terminologique, prononciation des tons.
\end{experience}

\newpage

\coursec{Méthodes et exercices résolus}

\coursec{Leçon 42 — Expression écrite : caractères complexes liés à la gestion des réseaux sociaux}

\courssub{Savoir-faire 1 : Produire un texte simple sur la gestion des réseaux sociaux}

\begin{methode}[Rédiger un texte simple sur les réseaux sociaux]
Pour produire un texte court (8 à 12 phrases) sur la gestion des réseaux sociaux, suis cette démarche :
\begin{enumerate}
\item \cle{Analyser le sujet} : souligne le thème (réseaux sociaux), le destinataire (un ami, un professeur) et le type de texte demandé (message, paragraphe d'opinion, récit).
\item \cle{Faire un plan en trois temps} : introduction (je me présente ou je pose le sujet avec \zh{现在} « maintenant » ou \zh{随着……} « avec le développement de… »), développement (2 ou 3 idées reliées par des connecteurs), conclusion (avis personnel avec \zh{我觉得} « je pense que »).
\item \cle{Mobiliser le lexique du thème} : \zh{社交媒体} « réseaux sociaux », \zh{发布} « publier », \zh{评论} « commenter », \zh{点赞} « liker », \zh{关注} « suivre / s'abonner », \zh{分享} « partager », \zh{隐私} « vie privée », \zh{账号} « compte », \zh{密码} « mot de passe ».
\item \cle{Rédiger avec des structures sûres} : phrases Sujet + Verbe + Complément ; \zh{因为……所以……} « parce que… donc… » ; \zh{虽然……但是……} « bien que… mais… » ; \zh{不但……而且……} « non seulement… mais aussi… ».
\item \cle{Relire} : vérifier les tons, les caractères, la place du complément de temps (avant le verbe) et l'emploi correct de \zh{了}.
\end{enumerate}
Réflexe : un connecteur par phrase au minimum ; jamais deux verbes principaux sans lien.
\end{methode}

\begin{definition}[Lexique de la gestion des réseaux sociaux — HSK 3-4]
\begin{center}
\begin{tabularx}{\linewidth}{|X|X|X|X|}
\hline
\rowcolor{popblueL} \textbf{Caractères} & \textbf{Pinyin} & \textbf{Nature} & \textbf{Traduction} \\ \hline
社交媒体 & shèjiāo méitǐ & nom & réseaux sociaux \\ \hline
账号 & zhànghào & nom & compte (utilisateur) \\ \hline
密码 & mìmǎ & nom & mot de passe \\ \hline
发布 & fābù & verbe & publier (contenu) \\ \hline
分享 & fēnxiǎng & verbe & partager \\ \hline
评论 & pínglùn & verbe / nom & commenter / commentaire \\ \hline
点赞 & diǎnzàn & verbe & liker, mettre « j'aime » \\ \hline
关注 & guānzhù & verbe & suivre, s'abonner à \\ \hline
隐私 & yǐnsī & nom & vie privée, intimité \\ \hline
视频 & shìpín & nom & vidéo \\ \hline
照片 & zhàopiàn & nom & photo \\ \hline
消息 & xiāoxi & nom & message, nouvelle \\ \hline
保护 & bǎohù & verbe & protéger \\ \hline
设置 & shèzhì & verbe & configurer, paramétrer \\ \hline
登录 & dēnglù & verbe & se connecter (log in) \\ \hline
注销 & zhùxiāo & verbe & se déconnecter / supprimer un compte \\ \hline
\end{tabularx}
\end{center}
\end{definition}

\begin{exoresolu}[Production : mon compte sur les réseaux sociaux]
Énoncé. Rédige un texte de 6 à 8 phrases en chinois : tu présentes ton compte de réseau social, ce que tu y publies, et ton avis sur l'usage des réseaux sociaux par les lycéens camerounais.
\tcblower
Solution détaillée.

\textbf{Étape 1 — Plan.} Introduction : j'ai un compte. Développement : ce que je publie et pourquoi. Conclusion : avis nuancé (avantages + prudence).

\textbf{Étape 2 — Choix des structures.} \zh{我有一个……账号} (présentation) ; \zh{我常常发布……} (habitude : \zh{常常} + verbe) ; \zh{虽然……但是……} (avis nuancé) ; \zh{我觉得……} (conclusion).

\textbf{Étape 3 — Rédaction.}

\zh{我有一个社交媒体账号，名字叫“小杜阿拉”。我常常发布我们学校足球队的照片，也分享喀麦隆的音乐。很多朋友给我点赞和评论，我很高兴。虽然社交媒体很有意思，但是我们不应该花太多时间。我觉得学生应该先学习，后玩手机。}

\emph{Wǒ yǒu yí ge shèjiāo méitǐ zhànghào, míngzi jiào “Xiǎo Dù'ālā”. Wǒ chángcháng fābù wǒmen xuéxiào zúqiú duì de zhàopiàn, yě fēnxiǎng Kāmàilóng de yīnyuè. Hěn duō péngyou gěi wǒ diǎnzàn hé pínglùn, wǒ hěn gāoxìng. Suīrán shèjiāo méitǐ hěn yǒu yìsi, dànshì wǒmen bù yīnggāi huā tài duō shíjiān. Wǒ juéde xuésheng yīnggāi xiān xuéxí, hòu wán shǒujī.}

« J'ai un compte de réseau social, qui s'appelle “Petit Douala”. Je publie souvent des photos de l'équipe de football de notre école, et je partage aussi de la musique camerounaise. Beaucoup d'amis me likent et me commentent, je suis très content. Bien que les réseaux sociaux soient très intéressants, nous ne devrions pas y passer trop de temps. Je pense que les élèves devraient d'abord étudier, ensuite jouer avec le téléphone. »

\textbf{Étape 4 — Vérification.} Classificateur \zh{个} devant \zh{账号} ; \zh{常常} placé avant le verbe \zh{发布} ; \zh{虽然……但是……} bien appariés ; pas de \zh{了} superflu (habitudes au présent). Conclusion : le texte respecte le plan, emploie 6 mots du lexique du thème et deux connecteurs : objectif atteint.
\end{exoresolu}

\courssub{Savoir-faire 2 : Écrire correctement les caractères complexes du thème}

\begin{methode}[Maîtriser traits, radicaux et caractères proches]
\begin{enumerate}
\item \cle{Identifier le radical (clé)} avant d'écrire : il donne le sens et souvent la position. Exemples du thème : \zh{讠} (parole) dans \zh{评} « commenter », \zh{讠} dans \zh{论} « discuter », \zh{心} (cœur) dans \zh{意} « idée », \zh{扌} (main) dans \zh{拍} « photographier » et \zh{按} « appuyer ».
\item \cle{Respecter l'ordre des traits} : de haut en bas, de gauche à droite, l'horizontal avant le vertical, l'extérieur avant l'intérieur, on ferme la boîte en dernier (\zh{园}, \zh{图}).
\item \cle{Décomposer les caractères complexes} : \zh{赞} = \zh{先} + \zh{先} + \zh{贝} (coquillage = valeur) ; \zh{隐} = \zh{阝} (colline gauche) + \zh{急} ; \zh{密} = \zh{宀} (toit) + \zh{必} + \zh{山}. Écrire composant par composant évite les oublis.
\item \cle{Distinguer les caractères proches (homophones ou quasi-homophones)} :
\begin{itemize}
\item \zh{密} (mì, secret, radical \zh{宀} toit) dans \zh{密码} vs \zh{蜜} (mì, miel, radical \zh{虫} insecte).
\item \zh{账} (zhàng, compte, radical \zh{贝} coquillage/argent) dans \zh{账号} vs \zh{帐} (zhàng, tente, radical \zh{巾} tissu).
\item \zh{论} (lùn, discuter, radical \zh{讠} parole) dans \zh{评论} vs \zh{伦} (lún, éthique, radical \zh{亻} personne).
\item \zh{赞} (zàn, louer/liker, radical \zh{贝} valeur) vs \zh{暂} (zàn, temporaire, radical \zh{日} soleil).
\end{itemize}
\item \cle{S'entraîner par séries} : écrire chaque caractère 5 fois en prononçant le pinyin à voix haute, puis l'employer dans une phrase.
\end{enumerate}
\end{methode}

\begin{popfigure}
\begin{tikzpicture}[scale=0.8, every node/.style={font=\small}]
% Décomposition de 赞
\node[draw=popblue, thick, rounded corners, fill=popblueL, align=center] (zan) at (0,0) {\textbf{赞} (zàn) — liker, louer\\[2pt] Radical : \zh{贝} (bèi, coquillage/valeur)};
\node[draw=poporange, thick, rounded corners, fill=poporangeL, align=center] (xian1) at (-2.5,-2) {\zh{先} (xiān)\\[2pt] haut : \zh{⺧} (piě)\\[2pt] bas : \zh{儿} (ér)};
\node[draw=poporange, thick, rounded corners, fill=poporangeL, align=center] (xian2) at (2.5,-2) {\zh{先} (xiān)};
\node[draw=popgreen, thick, rounded corners, fill=popgreenL, align=center] (bei) at (0,-4) {\zh{贝} (bèi)\\[2pt] coquillage / argent};
\draw[->, thick, popdark] (zan.south west) -- (xian1.north east);
\draw[->, thick, popdark] (zan.south east) -- (xian2.north west);
\draw[->, thick, popdark] (xian1.south) -- ++(0,-0.3) -| (bei.north west);
\draw[->, thick, popdark] (xian2.south) -- ++(0,-0.3) -| (bei.north east);
\node[align=center, text=popdark] at (0,-5.5) {Structure : \zh{先} (haut-gauche) + \zh{先} (haut-droit) + \zh{贝} (bas)};
\end{tikzpicture}
\legende{Analyse structurelle du caractère \zh{赞} (zàn) — radical \zh{贝} « valeur », double \zh{先} « premier/avant » suggérant l'approbation anticipée.}
\end{popfigure}

\begin{exoresolu}[Écriture : corriger les caractères fautifs]
Énoncé. Un élève a écrit la phrase suivante avec trois caractères faux. Retrouve-les, corrige-les et explique la confusion : \zh{我的帐号的蜜码很简单，别人很容易评伦我的照片。}
\tcblower
Solution détaillée.

\textbf{Erreur 1 : \zh{帐} au lieu de \zh{账}.} \zh{账号} (zhànghào, « compte ») s'écrit avec \zh{账}, radical \zh{贝} (coquillage, symbole d'argent : un compte concerne les biens). \zh{帐} avec \zh{巾} (tissu) signifie « tente, moustiquaire ». Même prononciation zhàng, sens différent : c'est un homophone-piège.

\textbf{Erreur 2 : \zh{蜜} au lieu de \zh{密}.} \zh{密码} (mìmǎ, « mot de passe ») s'écrit avec \zh{密} : \zh{宀} (toit) + \zh{必} + \zh{山}, idée de « caché, secret ». \zh{蜜} avec \zh{虫} (insecte) en bas signifie « miel » (l'abeille !). Homophones parfaits : seul le radical du bas les distingue.

\textbf{Erreur 3 : \zh{伦} au lieu de \zh{论}.} \zh{评论} (pínglùn, « commenter ») exige \zh{论} avec le radical de la parole \zh{讠}, car commenter, c'est parler. \zh{伦} (lún, avec \zh{亻} la personne) se trouve dans \zh{伦理} « éthique ».

\textbf{Phrase corrigée.} \zh{我的账号的密码很简单，别人很容易评论我的照片。}

\emph{Wǒ de zhànghào de mìmǎ hěn jiǎndān, biérén hěn róngyì pínglùn wǒ de zhàopiàn.}

« Le mot de passe de mon compte est très simple, les autres peuvent facilement commenter mes photos. »

Conclusion : pour les homophones, toujours vérifier le radical : \zh{贝} → argent/compte, \zh{讠} → parole, \zh{虫} → insecte/miel, \zh{宀} → caché/secret.
\end{exoresolu}

\courssub{Savoir-faire 3 : Structurer un texte avec les connecteurs}

\begin{propriete}[Tableau des connecteurs logiques utiles (niveau HSK 3-4)]
\begin{center}
\begin{tabularx}{\linewidth}{|X|X|X|}
\hline
\rowcolor{popblueL} \textbf{Relation logique} & \textbf{Structure chinoise} & \textbf{Exemple + Traduction} \\ \hline
Addition simple & \zh{也} (avant verbe) & \zh{我发布照片，也发布视频。} (Je publie des photos, je publie aussi des vidéos.) \\ \hline
Addition forte & \zh{而且} / \zh{不但……而且……} & \zh{他不但发布视频，而且点赞。} (Il non seulement publie des vidéos, mais en plus like.) \\ \hline
Cause $\rightarrow$ Conséquence & \zh{因为……所以……} & \zh{因为网络贵，所以我不常上网。} (Parce que le net est cher, je ne vais pas souvent en ligne.) \\ \hline
Concession / Opposition & \zh{虽然……但是……} / \zh{可是} & \zh{虽然我很忙，但是我回复消息。} (Bien que occupé, je réponds aux messages.) \\ \hline
Chronologie & \zh{先……然后……最后……} & \zh{先登录，然后看消息。} (D'abord se connecter, ensuite voir les messages.) \\ \hline
Condition & \zh{如果……就……} & \zh{如果密码简单，就不安全。} (Si le mot de passe est simple, ce n'est pas sûr.) \\ \hline
Avis / Conclusion & \zh{我觉得} / \zh{对我来说} / \zh{总之} & \zh{总之，保护隐私很重要。} (En somme, protéger sa vie privée est important.) \\ \hline
\end{tabularx}
\end{center}
\emph{Rappel crucial :} \zh{也}, \zh{都}, \zh{常常}, \zh{刚}, \zh{已经} se placent \textbf{avant le verbe} (ou l'auxiliaire), jamais en fin de phrase.
\end{propriete}

\begin{exoresolu}[Relier les phrases avec le bon connecteur]
Énoncé. Complète avec le connecteur qui convient (\zh{因为……所以……}, \zh{虽然……但是……}, \zh{先……然后……}, \zh{而且}) :

1. \zh{……网络很贵，……很多年轻人还是每天上网。}

2. \zh{我……完成作业，……看手机上的消息。}

3. \zh{他喜欢发布视频，……喜欢给朋友点赞。}
\tcblower
Solution détaillée.

\textbf{Phrase 1.} Le sens est une concession : « le réseau est cher » contre « les jeunes se connectent quand même ». Il faut \zh{虽然……但是……} : \zh{虽然网络很贵，但是很多年轻人还是每天上网。} \emph{Suīrán wǎngluò hěn guì, dànshì hěn duō niánqīngrén háishi měitiān shàngwǎng.} « Bien que l'Internet soit cher, beaucoup de jeunes se connectent quand même chaque jour. »

\textbf{Phrase 2.} Deux actions successives : il faut \zh{先……然后……} : \zh{我先完成作业，然后看手机上的消息。} \emph{Wǒ xiān wánchéng zuòyè, ránhòu kàn shǒujī shang de xiāoxi.} « Je fais d'abord mes devoirs, ensuite je regarde les messages sur mon téléphone. »

\textbf{Phrase 3.} Addition de deux goûts : \zh{而且} convient : \zh{他喜欢发布视频，而且喜欢给朋友点赞。} \emph{Tā xǐhuan fābù shìpín, érqiě xǐhuan gěi péngyou diǎnzàn.} « Il aime publier des vidéos, et de plus il aime liker ses amis. »

Conclusion : choisir le connecteur, c'est d'abord identifier la relation logique (concession, chronologie, addition), puis vérifier la place des mots dans la phrase.
\end{exoresolu}

\courssub{Savoir-faire 4 : Traduire du français vers le chinois (thème)}

\begin{methode}[Réussir le thème : checklist en 6 étapes]
\begin{enumerate}
\item \cle{Repérer le sujet et le verbe} de la phrase française, puis poser le squelette chinois : Sujet + (Temps/Lieu) + Verbe + Complément.
\item \cle{Placer les compléments de temps et de lieu AVANT le verbe} : « Je publie des photos le soir » → \zh{我晚上发布照片。} (pas \zh{*我发布照片晚上}).
\item \cle{Choisir le classificateur} : \zh{个} (général : \zh{一个账号}), \zh{张} (surfaces planes : \zh{一张照片}, \zh{一张桌子}), \zh{条} (longs/fins : \zh{一条消息}, \zh{一条视频}), \zh{首} (chansons/poèmes : \zh{一首歌}).
\item \cle{Gérer les mots-outils} : \zh{了} (action accomplie ponctuelle, pas habitude), \zh{的} (adjectif/pronom + nom), \zh{过} (expérience vécue).
\item \cle{Traduire les tournures idiomatiques} : « liker » → \zh{点赞} ; « s'abonner à » → \zh{关注} ; « passer du temps sur » → \zh{花时间在……上} ; « se connecter » → \zh{登录} / \zh{上线}.
\item \cle{Relire en comparant} : chaque mot français a-t-il un correspondant ? L'ordre chinois (SVO + compléments préverbaux) est-il respecté ?
\end{enumerate}
\end{methode}

\begin{exoresolu}[Thème : trois phrases sur les réseaux sociaux]
Énoncé. Traduis en chinois :

1. Ma sœur publie souvent des vidéos sur les réseaux sociaux.

2. Hier, j'ai partagé deux photos de Yaoundé avec mes amis.

3. Parce qu'il veut protéger sa vie privée, il a changé son mot de passe.
\tcblower
Solution détaillée.

\textbf{Phrase 1.} Squelette : Sujet \zh{我姐姐} + fréquence \zh{常常} (avant le verbe) + lieu \zh{在社交媒体上} (avant le verbe) + verbe \zh{发布} + complément \zh{视频}. Traduction : \zh{我姐姐常常在社交媒体上发布视频。} \emph{Wǒ jiějie chángcháng zài shèjiāo méitǐ shang fābù shìpín.} Pas de \zh{了} : c'est une habitude.

\textbf{Phrase 2.} « Hier » = complément de temps \zh{昨天}, placé après le sujet, avant le verbe. Action accomplie → \zh{了} après le verbe. « Deux photos » → classificateur \zh{张} : \zh{两张照片}. « Partager avec » → \zh{和……分享} ou \zh{跟……分享}. Traduction : \zh{昨天我和朋友们分享了两张雅温得的照片。} \emph{Zuótiān wǒ hé péngyoumen fēnxiǎng le liǎng zhāng Yǎwēndé de zhàopiàn.}

\textbf{Phrase 3.} « Parce que… » → \zh{因为……所以……} (on garde les deux en chinois). « Protéger sa vie privée » → \zh{保护隐私} (possesseur implicite) ou \zh{保护他的隐私}. « Changer » accompli → \zh{换了}. Traduction : \zh{因为他想保护隐私，所以他换了密码。} \emph{Yīnwèi tā xiǎng bǎohù yǐnsī, suǒyǐ tā huàn le mìmǎ.}

Conclusion : au thème, les trois points de contrôle sont la place des compléments (temps/lieu avant verbe), le classificateur et l'emploi (ou non) de \zh{了}.
\end{exoresolu}

\courssub{Savoir-faire 5 : Traduire du chinois vers le français (version)}

\begin{methode}[Réussir la version : démarche en 5 étapes]
\begin{enumerate}
\item \cle{Lire tout le texte} une première fois sans traduire, pour saisir le sens global (qui ? quoi ? quand ?).
\item \cle{Découper en segments} : Sujet / Compléments de temps et de lieu / Verbe (+ \zh{了}/\zh{过}/\zh{着}) / Complément d'objet.
\item \cle{Traduire segment par segment} en respectant l'ordre français : en français, les compléments de temps et de lieu vont souvent \textbf{après} le verbe, contrairement au chinois.
\item \cle{Rendre les mots-outils} : \zh{了} → passé composé / aspect accompli ; \zh{会} / \zh{要} → futur / obligation ; \zh{应该} → « devoir » (conseil moral) ; \zh{被} → voix passive ; \zh{才} → « seulement alors / ce n'est qu'ensuite que ».
\item \cle{Reformuler en français naturel} : pas de calque mot à mot ; \zh{花时间} → « passer du temps » ; \zh{上网} → « aller sur Internet / se connecter » ; \zh{比} → comparatif « plus… que ».
\end{enumerate}
\end{methode}

\begin{exoresolu}[Version : un court texte sur l'usage du téléphone]
Énoncé. Traduis en français :

\zh{现在很多喀麦隆学生每天都用社交媒体。他们早上先看手机，然后才吃饭。我觉得这样不好，因为学习比手机更重要。}
\tcblower
Solution détaillée.

\textbf{Segment 1.} \zh{现在} « maintenant / aujourd'hui » + \zh{很多喀麦隆学生} « beaucoup d'élèves camerounais » + \zh{每天都用} « utilisent chaque jour » (\zh{每天} avant le verbe, \zh{都} insiste sur la totalité) + \zh{社交媒体} « les réseaux sociaux ». → « Aujourd'hui, beaucoup d'élèves camerounais utilisent les réseaux sociaux tous les jours. »

\textbf{Segment 2.} \zh{他们早上先看手机，然后才吃饭。} Structure \zh{先……然后才……} : « d'abord… et seulement ensuite… ». \zh{才} insiste sur le retard de la deuxième action. → « Le matin, ils regardent d'abord leur téléphone, et ce n'est qu'ensuite qu'ils mangent. »

\textbf{Segment 3.} \zh{我觉得这样不好} : \zh{我觉得} « je pense que / à mon avis », \zh{这样} « ainsi, comme ça » → « Je trouve que ce n'est pas bien d'agir ainsi. » \zh{因为学习比手机更重要} : structure comparative \zh{A 比 B 更 + adjectif} « A est plus… que B » → « parce que les études sont plus importantes que le téléphone. »

\textbf{Traduction complète.} « Aujourd'hui, beaucoup d'élèves camerounais utilisent les réseaux sociaux tous les jours. Le matin, ils regardent d'abord leur téléphone, et ce n'est qu'ensuite qu'ils mangent. Je trouve que ce n'est pas bien d'agir ainsi, parce que les études sont plus importantes que le téléphone. »

Conclusion : en version, repérer les structures (\zh{先……然后才……}, \zh{比}, \zh{觉得}) avant de traduire évite les contresens.
\end{exoresolu}

\courssub{Savoir-faire 6 : Expression orale et débat — Gestion du temps d'écran}

\begin{experience}[Débat guidé : « Le téléphone à l'école : outil ou distraction ? »]
\textbf{Contexte :} Le proviseur du lycée envisage d'interdire les téléphones en classe. La vie scolaire organise un débat.
\textbf{Rôles (par groupes de 4) :}
\begin{itemize}
\item \textbf{Élève A (Pour l'interdiction) :} Utilise \zh{我觉得手机影响学习。} \zh{学生应该专心听课。} \zh{虽然手机方便，但是分心。}
\item \textbf{Élève B (Contre l'interdiction) :} Utilise \zh{手机是学习工具。} \zh{我们可以查资料、背单词。} \zh{如果管理得当，手机很有用。}
\item \textbf{Élève C (Modérateur) :} Utilise \zh{请问……} \zh{你同意吗？} \zh{总结一下：……}
\item \textbf{Élève D (Secrétaire) :} Note les arguments clés en caractères + pinyin, présente la synthèse.
\end{itemize}
\textbf{Contraintes linguistiques :} Minimum 3 connecteurs (\zh{因为……所以……}, \zh{虽然……但是……}, \zh{不但……而且……}) par intervenant. Utiliser le vocabulaire du thème (\zh{社交媒体}, \zh{分心}, \zh{查资料}, \zh{管理}).
\textbf{Durée :} 10 min préparation, 5 min débat, 3 min restitution.
\end{experience}

\courssub{Élément socioculturel : Cameroun — Chine}

\begin{savaistu}[Réseaux sociaux : WeChat / Weibo vs Facebook / WhatsApp]
\begin{itemize}
\item \textbf{Chine :} L'écosystème est dominé par \zh{微信} (Wēixìn, WeChat) : « super-application » tout-en-un (messagerie, paiement \zh{微信支付}, réseau social \zh{朋友圈}, services publics, taxi, livraison). \zh{微博} (Wēibó, Weibo) est l'équivalent de Twitter/X (microblog public). Accès à Facebook, Instagram, X, WhatsApp, YouTube bloqué par le « Grand Pare-feu » (\zh{防火长城}) sans VPN.
\item \textbf{Cameroun :} Facebook (et \zh{Facebook Lite}), WhatsApp, TikTok, Instagram dominent. L'usage mobile (data) est roi ; le paiement mobile (\zh{Mobile Money} : MTN MoMo, Orange Money) est distinct des apps sociales mais tout aussi vital.
\item \textbf{Point commun :} La place centrale du smartphone dans la vie quotidienne (paiement, info, lien social, école). Le débat sur la \zh{隐私} (vie privée) et la \zh{网络安全} (cybersécurité) est brûlant dans les deux pays.
\item \textbf{Vocabulaire bonus :} \zh{防火长城} (Grand Pare-feu), \zh{翻墙} (fānqiáng, « franchir le mur » = utiliser un VPN), \zh{实名制} (shímíngzhì, « identité réelle » obligatoire pour les comptes chinois), \zh{朋友圈} (péngyouquān, « cercle d'amis » = fil d'actualité WeChat).
\end{itemize}
\end{savaistu}

\courssub{Entraînement et évaluation}

\begin{exercice}[Entraînement écrit — Thème / Version / Production]
\textbf{Exercice 1 (Thème).} Traduis en chinois (caractères + pinyin) :
\begin{enumerate}
\item[a)] Mon mot de passe est trop simple, je dois le changer.
\item[b)] Hier soir, j'ai liké trois vidéos de mon ami.
c) Bien que les réseaux sociaux soient utiles, on ne doit pas y passer toute la journée.
\end{enumerate}

\textbf{Exercice 2 (Version).} Traduis en français :
\zh{我弟弟有一个账号，密码是“123456”。我对他说：“你的密码太简单了，别人容易进你的账号。你应该设置一个复杂的密码，保护你的隐私。”我弟弟听了我的建议，换了一个新密码。}

\textbf{Exercice 3 (Production écrite).} Rédige un message de 80–100 caractères chinois (environ 8–10 phrases) pour ton blog de classe sur le thème : « Mon utilisation des réseaux sociaux : avantages et précautions ». Utilise au moins 8 mots du lexique de la leçon et 3 connecteurs différents. Structure : Introduction → Habitudes (quoi, quand, comment) → Avantages → Risques / Précautions → Conclusion (conseil aux camarades).
\end{exercice}

\begin{corrige}[Exercice 1 — Thème]
\textbf{a)} \zh{我的密码太简单了，我得换一个。} \emph{Wǒ de mìmǎ tài jiǎndān le, wǒ děi huàn yí ge.} (Ou \zh{应该换} / \zh{必须换}).
\textbf{b)} \zh{昨天晚上我给朋友的三个视频点赞了。} \emph{Zuótiān wǎnshang wǒ gěi péngyou de sān ge shìpín diǎnzàn le.} (Structure : Sujet + Temps + \zh{给} Objet + Complément + Verbe + \zh{了}). Classificateur \zh{个} pour vidéo (ou \zh{条}).
\textbf{c)} \zh{虽然社交媒体有用，但是我们不应该整天花时间在上面。} \emph{Suīrán shèjiāo méitǐ yǒuyòng, dànshì wǒmen bù yīnggāi zhěngtiān huā shíjiān zài shangmiàn.} (\zh{整天} « toute la journée », \zh{在上面} « dessus / dedans »).
\end{corrige}

\begin{corrige}[Exercice 2 — Version]
« Mon petit frère a un compte, le mot de passe est "123456". Je lui ai dit : "Ton mot de passe est trop simple, les autres peuvent facilement entrer dans ton compte. Tu devrais configurer un mot de passe complexe, protéger ta vie privée." Mon petit frère a écouté mon conseil, il a changé pour un nouveau mot de passe. »

\textbf{Points de grammaire :} \zh{别人容易进……} (les autres entrent facilement… = risque de piratage) ; \zh{听了我的建议} (a écouté mon conseil → action accomplie) ; \zh{换了一个新密码} (en a changé un nouveau).
\end{corrige}

\begin{corrige}[Exercice 3 — Production écrite (Modèle de correction)]
\textbf{Modèle de réponse (≈ 90 caractères) :}

\zh{我有一个微信账号，名字叫“足球小将”。我每天早上和晚上都上线，常常发布训练的视频，也分享喀麦隆的新闻。很多同学关注我，给我点赞和评论，我们经常在网上讨论比赛。虽然社交媒体很有意思，可以交朋友、学中文，但是隐私很重要。我不发布家庭住址和电话号码，密码也很复杂。我觉得同学们应该管理好时间，先学习，后玩手机，保护好自己的账号安全。}

\emph{Wǒ yǒu yí ge Wēixìn zhànghào, míngzi jiào “Zúqiú Xiǎojiàng”. Wǒ měitiān zǎoshang hé wǎnshang dōu shàngxiàn, chángcháng fābù xùnliàn de shìpín, yě fēnxiǎng Kāmàilóng de xīnwén. Hěn duō tóngxué guānzhù wǒ, gěi wǒ diǎnzàn hé pínglùn, wǒmen jīngcháng zài wǎngshang tǎolùn bǐsài. Suīrán shèjiāo méitǐ hěn yǒu yìsi, kěyǐ jiāo péngyou, xué Zhōngwén, dànshì yǐnsī hěn zhòngyào. Wǒ bù fābù jiātíng zhǐzhǐ hé diànhuà hàomǎ, mìmǎ yě hěn fùzá. Wǒ juéde tóngxuemen yīnggāi guǎnlǐ hǎo shíjiān, xiān xuéxí, hòu wán shǒujī, bǎohù hǎo zìjǐ de zhànghào ānquán.}

\textbf{Analyse :} 10 phrases. Lexique : \zh{账号, 发布, 视频, 分享, 关注, 点赞, 评论, 隐私, 密码, 安全} (10/8 requis). Connecteurs : \zh{虽然……但是……}, \zh{也}, \zh{而且} (implicite dans la coordination), \zh{先……后……}. Structures : \zh{把} non requis mais \zh{管理好时间} (complément de résultat), \zh{保护好}. Respect de l'ordre des mots. Tons indiqués.
\end{corrige}

\begin{aretenir}[Grille d'auto-évaluation — Leçon 42]
\begin{center}
\begin{tabularx}{\linewidth}{|X|X|X|}
\hline
\rowcolor{popblueL} \textbf{Compétence} & \textbf{Indicateur de réussite (niveau 1ère A)} & \textbf{Auto-évaluation} \\ \hline
\textbf{Vocabulaire} & J'écris et j'utilise correctement 15+ mots du thème (caractères + pinyin + sens). & \square \square \square \\ \hline
\textbf{Caractères} & Je distingue les paires \zh{账/帐}, \zh{密/蜜}, \zh{论/伦}, \zh{赞/暂} ; j'écris l'ordre des traits de \zh{赞, 隐, 密, 评, 论}. & \square \square \square \\ \hline
\textbf{Grammaire / Connecteurs} & J'emploie correctement \zh{因为……所以……}, \zh{虽然……但是……}, \zh{先……然后……}, \zh{不但……而且……} dans des phrases complexes. & \square \square \square \\ \hline
\textbf{Thème (Fr $\rightarrow$ Zh)} & Je place le temps/lieu AVANT le verbe ; je choisis le bon classificateur (\zh{张, 条, 个}) ; j'utilise \zh{led} seulement pour l'accompli. & \square \square \square \\ \hline
\textbf{Version (Zh $\rightarrow$ Fr)} & Je repère les structures (\zh{比}, \zh{才}, \zh{被}, \zh{把}) et je rends en français naturel (pas de calque). & \square \square \square \\ \hline
\textbf{Production écrite} & Je produis un texte cohérent (80–100 car.) avec plan clair, connecteurs, lexique du thème, peu d'erreurs de caractères. & \square \square \square \\ \hline
\textbf{Expression orale} & Je participe au débat en argumentant avec le vocabulaire et les structures du thème ; prononciation tonale soignée. & \square \square \square \\ \hline
\end{tabularx}
\end{center}
\textbf{Objectif Bac :} Maîtriser l'écriture des 20 caractères « pièges » du module, réussir un thème/version de 5 phrases sans erreur de structure, produire un paragraphe d'opinion de 100 caractères en 20 min.
\end{aretenir}

\begin{attention}[Les 5 erreurs qui coûtent des points au Bac]
\begin{enumerate}
\item \cle{Tons négligés à l'écrit du pinyin} : \zh{mìmǎ} (密码, mot de passe) sans tons ou avec de mauvais tons est sanctionné ; entraîne-toi à écrire le pinyin tonalisé de chaque mot du thème : \zh{diǎnzàn}, \zh{guānzhù}, \zh{fēnxiǎng}, \zh{yǐnsī}.
\item \cle{Caractères homophones confondus} : \zh{账} (compte) / \zh{帐} (tente) ; \zh{密} (secret) / \zh{蜜} (miel) ; \zh{论} (discuter) / \zh{伦} (éthique). Réflexe : vérifier le radical (\zh{贝} = argent, \zh{讠} = parole, \zh{虫} = insecte, \zh{宀} = caché).
\item \cle{Complément de temps ou de lieu mal placé} : calque du français « Je publie des photos le soir » → faute \zh{我发布照片晚上。} Correct : \zh{我晚上发布照片。} Le temps et le lieu précèdent le verbe.
\item \cle{\zh{了} mis partout} : \zh{了} marque une action accomplie ponctuelle, jamais une habitude. Faute : \zh{我常常发布了照片。} Correct : \zh{我常常发布照片。}
\item \cle{Classificateur oublié ou faux} : on dit \zh{一张照片} (une photo, avec \zh{张}), \zh{一条消息} (un message, avec \zh{条}), \zh{一个账号} (un compte, avec \zh{个}). Écrire \zh{一照片} ou \zh{一个照片} coûte des points à chaque occurrence.
\end{enumerate}
\end{attention}

\newpage

\coursec{Exercices}

\courssub{Je vérifie mes connaissances}

\begin{exercice}[QCM : le bon caractère]
Pour chaque phrase, choisis le caractère correct parmi les trois propositions.
\begin{enumerate}
\item 我每天都要\_\_\_理我的账号。(Wǒ měitiān dōu yào \_\_\_ lǐ wǒ de zhànghào.) \\
a) 官 \quad b) 管 \quad c) 馆
\item 请不要发垃圾\_\_\_息。(Qǐng búyào fā lājī \_\_\_ xī.) \\
a) 消 \quad b) 悄 \quad c) 稍
\item 我把这条评论\_\_\_除了。(Wǒ bǎ zhè tiáo pínglùn \_\_\_ chú le.) \\
a) 删 \quad b) 册 \quad c) 栅
\item 他的\_\_\_号有五百个粉丝。(Tā de \_\_\_ hào yǒu wǔbǎi ge fěnsī.) \\
a) 张 \quad b) 账 \quad c) 帐
\end{enumerate}
\end{exercice}

\begin{exercice}[Vrai ou faux — justifie ta réponse]
Dis si chaque affirmation est vraie ou fausse, puis justifie en une phrase en français.
\begin{enumerate}
\item Le radical de \zh{管} (guǎn) est la clé du bambou \zh{竹}.
\item \zh{屏蔽} (píngbì) signifie « publier une photo ».
\item Dans la phrase \zh{我首先发布照片，然后回复评论。}, le connecteur \zh{然后} signifie « ensuite ».
\item On peut dire \zh{因为有很多垃圾信息，所以我很开心。} sans problème de sens.
\end{enumerate}
\end{exercice}

\begin{exercice}[Vocabulaire : traduis en chinois]
Donne le caractère et le pinyin des mots suivants.
\begin{enumerate}
\item un compte (sur un réseau social)
\item un commentaire
\item supprimer
\item un message indésirable (spam)
\item gérer, administrer
\end{enumerate}
\end{exercice}

\begin{exercice}[Associe le pinyin au caractère]
Relie chaque caractère à la bonne transcription pinyin.
\begin{center}
\begin{tabularx}{\linewidth}{|X|X|}\hline
\rowcolor{popblueL} Caractère & Pinyin \\ \hline
1. 账 & a) píngbì \\ \hline
2. 删 & b) guǎnlǐ \\ \hline
3. 屏蔽 & c) zhànghào \\ \hline
4. 管理 & d) shānchú \\ \hline
5. 回复 & e) huífù \\ \hline
\end{tabularx}
\end{center}
\end{exercice}

\courssub{Je m'entraîne}

\begin{exercice}[Caractères proches : complète]
Complète chaque mot avec le bon caractère : \zh{己} / \zh{已} ou \zh{未} / \zh{末}.
\begin{enumerate}
\item 自\_\_\_ (zì\_\_\_) : soi-même
\item \_\_\_经 (\_\_\_jīng) : déjà
\item 周\_\_\_ (zhōu\_\_\_) : week-end
\item \_\_\_来 (\_\_\_lái) : le futur
\end{enumerate}
\end{exercice}

\begin{exercice}[Remets les mots dans l'ordre]
Reconstitue une phrase correcte avec les mots suivants, puis traduis-la en français.
\begin{enumerate}
\item 每天 / 我 / 发布 / 照片 / 新的
\item 垃圾信息 / 因为 / 所以 / 有 / 我 / 屏蔽 / 他们
\item 网店的 / 爸爸 / 管理 / 账号 / 帮 / 我
\end{enumerate}
\end{exercice}

\begin{exercice}[Texte à trous : mots-outils et connecteurs]
Complète le texte suivant avec les mots de la liste : \zh{首先} — \zh{然后} — \zh{因为} — \zh{所以} — \zh{但是}.

\begin{textebox}[Mon compte]
Wǒ \_\_\_ fābù liǎng zhāng zhàopiàn, \_\_\_ huífù kèhù de pínglùn. \_\_\_ yǒu hěn duō lājī xìnxī, \_\_\_ wǒ měitiān dōu yào shānchú tāmen. Zhè ge gōngzuò hěn lèi, \_\_\_ wǒ hěn xǐhuan.
\end{textebox}
\end{exercice}

\begin{exercice}[Traduction rapide (version)]
Traduis les phrases suivantes en français.
\begin{enumerate}
\item \zh{我帮妈妈管理网店的账号。} (Wǒ bāng māma guǎnlǐ wǎngdiàn de zhànghào.)
\item \zh{请不要给我发垃圾信息。} (Qǐng búyào gěi wǒ fā lājī xìnxī.)
\item \zh{他删除了所有的评论。} (Tā shānchú le suǒyǒu de pínglùn.)
\end{enumerate}
\end{exercice}

\courssub{Je me prépare au Bac}

\begin{exercice}[Compréhension de l'écrit (type Bac)]
Lis le texte suivant, puis réponds aux questions en français.

\begin{textebox}[Je gère le compte de la boutique]
Wǒ jiào Lìmíng, wǒ zhù zài Yǎwēndé. Wǒ bāng bàba guǎnlǐ wǎngdiàn de zhànghào. Měitiān wǒ shǒuxiān fābù xīn de zhàopiàn, ránhòu huífù kèhù de pínglùn. Yīnwèi yǒu shíhou yǒu rén fā lājī xìnxī, suǒyǐ wǒ jiù píngbì tāmen. Zhè ge gōngzuò bù róngyì, dànshì wǒ xuéhuì le hěn duō dōngxi.\\
我帮爸爸管理网店的账号。每天我首先发布新的照片，然后回复客户的评论。因为有时候有人发垃圾信息，所以我就屏蔽他们。这个工作不容易，但是我学会了很多东西。
\end{textebox}

\begin{enumerate}
\item Où habite Lìmíng et que fait-il pour son père ? (2 pts)
\item Quelles sont les deux tâches qu'il effectue chaque jour, et dans quel ordre ? (2 pts)
\item Pourquoi bloque-t-il certaines personnes ? (2 pts)
\item Releve dans le texte deux connecteurs et donne leur traduction. (2 pts)
\item Question de langue : réécris la phrase \zh{我回复客户的评论} en la transformant avec \zh{把}. (2 pts)
\end{enumerate}
\end{exercice}

\begin{exercice}[Thème et version (type Bac)]
\textbf{A. Thème.} Traduis en chinois (caractères + pinyin) :
\begin{enumerate}
\item « D'abord, je publie des photos ; ensuite, je réponds aux commentaires des clients. » (3 pts)
\item « Parce qu'il y a des spams, je les supprime chaque jour. » (3 pts)
\end{enumerate}
\textbf{B. Version.} Traduis en français le texte suivant :
\begin{textebox}[À traduire]
Wǒ bāng bàba guǎnlǐ wǎngdiàn de zhànghào. Yǒu shíhou yǒu rén fā lājī xìnxī, wǒ jiù píngbì tāmen.\\
我帮爸爸管理网店的账号。有时候有人发垃圾信息，我就屏蔽他们。
\end{textebox}
(4 pts)
\end{exercice}

\begin{exercice}[Expression écrite (type Bac)]
\textbf{Sujet :} « Comment je gère mon compte sur les réseaux sociaux » (我怎样管理我的账号).

Rédige un texte en chinois d'\textbf{au moins 60 caractères} (environ 5 phrases). Tu dois :
\begin{enumerate}
\item utiliser obligatoirement les connecteurs \zh{首先}, \zh{然后}, \zh{但是} et la structure \zh{因为……所以……} ;
\item employer au moins quatre mots du lexique de la leçon : \zh{账号}, \zh{发布}, \zh{评论}, \zh{删除}, \zh{屏蔽}, \zh{垃圾信息}, \zh{管理} ;
\item écrire les caractères correctement, en respectant l'ordre des traits ;
\item donner le pinyin de ta première phrase seulement.
\end{enumerate}
Barème indicatif : respect du sujet et de la longueur (4 pts), exactitude des caractères (4 pts), connecteurs et structures (4 pts), richesse du vocabulaire (4 pts), pinyin correct (2 pts), présentation (2 pts).
\end{exercice}

\newpage

\coursec{Corrigés des exercices}

\begin{corrige}[Exercice 1 — QCM : le bon caractère]
\begin{enumerate}
\item \textbf{b) 管} — \zh{我每天都要管理我的账号。} (Wǒ měitiān dōu yào guǎnlǐ wǒ de zhànghào.) « Chaque jour, je dois gérer mon compte. » \zh{管} (clé du bambou \zh{竹}) signifie « gérer » ; \zh{官} = « fonctionnaire », \zh{馆} = « établissement » (饭馆, restaurant).
\item \textbf{a) 消} — \zh{请不要发垃圾消息。} (Qǐng búyào fā lājī xiāoxī.) « Merci de ne pas envoyer de messages indésirables. » \zh{消息} = « message, nouvelle » ; \zh{悄} (悄悄, silencieusement) et \zh{稍} (稍微, un peu) sont des homophones proches.
\item \textbf{a) 删} — \zh{我把这条评论删除了。} (Wǒ bǎ zhè tiáo pínglùn shānchú le.) « J'ai supprimé ce commentaire. » \zh{删} contient la clé du couteau \zh{刂} (action de couper, d'effacer) ; \zh{册} = « livre, registre », \zh{栅} = « grille ».
\item \textbf{b) 账} — \zh{他的账号有五百个粉丝。} (Tā de zhànghào yǒu wǔbǎi ge fěnsī.) « Son compte a cinq cents abonnés. » \zh{账} avec la clé \zh{贝} (coquillage = argent, biens) s'emploie pour les comptes ; \zh{张} est un classificateur (一张), \zh{帐} est une variante ancienne pour « tente / moustiquaire ».
\end{enumerate}
\textbf{Point de méthode :} face à des caractères proches, toujours identifier le \cle{radical} (clé sémantique) : il donne le domaine de sens et permet d'éliminer les mauvaises réponses.
\end{corrige}

\begin{corrige}[Exercice 2 — Vrai ou faux]
\begin{enumerate}
\item \textbf{Vrai.} \zh{管} se compose de la clé du bambou \zh{竹} (⺮) en haut et de \zh{官} (guān) en bas, qui donne le son.
\item \textbf{Faux.} \zh{屏蔽} (píngbì) signifie « bloquer, filtrer » (un utilisateur ou un contenu) ; « publier une photo » se dit \zh{发布照片} (fābù zhàopiàn).
\item \textbf{Vrai.} \zh{然后} (ránhòu) est le connecteur chronologique « ensuite, puis » ; il s'associe souvent à \zh{首先} (d'abord) : \zh{首先……，然后……。}
\item \textbf{Faux.} La phrase est grammaticalement correcte mais illogique : la structure \zh{因为……所以……} exige un rapport de cause à effet cohérent ; recevoir beaucoup de spams ne rend pas « content ». On dirait plutôt : \zh{因为有很多垃圾信息，所以我很生气。} « Parce qu'il y a beaucoup de spams, je suis en colère. »
\end{enumerate}
\end{corrige}

\begin{corrige}[Exercice 3 — Vocabulaire : traduis en chinois]
\begin{center}
\begin{tabularx}{\linewidth}{|X|X|X|X|}\hline
\rowcolor{popblueL} Français & Caractères & Pinyin & Remarque \\ \hline
un compte & 账号 & zhànghào & clé \zh{贝} (argent) \\ \hline
un commentaire & 评论 & pínglùn & clé de la parole \zh{讠} \\ \hline
supprimer & 删除 & shānchú & clé du couteau \zh{刂} \\ \hline
un message indésirable & 垃圾信息 & lājī xìnxī & litt. « information poubelle » \\ \hline
gérer, administrer & 管理 & guǎnlǐ & clé du bambou \zh{竹} \\ \hline
\end{tabularx}
\end{center}
\end{corrige}

\begin{corrige}[Exercice 4 — Associe le pinyin au caractère]
\begin{itemize}
\item 1 — c) : \zh{账} se lit zhàng, comme dans \zh{账号} (zhànghào).
\item 2 — d) : \zh{删} se lit shān, comme dans \zh{删除} (shānchú).
\item 3 — a) : \zh{屏蔽} se lit píngbì (« bloquer »).
\item 4 — b) : \zh{管理} se lit guǎnlǐ (« gérer »).
\item 5 — e) : \zh{回复} se lit huífù (« répondre »).
\end{itemize}
\textbf{Attention :} \zh{账} (zhàng, 4\ieme{} ton) ne se confond pas avec \zh{张} (zhāng, 1\ier{} ton) : même consonne, ton différent.
\end{corrige}

\begin{corrige}[Exercice 5 — Caractères proches : complète]
\begin{enumerate}
\item \zh{自己} (zìjǐ) : soi-même. \zh{己} : le trait vertical ne dépasse pas (caractère « ouvert »).
\item \zh{已经} (yǐjing) : déjà. \zh{已} : le trait dépasse à moitié.
\item \zh{周末} (zhōumò) : week-end. \zh{末} : la barre horizontale du \textbf{haut} est la plus longue (« la fin », la pointe de l'arbre).
\item \zh{未来} (wèilái) : le futur. \zh{未} : la barre horizontale du \textbf{bas} est la plus longue.
\end{enumerate}
\textbf{Moyen mnémotechnique :} pour \zh{未} / \zh{末}, « le futur est encore court en haut » (\zh{未}, petite barre en haut) ; « la fin est au sommet » (\zh{末}, grande barre en haut). Pour \zh{己} / \zh{已} / \zh{巳} : « soi-même tout ouvert, déjà à moitié ».
\end{corrige}

\begin{corrige}[Exercice 6 — Remets les mots dans l'ordre]
\begin{enumerate}
\item \zh{我每天发布新的照片。} (Wǒ měitiān fābù xīn de zhàopiàn.) « Chaque jour, je publie de nouvelles photos. » Ordre : Sujet + complément de temps + verbe + complément d'objet. Le complément de temps \zh{每天} se place avant le verbe, jamais à la fin.
\item \zh{因为有垃圾信息，所以我屏蔽他们。} (Yīnwèi yǒu lājī xìnxī, suǒyǐ wǒ píngbì tāmen.) « Parce qu'il y a des spams, je les bloque. » En chinois, on utilise souvent \zh{因为} \textbf{et} \zh{所以} ensemble, contrairement au français (« parce que » et « donc » ne se cumulent pas).
\item \zh{我帮爸爸管理网店的账号。} (Wǒ bāng bàba guǎnlǐ wǎngdiàn de zhànghào.) « J'aide papa à gérer le compte de la boutique en ligne. » Structure : Sujet + \zh{帮} + personne aidée + verbe. Le possessif \zh{的} relie \zh{网店} à \zh{账号}.
\end{enumerate}
\end{corrige}

\begin{corrige}[Exercice 7 — Texte à trous : mots-outils et connecteurs]
Texte complété :

\zh{我首先发布两张照片，然后回复客户的评论。因为有很多垃圾信息，所以我每天都要删除它们。这个工作很累，但是我很喜欢。}

(Wǒ \textbf{shǒuxiān} fābù liǎng zhāng zhàopiàn, \textbf{ránhòu} huífù kèhù de pínglùn. \textbf{Yīnwèi} yǒu hěn duō lājī xìnxī, \textbf{suǒyǐ} wǒ měitiān dōu yào shānchú tāmen. Zhè ge gōngzuò hěn lèi, \textbf{dànshì} wǒ hěn xǐhuan.)

Traduction : « D'abord, je publie deux photos, ensuite je réponds aux commentaires des clients. Parce qu'il y a beaucoup de spams, je dois les supprimer chaque jour. Ce travail est fatigant, mais je l'aime beaucoup. »

\textbf{Justification :} \zh{首先} / \zh{然后} ordonnent les actions ; \zh{因为} / \zh{所以} expriment la cause et la conséquence ; \zh{但是} introduit l'opposition finale (travail fatigant / plaisir).
\end{corrige}

\begin{corrige}[Exercice 8 — Traduction rapide (version)]
\begin{enumerate}
\item \zh{我帮妈妈管理网店的账号。} « J'aide maman à gérer le compte de la boutique en ligne. » (\zh{帮} + personne + verbe = « aider quelqu'un à faire ».)
\item \zh{请不要给我发垃圾信息。} « S'il te plaît, ne m'envoie pas de messages indésirables. » (\zh{不要} = négation de l'impératif ; \zh{给我} = « à moi ».)
\item \zh{他删除了所有的评论。} « Il a supprimé tous les commentaires. » (\zh{了} marque l'action accomplie ; \zh{所有的} = « tous les ».)
\end{enumerate}
\end{corrige}

\begin{corrige}[Exercice 9 — Compréhension de l'écrit (type Bac)]
\begin{enumerate}
\item Lìmíng habite à Yaoundé (雅温得). Il aide son père à gérer le compte de sa boutique en ligne. \emph{(2 pts : 1 pt pour le lieu, 1 pt pour l'activité.)}
\item Chaque jour, il publie d'abord de nouvelles photos (\zh{首先发布新的照片}), puis il répond aux commentaires des clients (\zh{然后回复客户的评论}). \emph{(2 pts : 1 pt par tâche, l'ordre doit être respecté.)}
\item Il bloque certaines personnes parce qu'il arrive qu'on envoie des messages indésirables (spams) sur le compte. \emph{(2 pts.)}
\item Exemples de connecteurs : \zh{首先} « d'abord », \zh{然后} « ensuite », \zh{因为……所以……} « parce que… donc… », \zh{但是} « mais ». \emph{(2 pts : 1 pt par connecteur correctement traduit.)}
\item Transformation avec \zh{把} : \zh{我把客户的评论回复了。} (Wǒ bǎ kèhù de pínglùn huífù le.) Structure : Sujet + \zh{把} + COD + verbe + \zh{了}. \emph{(2 pts : 1 pt pour la place de 把, 1 pt pour la phrase correcte.)}
\end{enumerate}
\textbf{Points de vigilance :} ne pas traduire mot à mot ; \zh{学会了很多东西} = « j'ai appris beaucoup de choses » ; dans la construction avec \zh{把}, le verbe ne peut jamais rester seul : il faut un élément après lui (ici \zh{了}).
\end{corrige}

\begin{corrige}[Exercice 10 — Thème et version (type Bac)]
\textbf{A. Thème.}
\begin{enumerate}
\item \zh{首先，我发布照片；然后，我回复客户的评论。} (Shǒuxiān, wǒ fābù zhàopiàn ; ránhòu, wǒ huífù kèhù de pínglùn.) \emph{(3 pts : connecteurs 1 pt, vocabulaire 1 pt, exactitude des caractères 1 pt.)}
\item \zh{因为有垃圾信息，所以我每天删除它们。} (Yīnwèi yǒu lājī xìnxī, suǒyǐ wǒ měitiān shānchú tāmen.) \emph{(3 pts : structure 因为…所以… 1 pt, place de 每天 1 pt, caractères 1 pt.)}
\end{enumerate}
\textbf{B. Version.} « J'aide papa à gérer le compte de la boutique en ligne. Parfois, des gens envoient des messages indésirables, alors je les bloque. » \emph{(4 pts : 2 pts par phrase, sens global et précision du lexique.)}

\textbf{Points de vigilance :} \zh{每天} se place avant le verbe ; \zh{就} dans \zh{我就屏蔽他们} marque la conséquence immédiate (« alors, du coup ») ; ne pas oublier le pronom \zh{他们} pour rendre « les ».
\end{corrige}

\begin{corrige}[Exercice 11 — Expression écrite (type Bac)]
\textbf{Proposition de rédaction modèle} (environ 90 caractères) :

\zh{我怎样管理我的账号呢？首先，我每天发布新的照片，然后回复朋友的评论。因为有时候有人发垃圾信息，所以我要删除这些信息，还要屏蔽他们。这个工作不容易，但是我很喜欢，因为我学会了很多东西。}

Pinyin de la première phrase : Wǒ zěnyàng guǎnlǐ wǒ de zhànghào ne ?

Traduction : « Comment est-ce que je gère mon compte ? D'abord, chaque jour, je publie de nouvelles photos, ensuite je réponds aux commentaires de mes amis. Parce qu'il arrive que des gens envoient des spams, je dois supprimer ces messages et aussi bloquer ces personnes. Ce travail n'est pas facile, mais je l'aime beaucoup, car j'ai appris beaucoup de choses. »

\textbf{Barème indicatif (20 pts) :}
\begin{itemize}
\item Respect du sujet et longueur (60 caractères minimum) : 4 pts ;
\item Exactitude des caractères (traits, radicaux, pas de confusion 账/张, 删/册) : 4 pts ;
\item Connecteurs et structures (\zh{首先}, \zh{然后}, \zh{但是}, \zh{因为……所以……}) : 4 pts ;
\item Richesse du vocabulaire (au moins 4 mots du lexique) : 4 pts ;
\item Pinyin correct de la première phrase (tons notés) : 2 pts ;
\item Présentation et ponctuation chinoise ，。 : 2 pts.
\end{itemize}
\textbf{Points de vigilance :} vérifier que les quatre mots du lexique sont bien présents ; ne pas écrire \zh{因为} sans \zh{所以} ; soigner les tons du pinyin (guǎnlǐ, zhànghào) ; utiliser la ponctuation pleine chasse et aucun espace entre les caractères.
\end{corrige}

\newpage

\coursec{Fiche bilan}

\begin{popfigure}
\begin{center}
\begin{tikzpicture}[mindmap, grow cyclic, every node/.style={concept}, root concept/.append style={concept color=popblue, font=\small\bfseries, text=white}, level 1/.append style={level distance=3.4cm, sibling angle=51}, level 2/.append style={level distance=2cm}]
\node[root concept, align=center]{Réseaux sociaux\\ \zh{社交媒体}}
  child[concept color=poporange]{ node[concept, align=center]{Lexique\\ \zh{词汇}} }
  child[concept color=popgreen]{ node[concept, align=center]{Écriture\\ \zh{写字}} }
  child[concept color=poppurple]{ node[concept, align=center]{Connecteurs\\ \zh{连接词}} }
  child[concept color=poppink]{ node[concept, align=center]{Thème\\ 法 $\to$ 中} }
  child[concept color=popgold]{ node[concept, align=center]{Version\\ 中 $\to$ 法} }
  child[concept color=popblue!70]{ node[concept, align=center]{Modèle de texte\\ \zh{范文}} }
  child[concept color=popgreen!70]{ node[concept, align=center]{Évaluation\\ \zh{评价}} };
\end{tikzpicture}
\end{center}
\legende{Carte mentale de la leçon 42 : expression écrite sur la gestion des réseaux sociaux.}
\end{popfigure}

\begin{bilan}
\textbf{1. Lexique clé de la gestion des réseaux sociaux}\par
\begin{center}
\begin{tabularx}{\linewidth}{|X|X|X|X|}
\hline
\rowcolor{popblueL} Caractères & Pinyin & Nature & Traduction \\
\hline
社交媒体 & shèjiāo méitǐ & n. & réseaux sociaux \\
\hline
账号 & zhànghào & n. & compte (utilisateur) \\
\hline
密码 & mìmǎ & n. & mot de passe \\
\hline
发布 & fābù & v. & publier, poster \\
\hline
评论 & pínglùn & v./n. & commenter ; commentaire \\
\hline
点赞 & diǎnzàn & v. & aimer, « liker » \\
\hline
转发 & zhuǎnfā & v. & partager, republier \\
\hline
关注 & guānzhù & v. & suivre (un compte) \\
\hline
粉丝 & fěnsī & n. & abonné, fan \\
\hline
隐私 & yǐnsī & n. & vie privée \\
\hline
网络安全 & wǎngluò ānquán & n. & sécurité en ligne \\
\hline
删除 & shānchú & v. & supprimer \\
\hline
\end{tabularx}
\end{center}

\textbf{2. Règles de grammaire à connaître par cœur}\par
\begin{center}
\begin{tabularx}{\linewidth}{|X|X|X|}
\hline
\rowcolor{popblueL} Structure & Exemple & Traduction \\
\hline
Sujet + 在 + réseau + 上 + verbe & \zh{我在微信上发布照片。} Wǒ zài Wēixìn shang fābù zhàopiàn. & Je publie des photos sur WeChat. \\
\hline
应该 / 不应该 + verbe & \zh{我们不应该泄露密码。} Wǒmen bù yīnggāi xièlòu mìmǎ. & Nous ne devrions pas révéler notre mot de passe. \\
\hline
要 + verbe + 了 (action future) & \zh{我要删除这条评论了。} Wǒ yào shānchú zhè tiáo pínglùn le. & Je vais supprimer ce commentaire. \\
\hline
先 … 然后 … & \zh{先登录，然后发布消息。} Xiān dēnglù, ránhòu fābù xiāoxi. & D'abord se connecter, ensuite publier un message. \\
\hline
越 … 越 … & \zh{粉丝越来越多。} Fěnsī yuè lái yuè duō. & Il y a de plus en plus d'abonnés. \\
\hline
\end{tabularx}
\end{center}

\textbf{3. Écriture des caractères : points de vigilance}\par
\begin{itemize}
  \item \zh{账} (compte) : clé 贝 (coquillage, argent) à gauche ; ne pas confondre avec \zh{帐} (tente).
  \item \zh{密} (secret) : clé 宀 (toit) en haut, puis 必, puis 山 en bas — écrire de haut en bas.
  \item \zh{隐} (caché) : clé 阝(阜) à gauche ; ne pas oublier le cœur 心 en bas à droite.
  \item \zh{删} (supprimer) : 册 à gauche, couteau 刂 à droite — gauche avant droite.
  \item \zh{赞} : deux 先 en haut, 贝 en bas ; compter les traits avec soin.
\end{itemize}

\textbf{4. Connecteurs utiles du texte}\par
\begin{itemize}
  \item Enchaînement : \zh{首先} shǒuxiān (d'abord), \zh{然后} ránhòu (ensuite), \zh{最后} zuìhòu (enfin).
  \item Opposition : \zh{但是} dànshì (mais), \zh{可是} kěshì (pourtant).
  \item Cause et conséquence : \zh{因为 … 所以 …} yīnwèi … suǒyǐ … (parce que … donc …).
  \item Conclusion : \zh{总之} zǒngzhī (en résumé).
\end{itemize}

\textbf{5. Méthodes express}\par
\begin{itemize}
  \item \textbf{Thème (français $\to$ chinois)} : repérer sujet + verbe + complément ; placer le complément de lieu avec 在 … 上 avant le verbe ; vérifier les tons et les classificateurs (\zh{一条评论} yì tiáo pínglùn, \zh{一个账号} yí ge zhànghào).
  \item \textbf{Version (chinois $\to$ français)} : traduire d'abord le sujet et le verbe, puis les compléments ; restituer 了 par un passé composé, 过 par une expérience (« avoir déjà »).
  \item \textbf{Production écrite} : plan en 3 parties (présentation du réseau, avantages, précautions) + connecteurs + phrase de conclusion avec \zh{总之}.
\end{itemize}
\end{bilan}

\begin{aretenir}[Checklist avant le Bac]
\begin{itemize}
  \item $\square$ Je connais par cœur le lexique clé : \zh{社交媒体、账号、密码、发布、评论、点赞、转发、关注、隐私、删除}.
  \item $\square$ Je sais écrire les caractères complexes \zh{账、密、隐、删、赞} avec le bon ordre des traits.
  \item $\square$ Je distingue les caractères proches : \zh{账/帐}, \zh{密/蜜}, \zh{己/已}.
  \item $\square$ Je place correctement le complément de lieu : \zh{在 + réseau + 上} avant le verbe.
  \item $\square$ J'utilise \zh{应该 / 不应该} pour donner un conseil sur la sécurité en ligne.
  \item $\square$ Je relie mes phrases avec \zh{首先、然后、最后、但是、因为 … 所以 …、总之}.
  \item $\square$ Je choisis le bon classificateur : \zh{条} pour un commentaire, \zh{个} pour un compte.
  \item $\square$ En version, je traduis \zh{了} par un passé composé et \zh{过} par « avoir déjà ».
  \item $\square$ Mon texte a une introduction, un développement et une conclusion avec \zh{总之}.
  \item $\square$ Je relis les tons du pinyin et la ponctuation chinoise pleine chasse （，。？！）.
  \item $\square$ Je reste simple : phrases courtes et correctes plutôt que longues et fautives.
  \item $\square$ Je gère mon temps : 5 min de plan, 20 min de rédaction, 5 min de relecture.
\end{itemize}
\end{aretenir}

\findecours

\end{document}
